\nonstopmode{}
\documentclass[a4paper]{book}
\usepackage[times,inconsolata,hyper]{Rd}
\usepackage{makeidx}
\makeatletter\@ifl@t@r\fmtversion{2018/04/01}{}{\usepackage[utf8]{inputenc}}\makeatother
% \usepackage{graphicx} % @USE GRAPHICX@
\makeindex{}
\begin{document}
\chapter*{}
\begin{center}
{\textbf{\huge Package `RomicsProcessor'}}
\par\bigskip{\large \today}
\end{center}
\ifthenelse{\boolean{Rd@use@hyper}}{\hypersetup{pdftitle = {RomicsProcessor: omics-oriented package for reproducible and sharable data processing pipelines}}}{}
\begin{description}
\raggedright{}
\item[Type]\AsIs{Package}
\item[Title]\AsIs{omics-oriented package for reproducible and sharable data
processing pipelines}
\item[Version]\AsIs{1.11.0}
\item[Author]\AsIs{@GeremyClair}
\item[Maintainer]\AsIs{The package maintainer }\email{geremy.clair@pnnl.gov}\AsIs{}
\item[Description]\AsIs{This package enables processing omics datasets using an object type called romics_object, enabling FAIR-compatible science. The romics_object contains all the details to reprocess the data identically and to create processing pipelines.}
\item[License]\AsIs{BSD 2-Clause License}
\item[Encoding]\AsIs{UTF-8}
\item[LazyData]\AsIs{true}
\item[RoxygenNote]\AsIs{7.3.3}
\item[biocViews]\AsIs{ComplexHeatmap}
\item[Imports]\AsIs{uuid, ggplot2, plotly, dplyr, tidyr, reshape2, scales,
ggrepel, ggalluvial, cowplot, ggdendro, dendextend,
RColorBrewer, viridis, parallel, doParallel, lme4, lmerTest,
emmeans, FactoMineR, missMDA, uwot, Rtsne, igraph, gplots,
ComplexHeatmap, circlize, gridExtra, grid, patchwork, FNN,
matrixStats}
\item[Depends]\AsIs{R (>= 4.0.0)}
\item[Suggests]\AsIs{pmartR, sva}
\end{description}
\Rdcontents{Contents}
\HeaderA{applyRomicsPipeline}{applyRomicsPipeline()}{applyRomicsPipeline}
%
\begin{Description}
Applies a pipeline created with the romicsCreatePipeline() function, pipelines can be edited prior to run them.
\end{Description}
%
\begin{Usage}
\begin{verbatim}
applyRomicsPipeline(romics_object, romics_pipeline)
\end{verbatim}
\end{Usage}
%
\begin{Arguments}
\begin{ldescription}
\item[\code{romics\_object}] A romics\_object created using createRomicsObject()

\item[\code{romics\_pipeline}] A pipeline created with the romicsCreatePipeline() function.
\end{ldescription}
\end{Arguments}
%
\begin{Value}
This function will return an romics\_object that has been processed through a pipeline
\end{Value}
\HeaderA{bubblePlotFeatures}{bubblePlotFeatures()}{bubblePlotFeatures}
%
\begin{Description}
will generate a bubble plot from a list of features the features will be clustered using a hclust, mean values will be represented as colors and percentage of completeness as dot size.
\end{Description}
%
\begin{Usage}
\begin{verbatim}
bubblePlotFeatures(
  romics_object,
  factor = "main",
  scale_feature = T,
  feature_list = c("a", "b"),
  gradient_colors = c("#fff1ed", "#fc896a", "#69010e")
)
\end{verbatim}
\end{Usage}
%
\begin{Arguments}
\begin{ldescription}
\item[\code{romics\_object}] A romics\_object created using romicsCreateObject()

\item[\code{factor}] has to be either 'main" or a factor of the romics\_object, the list of factor can be retrieved using the function romicsFactorNames()

\item[\code{scale\_feature}] has to be TRUE or FALSE to indicate if the calculated means have to be scaled by feature or not

\item[\code{feature\_list}] A list of features contained in the Romics\_object can be found using the function romicsFeatureNames()

\item[\code{gradient\_colors}] has to be a vector of lenght 3 containing three colors, they will be used for the low, med, and high values of the gradient respectively.
\end{ldescription}
\end{Arguments}
%
\begin{Details}
This function will generate a ggplot style figure.
\end{Details}
%
\begin{Value}
A ggplot styled figure
\end{Value}
%
\begin{Author}
Geremy Clair
\end{Author}
\HeaderA{combineRomicsObjects}{combineRomicsObjects()}{combineRomicsObjects}
%
\begin{Description}
Combine 2 or more Romics\_objects, containing the same <Samples>. It will remove the calculated stats ensure the log transformation are identical. For the combinations of multiple spatial objects use the function combineRomicsSpatialObjects()
\end{Description}
%
\begin{Usage}
\begin{verbatim}
combineRomicsObjects(...)
\end{verbatim}
\end{Usage}
%
\begin{Arguments}
\begin{ldescription}
\item[\code{...}] A succession of romics\_objects the object should have the exact same <sample> names (data column names).
\end{ldescription}
\end{Arguments}
%
\begin{Details}
This function will combine 2 or more Romics\_objects, containing the same <Samples>. It will remove the calculated stats ensure the log transformation are identical. For the combinations of multiple spatial objects use the function combineRomicsSpatialObjects()
\end{Details}
%
\begin{Value}
A combined romics\_object containing the normalized data from multiple romics\_objects
\end{Value}
%
\begin{Author}
Geremy Clair
\end{Author}
\HeaderA{createBulkMetadata}{createBulkMetadata()}{createBulkMetadata}
%
\begin{Description}
Create enhanced pseudobulk metadata with optimized processing
\end{Description}
%
\begin{Usage}
\begin{verbatim}
createBulkMetadata(
  romics_object,
  bulk_factor,
  bulk_levels,
  sample_names,
  keep_factors = NULL,
  bulk_by
)
\end{verbatim}
\end{Usage}
%
\begin{Arguments}
\begin{ldescription}
\item[\code{romics\_object}] Original romics object

\item[\code{bulk\_factor}] Factor values used for bulking (as factor or character vector)

\item[\code{bulk\_levels}] Unique levels of the bulking factor

\item[\code{sample\_names}] Names for the new pseudobulk samples

\item[\code{keep\_factors}] Character vector of factors to preserve in metadata. If NULL, preserves all factors.

\item[\code{bulk\_by}] Name of the bulking factor
\end{ldescription}
\end{Arguments}
%
\begin{Value}
Enhanced metadata data frame with optimized processing
\end{Value}
%
\begin{Author}
Geremy Clair
\end{Author}
\HeaderA{createRomicsObject}{createRomicsObject()}{createRomicsObject}
%
\begin{Description}
Create a romics\_object by combining a data and a metadata data frames
\end{Description}
%
\begin{Usage}
\begin{verbatim}
createRomicsObject(
  data,
  metadata,
  IDs,
  main_factor = "none",
  custom_colors = ROP_colors,
  omics_type = "unknown",
  omics_information = "unknown"
)
\end{verbatim}
\end{Usage}
%
\begin{Arguments}
\begin{ldescription}
\item[\code{data}] A data frame corresponding to the data.

\item[\code{metadata}] A data frame corresponding to the metadata, the columns must be the same as the data ones.

\item[\code{IDs}] A data frame corresponding to the list of alternative IDs for the romics\_object, this enables to use human readable ids for certain plots. The first column contains the same type of IDs as the imported data.

\item[\code{main\_factor}] The rowname of a metadata factor, by default the first row will be used as main factor.

\item[\code{custom\_colors}] A character vector containing the colors you want to use for your figures generated from an romics\_object.

\item[\code{omics\_type}] A character vector of length 1 indicating the type of omics data used.

\item[\code{omics\_information}] A table or vector containing omics information such as pre-processing details employed.
\end{ldescription}
\end{Arguments}
%
\begin{Value}
This function generate an romics\_object containing the following layers : data, metadata, missingdata, original\_data, original\_metadata, main\_factor, colors, steps, custom\_colors, omics\_type, omics\_information, uuid
\end{Value}
%
\begin{Author}
Geremy Clair
\end{Author}
%
\begin{Examples}
\begin{ExampleCode}
## Not run: 
ROP_romics_object <-
createRomicsObject(ROP_data, ROP_metadata, main_factor="growth_media")

## End(Not run)
\end{ExampleCode}
\end{Examples}
\HeaderA{createRomicsPipeline}{createRomicsPipeline()}{createRomicsPipeline}
%
\begin{Description}
Extracts a pipeline from the romics\_object. the pipeline can then be saved in a classic R object.
\end{Description}
%
\begin{Usage}
\begin{verbatim}
createRomicsPipeline(romics_object)
\end{verbatim}
\end{Usage}
%
\begin{Arguments}
\begin{ldescription}
\item[\code{romics\_object}] A romics\_object created using createRomicsObject()
\end{ldescription}
\end{Arguments}
%
\begin{Value}
This function will return character vector containing the stat columns previously calculated for a romics\_object, if no stats were previously calculated an error message will be displayed
\end{Value}
%
\begin{Author}
Geremy Clair
\end{Author}
\HeaderA{distribBoxplot}{distribBoxplot()}{distribBoxplot}
%
\begin{Description}
Plots the boxplot of the sample distribution using ggplot2. The colors used for the plotting will correspond to the main\_factor of the romics\_object.
\end{Description}
%
\begin{Usage}
\begin{verbatim}
distribBoxplot(romics_object, aggregate_by = "factor")
\end{verbatim}
\end{Usage}
%
\begin{Arguments}
\begin{ldescription}
\item[\code{romics\_object}] has to be a romics\_object

\item[\code{aggregate\_by}] has to be either NA (by default, in this case a Violin will be generated for each sample), or a factor from a romics\_object, the list of factor can be found using the function romicsFactorNames()
\end{ldescription}
\end{Arguments}
%
\begin{Details}
create a ggplot2 graphic output displaying the normalized or not intensites (or logged intensities) within each sample. ggplot2 methods can be used to visually adjust the plot.
\end{Details}
%
\begin{Value}
a plot generated using ggplot2 is generated with this function.
\end{Value}
%
\begin{Author}
Geremy Clair
\end{Author}
\HeaderA{distribFeature}{distribFeature()}{distribFeature}
%
\begin{Description}
Plots the distribution of a specific feature across all samples, with options to display quartiles and threshold lines.
\end{Description}
%
\begin{Usage}
\begin{verbatim}
distribFeature(
  romics_object,
  feature_name = NULL,
  bin = 0.1,
  show_quantiles = TRUE,
  threshold_value = NULL,
  threshold_color = "red",
  threshold_linetype = "dashed",
  threshold_size = 1,
  scale_x = NULL,
  scale_y = NULL,
  fill_color = NULL,
  alpha = 0.7,
  title = NULL,
  show_stats = TRUE
)
\end{verbatim}
\end{Usage}
%
\begin{Arguments}
\begin{ldescription}
\item[\code{romics\_object}] A romics\_object created using romicsCreateObject()

\item[\code{feature\_name}] Character string. Name of the feature to plot. Must match a row name in the data exactly.

\item[\code{bin}] Numeric. Width of the histogram bins. Default: 0.1

\item[\code{show\_quantiles}] Logical. Whether to display quartile lines (Q1, median, Q3) on the plot. Default: TRUE

\item[\code{threshold\_value}] Numeric. Value at which to draw a vertical threshold line. If NULL, no threshold line is drawn. Default: NULL

\item[\code{threshold\_color}] Character. Color of the threshold line. Default: "red"

\item[\code{threshold\_linetype}] Character. Line type for threshold ("solid", "dashed", "dotted"). Default: "dashed"

\item[\code{threshold\_size}] Numeric. Size/width of the threshold line. Default: 1

\item[\code{scale\_x}] Numeric vector of length 2. X-axis limits c(min, max). If NULL, uses data range. Default: NULL

\item[\code{scale\_y}] Numeric vector of length 2. Y-axis limits c(min, max). If NULL, auto-scales. Default: NULL

\item[\code{fill\_color}] Character. Fill color for histogram. If NULL, uses sample colors from romics\_object. Default: NULL

\item[\code{alpha}] Numeric. Transparency of histogram bars (0-1). Default: 0.7

\item[\code{title}] Character. Plot title. If NULL, uses feature name. Default: NULL

\item[\code{show\_stats}] Logical. Whether to display summary statistics on the plot. Default: TRUE
\end{ldescription}
\end{Arguments}
%
\begin{Details}
This function creates a histogram showing the distribution of a specific feature across all samples.
Quartile lines show Q1 (25th percentile), median (50th percentile), and Q3 (75th percentile).
The threshold line can be used to visualize cutoff values or reference points.
Automatic log-transform detection adjusts axis labels appropriately.
\end{Details}
%
\begin{Value}
A ggplot2 object
\end{Value}
%
\begin{Author}
Geremy Clair
\end{Author}
%
\begin{Examples}
\begin{ExampleCode}
## Not run: 
# Basic feature distribution
distribFeature(romics_obj, feature_name = "Gene1")

# With threshold line at value 5
distribFeature(romics_obj, feature_name = "Gene1", threshold_value = 5)

# Custom styling with no quartiles
distribFeature(romics_obj, feature_name = "Gene1",
               show_quantiles = FALSE, fill_color = "blue",
               threshold_value = 2.5, threshold_color = "green")

## End(Not run)
\end{ExampleCode}
\end{Examples}
\HeaderA{distribHistogram}{distribHistogram()}{distribHistogram}
%
\begin{Description}
Plots individual frequency plots showing the discribution of the data in each sample using grid.Extra and ggplots2. Please use the function distribHistogramGlobal() to plot the global distribution of the whole dataset.
\end{Description}
%
\begin{Usage}
\begin{verbatim}
distribHistogram(
  romics_object,
  bin = 1,
  scale_y = c(0, 100),
  scale_x = c(-10, 10),
  col = 3
)
\end{verbatim}
\end{Usage}
%
\begin{Arguments}
\begin{ldescription}
\item[\code{romics\_object}] has to be a log transformed romics\_object created using romicsCreateObject() and transformed using the function log2transform() or log10transform()

\item[\code{bin}] has to be a numerical value indicating the width of the frequency bins to use for the visualization

\item[\code{scale\_y}] has to be a numerical/double vector of length 2 indicating the y\_limits of each graph. if too low those will be automatically adjusted. by default  scale\_y=c(0,100) will be used

\item[\code{scale\_x}] has to be a numerical/double vector of length 2 indicating the x\_limits of each graph. if too low those will be automatically adjusted. by default  scale\_x=c(0,100) will be used

\item[\code{col}] has to be a double vector of lenght 1 indicating in how many columns the graphics will be displayed.
\end{ldescription}
\end{Arguments}
%
\begin{Details}
plots a complex graphic output displaying the data distribution within each sample
\end{Details}
%
\begin{Value}
a plot generated using ggplot2 is generated with this function.
\end{Value}
%
\begin{Author}
Geremy Clair
\end{Author}
\HeaderA{distribHistogramGlobal}{distribHistogramGlobal()}{distribHistogramGlobal}
%
\begin{Description}
Plots the frequency plot showing the discribution of the data in the whole data layer using ggplots2. Please use the function distribHistogram() to plot the global distribution for each sample.
\end{Description}
%
\begin{Usage}
\begin{verbatim}
distribHistogramGlobal(romics_object, bin = 1, show_quantiles = T)
\end{verbatim}
\end{Usage}
%
\begin{Arguments}
\begin{ldescription}
\item[\code{romics\_object}] has to be a log transformed romics\_object created using romicsCreateObject() and transformed using the function log2transform() or log10transform()

\item[\code{bin}] has to be a numerical value indicating the width of the frequency bins to use for the visualization
\end{ldescription}
\end{Arguments}
%
\begin{Details}
plots a complex graphic output displaying the data distribution within each sample
\end{Details}
%
\begin{Value}
a plot generated using ggplot2 is generated with this function.
\end{Value}
%
\begin{Author}
Geremy Clair
\end{Author}
\HeaderA{distribJitter}{distribJitter()}{distribJitter}
%
\begin{Description}
Plots the Jitter plot of the sample distribution using ggplot2. The colors used for the plotting will correspond to the main\_factor of the romics\_object.
\end{Description}
%
\begin{Usage}
\begin{verbatim}
distribJitter(romics_object, aggregate_by = "factor")
\end{verbatim}
\end{Usage}
%
\begin{Arguments}
\begin{ldescription}
\item[\code{romics\_object}] has to be a romics\_object

\item[\code{aggregate\_by}] has to be either NA (by default, in this case a Violin will be generated for each sample), or a factor from a romics\_object, the list of factor can be found using the function romicsFactorNames()
\end{ldescription}
\end{Arguments}
%
\begin{Details}
create a ggplot2 graphic output displaying the normalized or not intensites (or logged intensities) within each sample. ggplot2 methods can be used to visually adjust the plot.
\end{Details}
%
\begin{Value}
a plot generated using ggplot2 is generated with this function.
\end{Value}
%
\begin{Author}
Geremy Clair
\end{Author}
\HeaderA{distribRidges}{distribRidges()}{distribRidges}
%
\begin{Description}
Plots the ridges plot of the sample distribution using ggplot2. The colors used for the plotting will correspond to the main\_factor of the romics\_object.
\end{Description}
%
\begin{Usage}
\begin{verbatim}
distribRidges(romics_object)
\end{verbatim}
\end{Usage}
%
\begin{Arguments}
\begin{ldescription}
\item[\code{romics\_object}] has to be a romics\_object

\item[\code{aggregate\_by}] has to be either NA (by default, in this case a Violin will be generated for each sample), or a factor from a romics\_object, the list of factor can be found using the function romicsFactorNames()
\end{ldescription}
\end{Arguments}
%
\begin{Details}
create a ggplot2 graphic output displaying the normalized or not intensites (or logged intensities) within each sample. ggplot2 methods can be used to visually adjust the plot.
\end{Details}
%
\begin{Value}
a plot generated using ggplot2 is generated with this function.
\end{Value}
%
\begin{Author}
Geremy Clair
\end{Author}
\HeaderA{distribSina}{distribSina()}{distribSina}
%
\begin{Description}
Plots the Sina plot of the sample distribution using ggplot2. The colors used for the plotting will correspond to the main\_factor of the romics\_object.
\end{Description}
%
\begin{Usage}
\begin{verbatim}
distribSina(romics_object)
\end{verbatim}
\end{Usage}
%
\begin{Arguments}
\begin{ldescription}
\item[\code{romics\_object}] has to be a romics\_object

\item[\code{aggregate\_by}] has to be either NA (by default, in this case a Violin will be generated for each sample), or a factor from a romics\_object, the list of factor can be found using the function romicsFactorNames()
\end{ldescription}
\end{Arguments}
%
\begin{Details}
create a ggplot2 graphic output displaying the normalized or not intensites (or logged intensities) within each sample. ggplot2 methods can be used to visually adjust the plot.
\end{Details}
%
\begin{Value}
a plot generated using ggplot2 is generated with this function.
\end{Value}
%
\begin{Author}
Geremy Clair
\end{Author}
\HeaderA{distribViolin}{distribViolin()}{distribViolin}
%
\begin{Description}
Plots the violin plot of the sample distribution using ggplot2. The colors used for the plotting will correspond to the main\_factor of the romics\_object.
\end{Description}
%
\begin{Usage}
\begin{verbatim}
distribViolin(romics_object, aggregate_by = "factor")
\end{verbatim}
\end{Usage}
%
\begin{Arguments}
\begin{ldescription}
\item[\code{romics\_object}] has to be a romics\_object

\item[\code{aggregate\_by}] has to be either NA (by default, in this case a Violin will be generated for each sample), or a factor from a romics\_object, the list of factor can be found using the function romicsFactorNames()
\end{ldescription}
\end{Arguments}
%
\begin{Details}
create a ggplot2 graphic output displaying the normalized or not intensites (or logged intensities) within each sample. ggplot2 methods can be used to visually adjust the plot.
\end{Details}
%
\begin{Value}
a plot generated using ggplot2 is generated with this function.
\end{Value}
%
\begin{Author}
Geremy Clair
\end{Author}
\HeaderA{Example\_processed\_romics\_object}{Example\_unprocessed\_romics\_object}{Example.Rul.processed.Rul.romics.Rul.object}
\keyword{datasets}{Example\_processed\_romics\_object}
%
\begin{Description}
This dataset contains proteomics results of Bacillus cereus grown in 4 different media:
- Luria Bertani (lb)
- modified AOAC (maoac)
- autoclaved soil extract from INRAE avignon France (soil)
- autloclaved zucchinni puree (zucchinni)

This dataset contains proteomics results of Bacillus cereus grown in 4 different media:
- Luria Bertani (lb)
- modified AOAC (maoac)
- autoclaved soil extract from INRAE avignon France (soil)
- autloclaved zucchinni puree (zucchinni)
\end{Description}
%
\begin{Usage}
\begin{verbatim}
Example_processed_romics_object

Example_processed_romics_object
\end{verbatim}
\end{Usage}
%
\begin{Format}
A romics\_object

romics\_object
\end{Format}
%
\begin{Details}
This romics\_object was created by using the function :
romicsCreateObject(Example\_romics\_data,Example\_romics\_metadata, main\_factor = "media")

This romics\_object was created by using the function :
romicsCreateObject(Example\_romics\_data,Example\_romics\_metadata, main\_factor = "media")
and transformed using different romics functions.
to see the transformative step employed, use the following code:
romicsCreatePipeline(Example\_processed\_romics\_object)
\end{Details}
\HeaderA{Example\_romics\_data}{Example\_romics\_data}{Example.Rul.romics.Rul.data}
\keyword{datasets}{Example\_romics\_data}
%
\begin{Description}
This example is containing proteomics results of Bacillus cereus grown in 4 different media:
- Luria Bertani (lb)
- modified AOAC (maoac)
- autoclaved soil extract from INRAE avignon France (soil)
- autloclaved zucchinni puree (zucchinni)
\end{Description}
%
\begin{Usage}
\begin{verbatim}
Example_romics_data
\end{verbatim}
\end{Usage}
%
\begin{Format}
A data frame with 2459 rows and 13 columns
\end{Format}
%
\begin{Details}
All raw data are available using on MassIVE under the identifier MSV000085696 (ftp://massive.ucsd.edu/MSV000085696/)
The data was analyzed in MaxQuant and the iBAQ data was extracted using the function extractMaxQuant() on the proteinGroup.txt file
The first column contain protein identifiers while the other columns are the iBAQ intensities for the different samples
The row represent the different proteins identified.
\end{Details}
\HeaderA{Example\_romics\_metadata}{Example\_romics\_metadata}{Example.Rul.romics.Rul.metadata}
\keyword{datasets}{Example\_romics\_metadata}
%
\begin{Description}
This example is the metadata in data.frame format for the Example\_romics\_data
This dataset contains proteomics results of Bacillus cereus grown in 4 different media:
- Luria Bertani (lb)
- modified AOAC (maoac)
- autoclaved soil extract from INRAE avignon France (soil)
- autloclaved zucchinni puree (zucchinni)
\end{Description}
%
\begin{Usage}
\begin{verbatim}
Example_romics_metadata
\end{verbatim}
\end{Usage}
%
\begin{Format}
A data frame with 2 rows and 13 columns
\end{Format}
%
\begin{Details}
Metadata in romics have to be imputed in the following format:
The first column should contain the list of factors known/used for the analysis
Each column correspond to a given sample.
Note that the colnames of the metadata have to exactly match the ones of the data except for the first one.
\end{Details}
\HeaderA{extractDIANN}{extractDIANN()}{extractDIANN}
%
\begin{Description}
Extracts the quantification information from a DIANN report.pg\_matrix.tsv file
\end{Description}
%
\begin{Usage}
\begin{verbatim}
extractDIANN(file = "/filepath/report.pg_matrix.tsv")
\end{verbatim}
\end{Usage}
%
\begin{Arguments}
\begin{ldescription}
\item[\code{file}] This has to be the file path and file name of the FragPipe report.pg\_matrix.tsv file from which the information has to be extracted
\end{ldescription}
\end{Arguments}
%
\begin{Details}
This function will extracts the quantification information from one quantification time of a DIA-NN report.pg\_matrix.tsv file.
\end{Details}
%
\begin{Value}
it will return a data.frame with a first column containing the Protein IDs as first column, the other columns will be the Quantitative columns corresponding to the quantitation mode selected.
\end{Value}
%
\begin{Author}
Geremy Clair
\end{Author}
\HeaderA{extractDIANNIDs}{extractDIANNIDs()}{extractDIANNIDs}
%
\begin{Description}
Extracts the IDs information from a DIANN report.pg\_matrix.tsv file
\end{Description}
%
\begin{Usage}
\begin{verbatim}
extractDIANNIDs(file = "/filepath/report.pg_matrix.tsv")
\end{verbatim}
\end{Usage}
%
\begin{Arguments}
\begin{ldescription}
\item[\code{file}] This has to be the file path and file name of the FragPipe report.pg\_matrix.tsv file from which the information has to be extracted
\end{ldescription}
\end{Arguments}
%
\begin{Details}
This function will extracts the quantification information from one quantification time of a DIA-NN report.pg\_matrix.tsv file.
\end{Details}
%
\begin{Value}
it will return a data.frame with a first column containing the Protein IDs as first column, the other columns will be the Quantitative columns corresponding to the quantitation mode selected.
\end{Value}
%
\begin{Author}
Geremy Clair
\end{Author}
\HeaderA{extractFragPipeData}{extractFragPipeData()}{extractFragPipeData}
%
\begin{Description}
Extracts the quantification information from a FragPipe combined\_protein.csv file
\end{Description}
%
\begin{Usage}
\begin{verbatim}
extractFragPipeData(
  file = "/filepath/proteinGroups.txt",
  quantification_type = "MaxLFQ",
  remove_contaminants = TRUE
)
\end{verbatim}
\end{Usage}
%
\begin{Arguments}
\begin{ldescription}
\item[\code{file}] This has to be the file path and file name of the FragPipe combined\_protein.csv file from which the information has to be extracted

\item[\code{quantification\_type}] has to be one of the following options : 'MaxLFQ','Intensity','Total Spectral Count','Unique Spectral Count' Indicate what type of quantification needs to be extracted from the combined\_protein.csv table

\item[\code{remove\_contaminants}] has to be TRUE or FALSE, indicates if the contaminants have to be removed
\end{ldescription}
\end{Arguments}
%
\begin{Details}
This function will extracts the quantification information from one quantification time of a FragPipe combined\_protein.csv file.
\end{Details}
%
\begin{Value}
it will return a data.frame with a first column containing the Protein IDs as first column, the other columns will be the Quantitative columns corresponding to the quantitation mode selected.
\end{Value}
%
\begin{Author}
Geremy Clair
\end{Author}
\HeaderA{extractFragPipeIDs}{extractFragPipeIDs()}{extractFragPipeIDs}
%
\begin{Description}
Extracts the IDs from a FragPipe combined\_protein.csv file
\end{Description}
%
\begin{Usage}
\begin{verbatim}
extractFragPipeIDs(
  file = "/filepath/proteinGroups.txt",
  remove_contaminants = TRUE
)
\end{verbatim}
\end{Usage}
%
\begin{Arguments}
\begin{ldescription}
\item[\code{file}] This has to be the file path and file name of the FragPipe combined\_protein.csv file from which the information has to be extracted

\item[\code{remove\_contaminants}] has to be TRUE or FALSE, indicates if the contaminants have to be removed
\end{ldescription}
\end{Arguments}
%
\begin{Details}
This function will extracts the quantification information from one quantification time of a FragPipe combined\_protein.csv file.
\end{Details}
%
\begin{Value}
it will return a data.frame with a first column containing the Protein IDs as first column, the other columns will be the Quantitative columns corresponding to the quantitation mode selected.
\end{Value}
%
\begin{Author}
Geremy Clair
\end{Author}
\HeaderA{extractMaxQuant}{extractMaxQuant()}{extractMaxQuant}
%
\begin{Description}
Extracts the quantification information from a MaxQuant ProteinGroup.txt file
\end{Description}
%
\begin{Usage}
\begin{verbatim}
extractMaxQuant(
  file = "/filepath/proteinGroups.txt",
  quantification_type = "LFQ",
  remove_contaminants = TRUE,
  remove_site_only = TRUE,
  remove_reverse_matches = TRUE,
  min_peptides = 1,
  min_unique_peptides = 1,
  min_razor_peptides = 1
)
\end{verbatim}
\end{Usage}
%
\begin{Arguments}
\begin{ldescription}
\item[\code{file}] This has to be the file path and file name of the maxQuant proteinGroup.txt file from which the information has to be extracted

\item[\code{quantification\_type}] has to be one of the following options : 'LFQ','Intensity','iBAQ','MS.MS', "TMT","TMT.corrected" Indicate what type of quantification needs to be extracted from the ProteinGroup table, can be either

\item[\code{remove\_contaminants}] has to be TRUE or FALSE, indicates if the contaminant have to be removed

\item[\code{remove\_site\_only}] has to be TRUE or FALSE, indicates if proteins identified by site only have to be removed

\item[\code{remove\_reverse\_matches}] has to be TRUE or FALSE, indicates if the False Positive entries have to be removed

\item[\code{min\_peptides}] has to be numerical, indicates the minimum number of peptides for protein imporation
\end{ldescription}
\end{Arguments}
%
\begin{Details}
This function will extracts the quantification information from one quantification time of a MaxQuant ProteinGroup.txt file.
\end{Details}
%
\begin{Value}
it will return a data.frame with a first column containing the Protein IDs as first column, the other columns will be the Quantitative columns corresponding to the quantitation mode selected.
\end{Value}
%
\begin{Author}
Geremy Clair
\end{Author}
\HeaderA{extractMaxQuantIDs}{extractMaxQuantIDs()}{extractMaxQuantIDs}
%
\begin{Description}
Extracts the IDs information from a MaxQuant ProteinGroup.txt file
\end{Description}
%
\begin{Usage}
\begin{verbatim}
extractMaxQuantIDs(
  file = "/filepath/proteinGroups.txt",
  remove_contaminants = TRUE,
  remove_site_only = TRUE,
  remove_reverse_matches = TRUE,
  min_peptides = 1,
  min_unique_peptides = 1,
  min_razor_peptides = 1
)
\end{verbatim}
\end{Usage}
%
\begin{Arguments}
\begin{ldescription}
\item[\code{file}] This has to be the file path and file name of the maxQuant proteinGroup.txt file from which the information has to be extracted

\item[\code{remove\_contaminants}] has to be TRUE or FALSE, indicates if the contaminant have to be removed

\item[\code{remove\_site\_only}] has to be TRUE or FALSE, indicates if proteins identified by site only have to be removed

\item[\code{remove\_reverse\_matches}] has to be TRUE or FALSE, indicates if the False Positive entries have to be removed
\end{ldescription}
\end{Arguments}
%
\begin{Details}
This function will extracts the IDs information from one quantification time of a MaxQuant ProteinGroup.txt file.
\end{Details}
%
\begin{Value}
it will return a data.frame with a first column containing the Protein IDs the following columns will be the following : 'majority.protein.ids', 'fasta.headers', 'peptide.counts.all','peptide.counts.razor.unique','peptide.counts..unique.','fasta.headers','number.of.proteins','peptides','razor...unique.peptides', 'unique.peptides'
\end{Value}
%
\begin{Author}
Geremy Clair
\end{Author}
\HeaderA{extractScils}{extractScils()}{extractScils}
%
\begin{Description}
Extracts intensity data, metadata, and feature IDs from a SCiLS '.slx' file.
The function efficiently processes large datasets common in mass spectrometry imaging.

Extracts intensity data, metadata, and feature IDs from a SCiLS '.slx' file.
The function efficiently processes large datasets common in mass spectrometry imaging.
\end{Description}
%
\begin{Usage}
\begin{verbatim}
extractScils(
  file,
  scils_executable,
  port,
  feature_list,
  normId,
  metadata_fields = NULL
)

extractScils(
  file,
  scils_executable,
  port,
  feature_list,
  normId,
  metadata_fields = NULL
)
\end{verbatim}
\end{Usage}
%
\begin{Arguments}
\begin{ldescription}
\item[\code{file}] Character string. Path to the SCiLS '.slx' file.

\item[\code{scils\_executable}] Character string. Path to the SCiLS executable.
Default: 'C:/Program Files/SCiLS/SCiLS Lab/scilsMsServer.exe'

\item[\code{port}] Numeric. Port to use for the connection with the SCiLS environment.
Default: 8082

\item[\code{feature\_list}] Character string. Name of the feature list to import.
Default: "All Features"

\item[\code{normId}] Character string. Normalization ID to apply during import.
Default: "" (no normalization)

\item[\code{metadata\_fields}] Character vector. Names of metadata fields to extract.
Default: NULL (extracts all available metadata fields)
\end{ldescription}
\end{Arguments}
%
\begin{Details}
This function extracts data from a SCiLS '.slx' file through the SCiLS Lab API.
It requires both the '.slx' file and corresponding '.sbd' file to be present in the same directory.
The necessary companion '.sbd' file must be in the same directory as the '.slx' file.
The function was optimized for handling large datasets with potentially hundreds of thousands of
data points. It might still be slow to operate due to some limitation of the SCIls API.

It imports: the intensity matrix for the selected features, the Spot IDs and associated metadata, the Detailed information about each feature

This function extracts data from a SCiLS '.slx' file through the SCiLS Lab API.
It requires both the '.slx' file and corresponding '.sbd' file to be present in the same directory.

The function is optimized for handling large datasets with potentially hundreds of thousands of
data points. It imports:
\begin{itemize}

\item{} Intensity matrix for the selected features
\item{} Spot IDs and associated metadata
\item{} Detailed information about each feature

\end{itemize}


The necessary companion '.sbd' file must be in the same directory as the '.slx' file.
\end{Details}
%
\begin{Value}
A list with three components:
\begin{ldescription}
\item[\code{data}] Data frame containing the intensity matrix with features as rows and samples as columns
\item[\code{metadata}] Data frame containing spot IDs and other spot attributes
\item[\code{IDs}] Data frame containing detailed information about each feature

\end{ldescription}
A list with three components:
\begin{description}

\item[data] Data frame containing the intensity matrix with features as rows and samples as columns
\item[metadata] Data frame containing spot IDs and other spot attributes
\item[IDs] Data frame containing detailed information about each feature

\end{description}

\end{Value}
%
\begin{Author}
Geremy Clair, Brittney Gorman
\end{Author}
%
\begin{Examples}
\begin{ExampleCode}
## Not run: 
# Basic usage with default parameters
scils_data <- extractScils("path/to/your/file.slx")

# Custom parameters
scils_data <- extractScils(
  file = "path/to/your/file.slx",
  scils_executable = "C:/Custom/Path/scilsMsServer.exe",
  port = 8083,
  feature_list = "My Custom Feature List",
  normId = "TIC"
)

# Access the components
data <- scils_data$data
metadata <- scils_data$metadata
feature_ids <- scils_data$IDs

## End(Not run)

\end{ExampleCode}
\end{Examples}
\HeaderA{extractScilsOLD}{extractScilsOLD()}{extractScilsOLD}
%
\begin{Description}
Extracts intensity data, metadata, and feature IDs from a SCiLS '.slx' file.
The function efficiently processes large datasets common in mass spectrometry imaging.
\end{Description}
%
\begin{Usage}
\begin{verbatim}
extractScilsOLD(
  file,
  scils_executable,
  port,
  feature_list,
  normId,
  metadata_fields = NULL
)
\end{verbatim}
\end{Usage}
%
\begin{Arguments}
\begin{ldescription}
\item[\code{file}] Character string. Path to the SCiLS '.slx' file.

\item[\code{scils\_executable}] Character string. Path to the SCiLS executable.
Default: 'C:/Program Files/SCiLS/SCiLS Lab/scilsMsServer.exe'

\item[\code{port}] Numeric. Port to use for the connection with the SCiLS environment.
Default: 8082

\item[\code{feature\_list}] Character string. Name of the feature list to import.
Default: "All Features"

\item[\code{normId}] Character string. Normalization ID to apply during import.
Default: "" (no normalization)

\item[\code{metadata\_fields}] Character vector. Names of metadata fields to extract.
Default: NULL (extracts all available metadata fields)
\end{ldescription}
\end{Arguments}
%
\begin{Details}
This function extracts data from a SCiLS '.slx' file through the SCiLS Lab API.
It requires both the '.slx' file and corresponding '.sbd' file to be present in the same directory.

The function is optimized for handling large datasets with potentially hundreds of thousands of
data points. It imports:
\begin{itemize}

\item{} Intensity matrix for the selected features
\item{} Spot IDs and associated metadata
\item{} Detailed information about each feature

\end{itemize}


The necessary companion '.sbd' file must be in the same directory as the '.slx' file.
\end{Details}
%
\begin{Value}
A list with three components:
\begin{description}

\item[data] Data frame containing the intensity matrix with features as rows and samples as columns
\item[metadata] Data frame containing spot IDs and other spot attributes
\item[IDs] Data frame containing detailed information about each feature

\end{description}

\end{Value}
%
\begin{Author}
Geremy Clair, Brittney Gorman
\end{Author}
%
\begin{Examples}
\begin{ExampleCode}
## Not run: 
# Basic usage with default parameters
scils_data <- extractScils("path/to/your/file.slx")

# Custom parameters
scils_data <- extractScils(
  file = "path/to/your/file.slx",
  scils_executable = "C:/Custom/Path/scilsMsServer.exe",
  port = 8083,
  feature_list = "My Custom Feature List",
  normId = "TIC"
)

# Access the components
data <- scils_data$data
metadata <- scils_data$metadata
feature_ids <- scils_data$IDs

## End(Not run)

\end{ExampleCode}
\end{Examples}
\HeaderA{factorCountLevels}{factorCountLevels()}{factorCountLevels}
%
\begin{Description}
Indicates in a vector how many times each level of a factor is detected, with optional visualization
\end{Description}
%
\begin{Usage}
\begin{verbatim}
factorCountLevels(
  romics_object = romics_object,
  factor = "factor",
  plot = FALSE,
  textsize = 3,
  fill_color = "gray70",
  order = FALSE
)
\end{verbatim}
\end{Usage}
%
\begin{Arguments}
\begin{ldescription}
\item[\code{romics\_object}] A romics\_object created using romicsCreateObject().

\item[\code{factor}] A character vector (length=1) containing a factor of the romics\_object. The list of factors can be identified by using the function romicsFactorNames().

\item[\code{plot}] Logical indicating whether to return a ggplot bar chart. Default is FALSE (returns numeric vector).

\item[\code{textsize}] Numeric value for text size in the plot. Default is 3. Only used when plot=TRUE.

\item[\code{fill\_color}] Color for the bars in the plot. Default is "gray70". Only used when plot=TRUE.

\item[\code{order}] Logical indicating whether to order bars by count (descending). Default is FALSE. Only used when plot=TRUE.
\end{ldescription}
\end{Arguments}
%
\begin{Details}
Generates a count of how many times each level is seen in a factor. Can return either a named numeric vector or a ggplot bar chart.
\end{Details}
%
\begin{Value}
If plot=FALSE: A named numeric vector with counts for each factor level. If plot=TRUE: A ggplot object showing the counts as a bar chart.
\end{Value}
%
\begin{Author}
Geremy Clair
\end{Author}
%
\begin{Examples}
\begin{ExampleCode}
## Not run: 
# Get counts as a named vector
counts <- factorCountLevels(romics_object = my_romics, factor = "Treatment")

# Generate a plot
factorCountLevels(romics_object = my_romics, factor = "Treatment", plot = TRUE)

# Customize the plot
factorCountLevels(romics_object = my_romics, factor = "Sex",
                  plot = TRUE, textsize = 4, fill_color = "gray50", order = TRUE)

# Access specific level count
counts <- factorCountLevels(romics_object = my_romics, factor = "Sex")
counts["Male"]

## End(Not run)
\end{ExampleCode}
\end{Examples}
\HeaderA{FeatureCorrelation}{FeatureCorrelation()}{FeatureCorrelation}
%
\begin{Description}
calculate the correlation matrix for the features of a romics\_object
\end{Description}
%
\begin{Usage}
\begin{verbatim}
FeatureCorrelation(romics_object, corr_type = "pearson", use = "everything")
\end{verbatim}
\end{Usage}
%
\begin{Arguments}
\begin{ldescription}
\item[\code{romics\_object}] an object of type romics\_object

\item[\code{corr\_type}] has to be either 'pearson', 'kendall',or 'spearman', indicates the type of correlation calculated.

\item[\code{use}] has to be either 'everything', 'all.obs', 'complete.obs', 'na.or.complete', or 'pairwise.complete.obs'
\end{ldescription}
\end{Arguments}
%
\begin{Details}
calculate a correlation matrix for the features of an Romics\_object
\end{Details}
%
\begin{Value}
This function will return a correlation matrix for the features of an romics\_object
\end{Value}
%
\begin{Author}
Geremy Clair
\end{Author}
\HeaderA{FeatureCorrelationHclust}{FeatureCorrelationHclust()}{FeatureCorrelationHclust}
%
\begin{Description}
calculate the feature correlation matrix for an romics\_object and generate a hclust object.
\end{Description}
%
\begin{Usage}
\begin{verbatim}
FeatureCorrelationHclust(
  romics_object,
  corr_type = "pearson",
  use = "everything",
  hclust_method = "ward.D",
  dist_method = "euclidean"
)
\end{verbatim}
\end{Usage}
%
\begin{Arguments}
\begin{ldescription}
\item[\code{romics\_object}] an object of type romics\_object

\item[\code{corr\_type}] has to be either 'pearson', 'kendall',or 'spearman', indicates the type of correlation calculated.

\item[\code{use}] has to be either 'everything', 'all.obs', 'complete.obs', 'na.or.complete', or 'pairwise.complete.obs'

\item[\code{hclust\_method}] has to be in 'ward.D', 'ward.D2', 'single', 'complete', 'average', 'mcquitty', 'median', or 'centroid'

\item[\code{dist\_method}] has to be in 'correlation.dissimilarity','euclidean', 'maximum', 'manhattan', 'canberra', 'binary' or 'minkowski'.
\end{ldescription}
\end{Arguments}
%
\begin{Details}
calculate the feature correlation matrix for the samples of an romics object and generate a hclust object.
\end{Details}
%
\begin{Value}
This function will return a hclust result that can be used to generate hclust plots.
\end{Value}
%
\begin{Author}
Geremy Clair
\end{Author}
\HeaderA{FeatureCorrelationHclustPlot}{FeatureCorrelationHclustPlot()}{FeatureCorrelationHclustPlot}
%
\begin{Description}
calculate the Feature correlation matrix for the samples of an romics\_object and generate a hclust plot.
\end{Description}
%
\begin{Usage}
\begin{verbatim}
FeatureCorrelationHclustPlot(
  romics_object,
  corr_type = "pearson",
  use = "everything",
  hclust_method = "ward.D",
  dist_method = "euclidean",
  plot_type = "dendrogram",
  number_of_clusters = 0,
  cex = 1,
  cluster_colors = c("c1", "c2", "c3", "c4")
)
\end{verbatim}
\end{Usage}
%
\begin{Arguments}
\begin{ldescription}
\item[\code{romics\_object}] an object of type romics\_object

\item[\code{corr\_type}] has to be either 'pearson', 'kendall',or 'spearman', indicates the type of correlation calculated.

\item[\code{use}] has to be either 'everything', 'all.obs', 'complete.obs', 'na.or.complete', or 'pairwise.complete.obs'

\item[\code{hclust\_method}] has to be in 'ward.D', 'ward.D2', 'single', 'complete', 'average', 'mcquitty', 'median', or 'centroid'

\item[\code{dist\_method}] has to be in 'correlation.dissimilarity','euclidean', 'maximum', 'manhattan', 'canberra', 'binary' or 'minkowski'.

\item[\code{plot\_type}] has to be either 'dendrogram', 'unrooted', or 'fan'

\item[\code{number\_of\_clusters}] has to be a numeric value to indicate the number of clusters

\item[\code{cex}] has to be a numeric value to indicate the label font size

\item[\code{cluster\_colors}] has to be a character vector containing the list of colors used for the clusters.
\end{ldescription}
\end{Arguments}
%
\begin{Details}
calculate the Feature correlation matrix for the samples of an romics object and generate a hclust plot
\end{Details}
%
\begin{Value}
This function will return a hclust plot generated with the base R plot() function.
\end{Value}
%
\begin{Author}
Geremy Clair
\end{Author}
\HeaderA{FeatureCorrelationHeatmap}{FeatureCorrelationHeatmap()}{FeatureCorrelationHeatmap}
%
\begin{Description}
calculate the feature correlation matrix of an romics\_object and hclust the results are displayed as an heatmap.
\end{Description}
%
\begin{Usage}
\begin{verbatim}
FeatureCorrelationHeatmap(
  romics_object,
  corr_type = "pearson",
  use = "everything",
  number_of_clusters = 0,
  cex = 1,
  cluster_colors = c("c1", "c2", "c3", "c4"),
  heatmap_col = viridis(20),
  dist_method = "euclidean",
  hclust_method = "ward.D",
  cellnote = F
)
\end{verbatim}
\end{Usage}
%
\begin{Arguments}
\begin{ldescription}
\item[\code{romics\_object}] an object of type romics\_object

\item[\code{corr\_type}] has to be either 'pearson', 'kendall',or 'spearman', indicates the type of correlation calculated.

\item[\code{use}] has to be either 'everything', 'all.obs', 'complete.obs', 'na.or.complete', or 'pairwise.complete.obs'

\item[\code{number\_of\_clusters}] has to be a numeric value to indicate the number of clusters

\item[\code{cex}] has to be a numeric value to indicate the label font size

\item[\code{cluster\_colors}] has to be a character vector containing the list of colors used for the clusters.

\item[\code{heatmap\_col}] has to be a color gradient (e.g.viridis(20), greenred(50))

\item[\code{dist\_method}] has to be in 'correlation.dissimilarity','euclidean', 'maximum', 'manhattan', 'canberra', 'binary' or 'minkowski'.

\item[\code{hclust\_method}] has to be in 'ward.D', 'ward.D2', 'single', 'complete', 'average', 'mcquitty', 'median', or 'centroid'

\item[\code{cellnote}] has to be TRUE or FALSE to indicate if the correlation numbers should be displayed onto the heatmap.
\end{ldescription}
\end{Arguments}
%
\begin{Value}
This function will return a an heatmap generated with heatmap.2().
\end{Value}
%
\begin{Author}
Geremy Clair
\end{Author}
\HeaderA{findStatisticalColumns}{Helper function to find matching statistical columns}{findStatisticalColumns}
%
\begin{Description}
Helper function to find matching statistical columns
\end{Description}
%
\begin{Usage}
\begin{verbatim}
findStatisticalColumns(romics_object, factor_levels, pval_type)
\end{verbatim}
\end{Usage}
%
\begin{Arguments}
\begin{ldescription}
\item[\code{romics\_object}] A romics\_object with statistics

\item[\code{factor\_levels}] Vector of factor levels to compare

\item[\code{pval\_type}] Either "p" or "padj"
\end{ldescription}
\end{Arguments}
%
\begin{Value}
Data frame with combinations and their corresponding column names
\end{Value}
\HeaderA{getTrend}{getTrend()}{getTrend}
%
\begin{Description}
This function extract the linear and quadratic trend pvalues and direction based on the pvalue it attributes the best fitted trend
\end{Description}
%
\begin{Usage}
\begin{verbatim}
getTrend(sData, m = NULL, p = 0.05, type = "both")
\end{verbatim}
\end{Usage}
%
\begin{Arguments}
\begin{ldescription}
\item[\code{sData}] is a serie of data. It should be a vector of numerical values corresponding to the different points of the serie (e.g. time serie)

\item[\code{m}] is the metadata associated to the serie (e.g. for a time-serie, it would be the time) if m is null a numeric sequence from 1:n will be used.

\item[\code{p}] is the minimum fitting pvalue (by default p = 0.05)

\item[\code{type}] indicates the type of functions to be fitted for the data serie. Has to be either "quadratic" (quadratic only), "linear" (linear only), or "both" (linear+quardatic)
\end{ldescription}
\end{Arguments}
%
\begin{Details}
this function enables to calculate and extract the linear or quardratic trend p values and the best fitted trend (e.g. linear\_increasing, linear\_decreasing, quadratic\_concave,quadratic\_convex)
\end{Details}
%
\begin{Value}
This function returns a list
\end{Value}
%
\begin{Author}
Geremy Clair and Feng Song
\end{Author}
\HeaderA{heatmapFeatures}{heatmapFeatures()}{heatmapFeatures}
%
\begin{Usage}
\begin{verbatim}
heatmapFeatures(
  romics_object,
  factor = "main",
  scale_feature = TRUE,
  feature_list = NULL,
  filter_by_stat_column = NULL,
  stat_column_filter = NULL,
  gradient_colors = NULL,
  viridis_option = "viridis",
  show_completeness = TRUE,
  show_dendrogram = TRUE,
  show_feature_names = TRUE,
  show_clusters = FALSE,
  show_cluster_legend = FALSE,
  n_clusters = NULL,
  clustering_method = "ward.D"
)
\end{verbatim}
\end{Usage}
%
\begin{Arguments}
\begin{ldescription}
\item[\code{romics\_object}] A romics object containing \$data, \$statistics, and \$metadata layers

\item[\code{factor}] Character string specifying the factor to use for grouping samples. Default: "main" uses the main factor. Use romicsFactorNames() to see available factors

\item[\code{scale\_feature}] TRUE or FALSE, if TRUE scales features by rows (z-score normalization). Default: TRUE

\item[\code{feature\_list}] Character vector of features to include in heatmap. If NULL, uses all features. Default: NULL

\item[\code{filter\_by\_stat\_column}] Character string specifying the statistics column to filter features. If NULL, no filtering applied. Example: "Healthy\_vs\_Disease\_Ttest\_padj". Default: NULL

\item[\code{stat\_column\_filter}] Character string specifying the filter expression using operators <, >, <=, >=, ==, !=. Example: "<0.05" or ">1.5". Only used if filter\_by\_stat\_column is specified. Default: NULL

\item[\code{gradient\_colors}] Character vector of colors for the heatmap gradient. If NULL, uses viridis color scale. For 3 colors: c(low, mid, high). For 2 colors: c(low, high). Default: NULL

\item[\code{viridis\_option}] Character string specifying the viridis color scale option. Only used if gradient\_colors is NULL. Accepted values: "viridis", "plasma", "inferno", "magma", "cividis", "rocket", "mako", "turbo". Default: "viridis"

\item[\code{show\_completeness}] TRUE or FALSE, if TRUE shows data completeness as tile transparency (alpha). If FALSE, completely missing data (0

\bsl{}itemshow\_dendrogramTRUE or FALSE, if TRUE displays hierarchical clustering dendrogram on the left side of the heatmap. Requires ggdendro package. Default: TRUE

\bsl{}itemshow\_feature\_namesTRUE or FALSE, if TRUE displays feature names on the right side of the heatmap. Default: TRUE

\bsl{}itemshow\_clustersTRUE or FALSE, if TRUE displays cluster assignments as a sidebar. Default: FALSE

\bsl{}itemshow\_cluster\_legendTRUE or FALSE, if TRUE displays the cluster legend. Only used if show\_clusters is TRUE. Default: FALSE

\bsl{}itemn\_clustersInteger specifying the number of clusters to cut the dendrogram into. Only used if show\_clusters is TRUE. Default: NULL (no clustering)

\bsl{}itemclustering\_methodCharacter string specifying hierarchical clustering method. Accepted values: "complete", "average", "single", "ward.D", "ward.D2". Default: "complete"

\end{ldescription}

A ggplot/cowplot object containing the heatmap visualization (with optional dendrogram and cluster sidebar).
Clustering results are attached as attributes (accessible via attr(plot, "clustering\_results")).


Creates an interactive heatmap visualization from a romics object with hierarchical clustering and optional statistical filtering


The function creates an interactive heatmap with the following features:
- Features are automatically ordered by hierarchical clustering on an imputed dataset
- Imputation uses values 2 standard deviations below the minimum observed value, ensuring robust clustering
- The displayed heatmap shows only real data (no imputation artifacts)
- Tiles can be colored by mean expression with optional row/column scaling
- Data completeness can be visualized via transparency or as black tiles for missing data
- Optional statistical filtering reduces features to significant results
- Interactive tooltips show exact values on hover
- Dendrogram provides visual representation of feature relationships
- Viridis color scale can be customized with different options (viridis, plasma, inferno, magma, cividis, rocket, mako, turbo)
- Optionally cut dendrogram into clusters and display cluster assignments

If feature\_list is provided, the heatmap will only display those features.
If filter\_by\_stat\_column and stat\_column\_filter are both specified, features are filtered before visualization.
If gradient\_colors is provided, viridis\_option is ignored.
If show\_clusters is TRUE and n\_clusters is specified, the dendrogram is cut into clusters and a sidebar shows cluster assignments.


Geremy Clair

\end{Arguments}
\HeaderA{imputeMinDiv2}{imputeMinDiv2()}{imputeMinDiv2}
%
\begin{Description}
Imputes the data layer of the romics\_object using the minimum value of the table divided by 2. this function will work if the data only contains positive values.
\end{Description}
%
\begin{Usage}
\begin{verbatim}
imputeMinDiv2(romics_object)
\end{verbatim}
\end{Usage}
%
\begin{Arguments}
\begin{ldescription}
\item[\code{romics\_object}] has to be an romics\_object created using romicsCreateObject()
\end{ldescription}
\end{Arguments}
%
\begin{Details}
This function will use the minimum value of the data table divided by 2 to impute the missing values of the data layer
\end{Details}
%
\begin{Value}
The function will return a modified romics\_object that will have imputed data, however the missingdata layer will conserve the location of the missingness, the missingness can subsequently be restored using the function romicsRestoreMissing().
\end{Value}
%
\begin{Author}
Geremy Clair
\end{Author}
\HeaderA{imputeMissing}{imputeMissing()}{imputeMissing}
%
\begin{Description}
Imputes the data using a normal distribution down-shifted from the median by a user defined number of standard deviations and a user defined width. the distribution of the imputed data can be evaluated using the function imputeMissingEval().
\end{Description}
%
\begin{Usage}
\begin{verbatim}
imputeMissing(romics_object, nb_stdev = 1.8, width_stdev = 0.5, seed = 42)
\end{verbatim}
\end{Usage}
%
\begin{Arguments}
\begin{ldescription}
\item[\code{romics\_object}] has to be an romics\_object created using romicsCreateObject()

\item[\code{nb\_stdev}] Numeric. Number of standard deviations to shift the imputation distribution from the median. Default: 1.8

\item[\code{width\_stdev}] Numeric. Width of the imputation distribution in standard deviations. Default: 0.5

\item[\code{seed}] Integer. Random seed for reproducible results. Default: 42
\end{ldescription}
\end{Arguments}
%
\begin{Details}
This function will impute the data using the method described in the Perseus paper by Tyranova et al. 2016 (Nature Method). By default the data is imputed with values in a normal distribution 1.8 standard deviation away from the median and a width of distribution of 0.5.
\end{Details}
%
\begin{Value}
The function will return a modified romics\_object that will have imputed data, however the missingdata layer will conserve the location of the missingness, the missingness can subsequently be restored using the function romicsRestoreMissing().
\end{Value}
%
\begin{Author}
Geremy Clair
\end{Author}
\HeaderA{imputeMissingEval}{imputeMissingEval()}{imputeMissingEval}
%
\begin{Description}
Plots the distribution of the data to be imputed using the imputeMissing() function. Enables  to optimize the parameters prior to apply the imputeMissing() function.
\end{Description}
%
\begin{Usage}
\begin{verbatim}
imputeMissingEval(
  romics_object,
  nb_stdev = 1.8,
  width_stdev = 0.5,
  bin = 1,
  scale_x = "auto",
  sample_size = 1e+05,
  set.seed = TRUE,
  seed = 42
)
\end{verbatim}
\end{Usage}
%
\begin{Arguments}
\begin{ldescription}
\item[\code{romics\_object}] has to be an romics\_object created using romicsCreateObject()

\item[\code{nb\_stdev}] number of standard deviations away from the median to place the missingness. Default: 1.8

\item[\code{width\_stdev}] width of the random values used for imputation expressed in standard deviation from the data. Default: 0.5

\item[\code{bin}] binwidth for the histogram. Default: 1

\item[\code{scale\_x}] numeric vector of length 2 for x-axis limits. Default: c(-10,10)

\item[\code{sample\_size}] maximum number of data points to use for visualization (for performance with large datasets). Default: 100000

\item[\code{set.seed}] logical indicating whether to set a seed for reproducible results. Default: TRUE

\item[\code{seed}] numeric value for the random seed when set.seed=TRUE. Default: 42
\end{ldescription}
\end{Arguments}
%
\begin{Details}
This function does not alter the romics\_object, it plots the distribution of the whole data and of the imputed values using the method described in the Perseus paper by Tyranova et al. 2016 (Nature Method). By default the data is imputed with values in a normal distribution 1.8 standard deviation away from the median  and a width of distribution of 0.5.
\end{Details}
%
\begin{Value}
This function will return a ggplot2 geom\_histogram plot. it can then be further visually adjusted using the ggplot2 commands
\end{Value}
%
\begin{Author}
Geremy Clair
\end{Author}
\HeaderA{is.romicsObject}{is.romicsObject()}{is.romicsObject}
%
\begin{Description}
Enables to check if the romics\_object is in the appropriate format. Stops execution if UUID is missing.
\end{Description}
%
\begin{Usage}
\begin{verbatim}
is.romicsObject(romics_object)
\end{verbatim}
\end{Usage}
%
\begin{Arguments}
\begin{ldescription}
\item[\code{romics\_object}] A romics\_object created using createRomicsObject()
\end{ldescription}
\end{Arguments}
%
\begin{Value}
This function will return TRUE if the object is a valid romics\_object, or stop execution with an error message
\end{Value}
%
\begin{Author}
Geremy Clair
\end{Author}
\HeaderA{lmp}{lmp()}{lmp}
%
\begin{Description}
This function extract the overall p-value of a linear model
\end{Description}
%
\begin{Usage}
\begin{verbatim}
lmp(modelobject)
\end{verbatim}
\end{Usage}
%
\begin{Arguments}
\begin{ldescription}
\item[\code{modelobject}] as to be an object of the class 'lm'
\end{ldescription}
\end{Arguments}
%
\begin{Details}
This function extract the p-value of an object of class 'lm'
\end{Details}
%
\begin{Value}
This function returns the p-value of an object of class 'lm'
\end{Value}
%
\begin{Author}
Geremy Clair and Feng Song
\end{Author}
\HeaderA{log10transform}{log10transform()}{log10transform}
%
\begin{Description}
log10-tranforms the romics\_object data layer.
\end{Description}
%
\begin{Usage}
\begin{verbatim}
log10transform(romics_object)
\end{verbatim}
\end{Usage}
%
\begin{Arguments}
\begin{ldescription}
\item[\code{romics\_object}] has to be an romics\_object created using romicsCreateObject()
\end{ldescription}
\end{Arguments}
%
\begin{Details}
will log10 transform the romics object
\end{Details}
%
\begin{Value}
This function returns the transformed romics\_object with updated data layer
\end{Value}
%
\begin{Author}
Geremy Clair
\end{Author}
\HeaderA{log2transform}{log2transform()}{log2transform}
%
\begin{Description}
log2-tranforms the romics\_object data layer.
\end{Description}
%
\begin{Usage}
\begin{verbatim}
log2transform(romics_object)
\end{verbatim}
\end{Usage}
%
\begin{Arguments}
\begin{ldescription}
\item[\code{romics\_object}] has to be an romics\_object created using romicsCreateObject()
\end{ldescription}
\end{Arguments}
%
\begin{Details}
will log2 transform the romics object
\end{Details}
%
\begin{Value}
This function returns the transformed romics\_object with updated data layer
\end{Value}
%
\begin{Author}
Geremy Clair
\end{Author}
\HeaderA{medianCenterFeature}{medianCenterFeature()}{medianCenterFeature}
%
\begin{Description}
Normalizes the features by their median. Zero is used as the median alignment center.
\end{Description}
%
\begin{Usage}
\begin{verbatim}
medianCenterFeature(romics_object)
\end{verbatim}
\end{Usage}
%
\begin{Arguments}
\begin{ldescription}
\item[\code{romics\_object}] has to be an romics\_object created using romicsCreateObject()
\end{ldescription}
\end{Arguments}
%
\begin{Details}
Median normalize within each feature so that the median will be zero centered
\end{Details}
%
\begin{Value}
This function returns the transformed romics\_object with updated data layer.
\end{Value}
%
\begin{Author}
Geremy Clair
\end{Author}
\HeaderA{medianCenterSample}{medianCenterSample()}{medianCenterSample}
%
\begin{Description}
Normalizes the samples by their median. Zero is used as the median alignment center.
\end{Description}
%
\begin{Usage}
\begin{verbatim}
medianCenterSample(romics_object)
\end{verbatim}
\end{Usage}
%
\begin{Arguments}
\begin{ldescription}
\item[\code{romics\_object}] has to be an romics\_object created using romicsCreateObject()
\end{ldescription}
\end{Arguments}
%
\begin{Details}
Median normalize within each sample the median will be zero centered
\end{Details}
%
\begin{Value}
This function returns the transformed romics\_object with updated data layer.
\end{Value}
%
\begin{Author}
Geremy Clair
\end{Author}
\HeaderA{medianNormFactor}{medianNormFactor()}{medianNormFactor}
%
\begin{Description}
Normalizes the samples by their median within a given factor. The median of the median within this factor will be used as factor-specific median center.
\end{Description}
%
\begin{Usage}
\begin{verbatim}
medianNormFactor(romics_object, factor = "main")
\end{verbatim}
\end{Usage}
%
\begin{Arguments}
\begin{ldescription}
\item[\code{romics\_object}] has to be an romics\_object created using romicsCreateObject()

\item[\code{factor}] Character. Factor name from romics\_object to use for grouping. Use "main" for the main factor. Default: "main"
\end{ldescription}
\end{Arguments}
%
\begin{Details}
median normalize the samples within a given factor the median of the median of this factor will be used as new factor median.
\end{Details}
%
\begin{Value}
This function returns the transformed romics\_object with updated data layer.
\end{Value}
%
\begin{Author}
Geremy Clair
\end{Author}
\HeaderA{medianNormSample}{medianNormSample()}{medianNormSample}
%
\begin{Description}
Normalizes the samples by their median. The median of the medians of all the samples is used as the alignment point.
\end{Description}
%
\begin{Usage}
\begin{verbatim}
medianNormSample(romics_object)
\end{verbatim}
\end{Usage}
%
\begin{Arguments}
\begin{ldescription}
\item[\code{romics\_object}] has to be an romics\_object created using romicsCreateObject()
\end{ldescription}
\end{Arguments}
%
\begin{Details}
Median normalize within each sample the median of each sample median will be used as alignment point. If you waht to center the median at 0 please use the function medianCenterSample().
\end{Details}
%
\begin{Value}
This function returns the transformed romics\_object with updated data layer.
\end{Value}
%
\begin{Author}
Geremy Clair
\end{Author}
\HeaderA{multiFeatureMap}{multiFeatureMap()}{multiFeatureMap}
%
\begin{Description}
Creates a spatial map displaying multiple features from a romics\_object simultaneously, with each feature represented by different colors and intensities mapped to transparency.
\end{Description}
%
\begin{Usage}
\begin{verbatim}
multiFeatureMap(
  romics_object,
  features = NULL,
  colors = NULL,
  scale_intensities = TRUE,
  ROI = NULL,
  point_size = 1,
  alpha_range = c(0.1, 0.8),
  alpha_limits = NULL,
  facet = FALSE,
  na_value = "gray90",
  plotly = FALSE,
  title = NULL,
  white_background = TRUE,
  show_point_locations = FALSE
)
\end{verbatim}
\end{Usage}
%
\begin{Arguments}
\begin{ldescription}
\item[\code{romics\_object}] A romics\_object created using romicsCreateObject(), the metadata must contain 'x' and 'y' coordinates.

\item[\code{features}] A character vector containing the feature names to plot. Must match row names in the data.

\item[\code{colors}] A character vector of colors to represent different features. If NULL, uses default color palette. Default: NULL

\item[\code{scale\_intensities}] Logical. Whether to scale feature values (recommended for multiple features). Default: TRUE

\item[\code{ROI}] A numeric vector of length 3: c(x\_min, y\_min, distance) defining the region of interest. If NULL, uses full data range. Default: NULL

\item[\code{point\_size}] Numeric. Size of the points. Default: 1

\item[\code{alpha\_range}] Numeric vector of length 2. Min and max transparency values c(min, max). Values between 0-1. Default: c(0.1, 0.8)

\item[\code{alpha\_limits}] Numeric vector of length 2. Intensity values corresponding to min/max transparency. If NULL, uses data range. Default: NULL

\item[\code{facet}] Logical. Whether to create separate panels for each feature. Default: FALSE

\item[\code{na\_value}] Color for NA values. Default: "gray90"

\item[\code{plotly}] Logical. If TRUE, creates an interactive plotly plot. Default: FALSE

\item[\code{title}] Character. Plot title. If NULL, uses auto-generated title. Default: NULL

\item[\code{white\_background}] Logical. If TRUE, uses white background with black grid. If FALSE, uses black background with white grid. Default: TRUE

\item[\code{show\_point\_locations}] Logical. Whether to show background points indicating all point locations. Default: FALSE
\end{ldescription}
\end{Arguments}
%
\begin{Details}
This function overlays multiple features on a single spatial map, using color to distinguish features and transparency to represent intensity values. Feature names with spaces and special characters are properly handled.
\end{Details}
%
\begin{Value}
Returns either a ggplot2 or plotly interactive plot showing the spatial distribution of multiple features.
\end{Value}
%
\begin{Author}
Geremy Clair
\end{Author}
\HeaderA{multipleFeatureComparisonPlot}{multipleFeatureComparisonPlot()}{multipleFeatureComparisonPlot}
%
\begin{Description}
Creates a single plot comparing multiple features from a romics\_object across groups.
\end{Description}
%
\begin{Usage}
\begin{verbatim}
multipleFeatureComparisonPlot(
  romics_object,
  features,
  point_type = "scatter",
  factor = "main",
  group_by = "feature",
  normalize = FALSE,
  group_spacing = 0.5,
  point_size = 2,
  element_width = 0.7,
  element_alpha = 0.7,
  title = "Feature Comparison",
  feature_labels = NULL,
  x_axis_label = NULL,
  y_axis_label = "Intensity",
  legend_title = NULL,
  palette = NULL,
  show_dividers = TRUE,
  divider_color = "gray70",
  divider_linetype = "dashed",
  group_order = NULL
)
\end{verbatim}
\end{Usage}
%
\begin{Arguments}
\begin{ldescription}
\item[\code{romics\_object}] A romics\_object created using romicsCreateObject().

\item[\code{features}] Character vector. Names of features to plot. Must match exactly (case-sensitive).

\item[\code{point\_type}] Character. How to represent individual data points: "scatter", "boxplot", "violin", "bar". Default: "scatter"

\item[\code{factor}] Character. Factor name from romics\_object to use for grouping. Use "main" for the main factor. Default: "main"

\item[\code{group\_by}] Character. How to organize the plot: "feature" (features together) or "factor" (factors together). Default: "feature"

\item[\code{normalize}] Logical or character. Whether/how to normalize feature values:
FALSE (no normalization), "zscore", "minmax", or "percent\_max". Default: FALSE

\item[\code{group\_spacing}] Numeric. Space between groups (features or factors). Default: 0.5

\item[\code{point\_size}] Numeric. Size of points for scatter representation. Default: 2

\item[\code{element\_width}] Numeric. Width of boxplots/violins. Default: 0.7

\item[\code{element\_alpha}] Numeric. Transparency for elements (0-1). Default: 0.7

\item[\code{title}] Character. Plot title. Default: "Feature Comparison"

\item[\code{feature\_labels}] Character vector. Custom labels for features (same length as features). Default: NULL (use feature names)

\item[\code{x\_axis\_label}] Character. Label for the x-axis. Default: NULL (auto-set based on group\_by)

\item[\code{y\_axis\_label}] Character. Label for the y-axis. Default: "Intensity" (changes based on normalization)

\item[\code{legend\_title}] Character. Title for the legend. Default: NULL (auto-set based on group\_by)

\item[\code{palette}] Character vector. Custom color palette for features/groups. Default: NULL (uses built-in palette)

\item[\code{show\_dividers}] Logical. Whether to show dividing lines between groups. Default: TRUE

\item[\code{divider\_color}] Character. Color for the dividing lines. Default: "gray70"

\item[\code{divider\_linetype}] Character. Line type for dividers: "solid", "dashed", etc. Default: "dashed"
\end{ldescription}
\end{Arguments}
%
\begin{Details}
This function allows plotting multiple features within the same plot for direct comparison.
Features can be normalized to enable comparison on the same scale. Various representation types are
supported for individual data points, including scatter, boxplots, violins, and bars.
\end{Details}
%
\begin{Value}
A ggplot2 object
\end{Value}
%
\begin{Author}
Your Name
\end{Author}
\HeaderA{multipleFeaturePlot}{multipleFeaturePlot()}{multipleFeaturePlot}
%
\begin{Description}
Creates a grid of feature plots from a romics\_object with optional statistical comparison brackets.
\end{Description}
%
\begin{Usage}
\begin{verbatim}
multipleFeaturePlot(
  romics_object,
  features,
  plot_type = "jb",
  factor = "main",
  ncol = NULL,
  limits = "free",
  title = NULL,
  pval = "none",
  y_bracket_pos = 0.5,
  test_priority = c("Ttest", "ttest", "wilcoxon", "glm_binomial", "ANOVA"),
  significance_threshold = 0.05,
  feature_titles = TRUE,
  legend_position = "bottom"
)
\end{verbatim}
\end{Usage}
%
\begin{Arguments}
\begin{ldescription}
\item[\code{romics\_object}] A romics\_object created using romicsCreateObject().

\item[\code{features}] Character vector. Names of features to plot. Must match exactly (case-sensitive).

\item[\code{plot\_type}] Character. Plot type: "jitter", "boxplot", "violin", "jb" (jitter+boxplot), or "jv" (jitter+violin). Default: "jb"

\item[\code{factor}] Character. Factor name from romics\_object to use for grouping. Use "main" for the main factor. Default: "main"

\item[\code{ncol}] Integer. Number of columns in the plot grid. Default: NULL (automatic calculation)

\item[\code{limits}] Numeric vector of length 2 or "free". Y-axis limits c(min, max) or "free" for variable scales. Default: "free"

\item[\code{title}] Character. Overall plot title. Default: NULL (no title)

\item[\code{pval}] Character. P-value type for statistical brackets: "none", "p", or "padj". Default: "none"

\item[\code{y\_bracket\_pos}] Numeric. Vertical spacing between statistical brackets. Default: 0.5

\item[\code{test\_priority}] Character vector. Priority order for statistical tests. Default: c("Ttest", "ttest", "wilcoxon", "glm\_binomial", "ANOVA")

\item[\code{significance\_threshold}] Numeric. P-value threshold for highlighting significant results in red. Default: 0.05

\item[\code{feature\_titles}] Logical. Whether to display feature names as subplot titles. Default: TRUE

\item[\code{legend\_position}] Character. Position of the legend: "right", "bottom", "none", etc. Default: "bottom"
\end{ldescription}
\end{Arguments}
%
\begin{Details}
The function creates a grid of plots, each showing a different feature from the romics\_object.
Statistical comparisons are added as brackets if requested. The grid layout can be customized.
\end{Details}
%
\begin{Value}
A ggplot2 object with multiple feature plots arranged in a grid
\end{Value}
%
\begin{Author}
Brittney Gorman
\end{Author}
\HeaderA{pFrequencyPlot}{pFrequencyPlot()}{pFrequencyPlot}
%
\begin{Description}
makes a frequency plot of the pvalues and adjustedpvalues (columns of the statistics layer ending by '\_p' and '\_padj')
\end{Description}
%
\begin{Usage}
\begin{verbatim}
pFrequencyPlot(romics_object, p_columns = "all", p = 0.05, bin_width = 0.01)
\end{verbatim}
\end{Usage}
%
\begin{Arguments}
\begin{ldescription}
\item[\code{romics\_object}] A romics\_object created with the function romicsCreateObject()

\item[\code{p\_columns}] 'all' by default, otherwise it can be a text vector containing the columns to be plotted.

\item[\code{p}] indicate the target pvalue to be plotted with a red dotted bar.

\item[\code{bin\_width}] numeric vector, by default 0.01 indicate the width of the frequency bins
\end{ldescription}
\end{Arguments}
%
\begin{Details}
plot all or a specified list of pvalue and adjusted pvalues frequency plots
\end{Details}
%
\begin{Value}
returns one or multiple plots
\end{Value}
%
\begin{Author}
Geremy Clair
\end{Author}
\HeaderA{plotMap}{plotMap()}{plotMap}
%
\begin{Description}
This function enables to create a spatial map of the features contained in an romics\_object. It can either represent segments contained in a specific factor or features data. To ensure that the maps conserve appropriate proportionality all maps generated using this function are squares.
\end{Description}
%
\begin{Usage}
\begin{verbatim}
plotMap(
  romics_object,
  factor = NULL,
  factor_colors = ROP_colors,
  ROI = NULL,
  feature = NULL,
  point_size = 1,
  point_alpha = 1,
  scale_intensities = TRUE,
  scale_colors = viridis::viridis(3),
  scale_limits = NULL,
  na_value = "gray20",
  plotly = FALSE
)
\end{verbatim}
\end{Usage}
%
\begin{Arguments}
\begin{ldescription}
\item[\code{romics\_object}] A romics\_object created using romicsCreateObject(), the metadata has to contain at least a 'x' and 'y' coordinate in the factor layer in order to plot the map

\item[\code{factor}] The rowname of a metadata factor to be used for the plotting of regions associated with a given factor, if a factor is selected the feature map won't be shown.

\item[\code{factor\_colors}] A character vector containing a list of colors to be used to represent the different levels of the factor. Default romics\_colors will be used if not specified.

\item[\code{ROI}] A numeric vector containing 3 parameters, c(i,ii,iii): (i) 'x axis minimum' indicating the starting of the map on the x axis, (ii) 'y axis minimum' indicates the starting of the map on the y axis, and (iii) the lateral distance indicating the width and height of the map, to ensure that the maps conserve appropriate proportionality all maps generated using this function are squares.

\item[\code{feature}] A character vector containing the feature name. The feature name has to be the rowname of the 'romics\_object\$data' layer.

\item[\code{point\_size}] A numeric vector indicating the size of the points generated for the map (1 by default)

\item[\code{point\_alpha}] A numeric vector indicating the alpha of the points generated for the map (has to be comprised between 0 and 1, will be 1 = opaque by default).

\item[\code{scale\_intensities}] Boolean. Whether to scale feature intensities. Default: TRUE

\item[\code{scale\_colors}] A character vector that indicates the colors that should be used for intensity maps by default the viridis colors will be used

\item[\code{scale\_limits}] A numerical vector that indicates the limits to be used for the color\_scale, the function will automatically use the "squish" option that will make the points below the minimum or above the maximum appear with the extreme value of the scale.

\item[\code{na\_value}] A color that indicates how NA values should be displayed when a color scale is used.

\item[\code{plotly}] Boolean. If TRUE, generates an interactive plotly plot instead of a static ggplot. Default: FALSE
\end{ldescription}
\end{Arguments}
%
\begin{Value}
This function generate spatial maps of spatial omics data as ggplot2 or plotly interactive plots.
\end{Value}
%
\begin{Author}
Geremy Clair
\end{Author}
\HeaderA{quantileNormSample}{quantileNormSample()}{quantileNormSample}
%
\begin{Description}
Quantile-normalizes the samples.
\end{Description}
%
\begin{Usage}
\begin{verbatim}
quantileNormSample(romics_object)
\end{verbatim}
\end{Usage}
%
\begin{Arguments}
\begin{ldescription}
\item[\code{romics\_object}] has to be an romics\_object created using romicsCreateObject()
\end{ldescription}
\end{Arguments}
%
\begin{Details}
Quantile-normalizes within each sample.
\end{Details}
%
\begin{Value}
This function returns the transformed romics\_object with updated data layer.
\end{Value}
%
\begin{Author}
Geremy Clair
\end{Author}
\HeaderA{resetRomicsObject}{resetRomicsObject()}{resetRomicsObject}
%
\begin{Description}
Reset a romics\_object to its original state using createRomicsObject and preserving steps/UUID
\end{Description}
%
\begin{Usage}
\begin{verbatim}
resetRomicsObject(romics_object, main_factor = NULL)
\end{verbatim}
\end{Usage}
%
\begin{Arguments}
\begin{ldescription}
\item[\code{romics\_object}] A romics\_object created using createRomicsObject()

\item[\code{main\_factor}] Character string indicating the metadata row to use as main factor. If not provided, the first row will be used.
\end{ldescription}
\end{Arguments}
%
\begin{Details}
Restores the romics\_object to its original state by recreating it with createRomicsObject
and then restoring the original steps (up to createRomicsObject) and UUID
\end{Details}
%
\begin{Value}
A reset romics\_object in its original state
\end{Value}
%
\begin{Author}
Geremy Clair
\end{Author}
\HeaderA{romicsAddDependency}{romicsAddDependency()}{romicsAddDependency}
%
\begin{Description}
Adds a package to the list of dependencies of the romics\_object. Enables developers to add automatically a dependency when their function has been applied by the user on their romics\_object.
\end{Description}
%
\begin{Usage}
\begin{verbatim}
romicsAddDependency(
  romics_object,
  new_dependency = c("package_1", "package_2")
)
\end{verbatim}
\end{Usage}
%
\begin{Arguments}
\begin{ldescription}
\item[\code{romics\_object}] A romics\_object created using createRomicsObject()

\item[\code{new\_dependency}] The name of one or more R package
\end{ldescription}
\end{Arguments}
%
\begin{Value}
This function will return an romics\_object updated with the new dependency.
\end{Value}
%
\begin{Author}
Geremy Clair
\end{Author}
\HeaderA{romicsANOVA}{romicsANOVA()}{romicsANOVA}
%
\begin{Description}
Performs the ANOVA for each variable contained in the data layer of the romics\_object. The factor of the romics\_object to be used for the analysis can be selected.
\end{Description}
%
\begin{Usage}
\begin{verbatim}
romicsANOVA(romics_object, padj = TRUE, padj_method = "BH", factor = "main")
\end{verbatim}
\end{Usage}
%
\begin{Arguments}
\begin{ldescription}
\item[\code{romics\_object}] A romics\_object created with the function romicsCreateObject(),

\item[\code{padj}] Boolean indincating wheter to perform or not adjustment of pvalues

\item[\code{padj\_method}] correction method. Must be in  c("holm", "hochberg", "hommel", "bonferroni", "BH", "BY", "fdr")

\item[\code{factor}] A character string indicating the factor to use for the test, the list of the available factor can be obtained by using the function romics\_factors(), if missing the function will use the main factor of the object.
\end{ldescription}
\end{Arguments}
%
\begin{Details}
perform an ANOVA for each variable of the romics\_object the factor used will be the main\_factor of the romics\_object unless specified differently.
\end{Details}
%
\begin{Value}
an romics\_object with the statistical layer containing the newly generated ANOVA columns
\end{Value}
%
\begin{Author}
Geremy Clair
\end{Author}
\HeaderA{romicsAttributeUUID}{romicsAttributeUUID()}{romicsAttributeUUID}
%
\begin{Description}
Add or update a UUID to a romics\_object
\end{Description}
%
\begin{Usage}
\begin{verbatim}
romicsAttributeUUID(romics_object, force_new = FALSE)
\end{verbatim}
\end{Usage}
%
\begin{Arguments}
\begin{ldescription}
\item[\code{romics\_object}] A romics\_object to which a UUID will be added or updated

\item[\code{force\_new}] Logical, if TRUE will generate a new UUID even if one already exists (default: FALSE)
\end{ldescription}
\end{Arguments}
%
\begin{Value}
The romics\_object with a UUID added to the "uuid" layer
\end{Value}
%
\begin{Author}
Geremy Clair
\end{Author}
%
\begin{Examples}
\begin{ExampleCode}
## Not run: 
# Add UUID to existing romics_object
romics_object <- romicsAttributeUUID(romics_object)

# Force generation of new UUID
romics_object <- romicsAttributeUUID(romics_object, force_new = TRUE)

## End(Not run)
\end{ExampleCode}
\end{Examples}
\HeaderA{romicsBatchCorrection}{romicsBatchCorrection()}{romicsBatchCorrection}
%
\begin{Description}
Performs the sva::ComBat() batch correction on the data layer of the romics\_object. The data layer must not contain missing values and the factor utilized will be the one used for the correction.
\end{Description}
%
\begin{Usage}
\begin{verbatim}
romicsBatchCorrection(
  romics_object,
  batch_factor,
  method = "ComBat",
  verbose = TRUE,
  ...
)
\end{verbatim}
\end{Usage}
%
\begin{Arguments}
\begin{ldescription}
\item[\code{romics\_object}] A romics\_object created using romicsCreateObject()

\item[\code{batch\_factor}] A character string specifying the factor contained in the romics\_object that will serve as batch covariate. To obtain the list of factors, use romicsFactorNames()

\item[\code{method}] A character string that must be either 'ComBat' or 'mean.only' to indicate how the ComBat function should be run.

\item[\code{verbose}] Logical. If TRUE, prints detailed information about the batch correction process. Default: TRUE

\item[\code{...}] Additional parameters passed to sva::ComBat(). See sva::ComBat() documentation for more details.
\end{ldescription}
\end{Arguments}
%
\begin{Details}
This function performs ComBat batch correction on a romics\_object. It can be performed using the ComBat method or using a mean.only method. The batch factor is automatically converted to the appropriate format for ComBat. See sva::ComBat() documentation for more details.
\end{Details}
%
\begin{Value}
This function returns a transformed romics\_object with batch-corrected data.
\end{Value}
%
\begin{Author}
Geremy Clair
\end{Author}
\HeaderA{romicsCalculatedStats}{romicsCalculatedStats()}{romicsCalculatedStats}
%
\begin{Description}
Indicates the stat columns calculated for the romics\_object
\end{Description}
%
\begin{Usage}
\begin{verbatim}
romicsCalculatedStats(romics_object)
\end{verbatim}
\end{Usage}
%
\begin{Arguments}
\begin{ldescription}
\item[\code{romics\_object}] A romics\_object created using createRomicsObject()
\end{ldescription}
\end{Arguments}
%
\begin{Value}
This function will return character vector containing the stat columns previously calculated for a romics\_object, if no stats were previously calculated an error message will be displayed
\end{Value}
%
\begin{Author}
Geremy Clair
\end{Author}
\HeaderA{romicsChangeColors}{romicsChangeColors()}{romicsChangeColors}
%
\begin{Description}
Change the colors of an existing romics\_object based on the main factor levels (performance optimized)
\end{Description}
%
\begin{Usage}
\begin{verbatim}
romicsChangeColors(romics_object, new_colors)
\end{verbatim}
\end{Usage}
%
\begin{Arguments}
\begin{ldescription}
\item[\code{romics\_object}] A romics\_object created using createRomicsObject()

\item[\code{new\_colors}] A character vector containing the new colors you want to use. The colors will be applied in the order of the main factor levels.
\end{ldescription}
\end{Arguments}
%
\begin{Details}
This function updates the colors in a romics\_object by replacing the existing color scheme with new colors. Uses vectorized operations for optimal performance with large datasets.
\end{Details}
%
\begin{Value}
This function returns the romics\_object with updated colors in the metadata, colors, and custom\_colors layers
\end{Value}
%
\begin{Author}
Geremy Clair
\end{Author}
\HeaderA{romicsChangeFactor}{romicsChangeFactor()}{romicsChangeFactor}
%
\begin{Description}
Change the main factor of the romics\_object
\end{Description}
%
\begin{Usage}
\begin{verbatim}
romicsChangeFactor(romics_object, main_factor = "none")
\end{verbatim}
\end{Usage}
%
\begin{Arguments}
\begin{ldescription}
\item[\code{romics\_object}] A object created using the function createRomicsObject().

\item[\code{main\_factor}] Either 'none' OR a factor from the romics\_object (corresponding to a row from the original metadata file), the list of factors from an romics object can be obtained using the function romicsFactorNames().
\end{ldescription}
\end{Arguments}
%
\begin{Details}
changes the main\_factor of an romics\_object and updates the colors to this new factor. Also reorders the embeddings layer if present.
\end{Details}
%
\begin{Value}
This function returns a modified romics object, please see the create\_romics\_object() documentation.
\end{Value}
%
\begin{Author}
Geremy Clair
\end{Author}
\HeaderA{romicsChangeIDs}{romicsChangeIDs()}{romicsChangeIDs}
%
\begin{Description}
Change the main IDs of the romics\_object using a specified column from the IDs layer
\end{Description}
%
\begin{Usage}
\begin{verbatim}
romicsChangeIDs(romics_object, newIDs)
\end{verbatim}
\end{Usage}
%
\begin{Arguments}
\begin{ldescription}
\item[\code{romics\_object}] A romics\_object created using createRomicsObject()

\item[\code{newIDs}] Character string. Must match exactly one of the column names of the 'IDs' layer
\end{ldescription}
\end{Arguments}
%
\begin{Details}
Changes the row names (IDs) of the data, missingdata, and statistics layers using the specified column
from the IDs layer. Always stores the conversion history with "original" and "converted" columns.
Updates the IDs table with "Original\_IDs" and "Current\_IDs" columns to enable reverse mapping.
Missing mappings are handled with "Undefined\_ID\_matched\_to\_" prefix.
\end{Details}
%
\begin{Value}
This function returns a modified romics object with updated IDs and conversion history stored.
\end{Value}
%
\begin{Author}
Geremy Clair
\end{Author}
\HeaderA{romicsChangeLevelName}{romicsChangeLevelName()}{romicsChangeLevelName}
%
\begin{Description}
Changes the name(s) of one or more levels in a factor of the romics\_object metadata layer.
\end{Description}
%
\begin{Usage}
\begin{verbatim}
romicsChangeLevelName(romics_object, factor_name, old_names, new_names)
\end{verbatim}
\end{Usage}
%
\begin{Arguments}
\begin{ldescription}
\item[\code{romics\_object}] A romics\_object created using romicsCreateObject().

\item[\code{factor\_name}] Character string. The name of the factor whose level names should be changed.

\item[\code{old\_names}] Character vector. The current name(s) of the level(s) to be changed.

\item[\code{new\_names}] Character vector. The new name(s) for the level(s). Must be the same length as old\_names.
\end{ldescription}
\end{Arguments}
%
\begin{Details}
This function allows renaming of factor levels in the metadata layer of a romics\_object.
\end{Details}
%
\begin{Value}
A romics\_object with updated factor level names in the metadata layer.
\end{Value}
%
\begin{Author}
Geremy Clair
\end{Author}
\HeaderA{romicsChangeSampleNames}{romicsChangeSampleNames()}{romicsChangeSampleNames}
%
\begin{Description}
Changes the sample identifiers using values contained in one of the factors from the romics\_object metadata.
The new sample names will be taken from the specified factor, which must contain unique values for each sample.
\end{Description}
%
\begin{Usage}
\begin{verbatim}
romicsChangeSampleNames(
  romics_object,
  factor_name,
  original_names = "keep",
  original_names_factor = "original_sample_names"
)
\end{verbatim}
\end{Usage}
%
\begin{Arguments}
\begin{ldescription}
\item[\code{romics\_object}] A romics\_object created using createRomicsObject()

\item[\code{factor\_name}] A character string specifying the name of a factor contained in the metadata of the romics\_object.
This factor must contain unique values that will become the new sample names.
To obtain the list of available factors, use romicsFactorNames().

\item[\code{original\_names}] Either 'keep' or 'drop'. Indicates what to do with the original sample names:
'keep' will store them as a new factor in the metadata, 'drop' will discard them completely. Default: 'keep'

\item[\code{original\_names\_factor}] The name to give to the new factor that will store the original sample names
(only used when original\_names = 'keep'). Default: 'original\_sample\_names'
\end{ldescription}
\end{Arguments}
%
\begin{Details}
This function enables quick renaming of sample names using values from a factor contained in the metadata.
The specified factor must contain only unique values to ensure each sample gets a distinct name.
Sample names will be updated across all layers of the romics\_object (data, metadata, missingdata, embeddings).

When original\_names = 'keep', the current sample names are preserved as a new factor in the metadata
before being replaced. When original\_names = 'drop', the original names are permanently lost.

The function works with all romics\_object layers including embeddings from PCA, UMAP, or t-SNE analyses.
\end{Details}
%
\begin{Value}
A romics\_object with updated sample names across all layers. If original\_names = 'keep',
a new factor containing the original sample names will be added to the metadata.
\end{Value}
%
\begin{Author}
Geremy Clair
\end{Author}
\HeaderA{romicsCombineFactors}{romicsCombineFactors}{romicsCombineFactors}
%
\begin{Description}
This function will use multiple factor of an romics object to create a new one.
\end{Description}
%
\begin{Usage}
\begin{verbatim}
romicsCombineFactors(
  romics_object = romics_object,
  list_of_factors = c("f1", "f2"),
  separator = ";",
  replace_missing = "unknown",
  factor_name = NULL
)
\end{verbatim}
\end{Usage}
%
\begin{Arguments}
\begin{ldescription}
\item[\code{romics\_object}] A romics\_object created using romicsCreateObject().

\item[\code{list\_of\_factors}] A character vector (length=>1) containing names of factors of the romics\_object in the desired order. The list of factor names can be identified by using the function romicsFactorNames().

\item[\code{separator}] A character vector (length=1) that indicates the separator of the different factors to be used (when multiple factors to be created are contained in a given factor)

\item[\code{replace\_missing}] A character string to replace missing values. Default: "unknown"

\item[\code{factor\_name}] A character string for the name of the new factor. If not provided, uses the collapsed names of input factors.
\end{ldescription}
\end{Arguments}
%
\begin{Details}
Generate new factor from a list of existing factors in the specified order. The new factor will be the aggregation of the different factors content in the order provided.
\end{Details}
%
\begin{Value}
A Romics\_object with a new factor. The name of the factor will be the collapsed name of the factors aggregated (separated with the <separator>), the content of the factor will be the aggregation of the different factors content in the specified order.
\end{Value}
%
\begin{Author}
Geremy Clair, Naina Beishembieva
\end{Author}
\HeaderA{romicsCombineSpatialObjects}{romicsCombineSpatialObjects()}{romicsCombineSpatialObjects}
%
\begin{Description}
Combines multiple spatial romics\_objects into a single object by arranging them in a grid layout.
This function is designed for spatial omics data where samples have x,y coordinates. Each romics\_object
is positioned in a grid with specified spacing, and coordinates are adjusted to prevent overlap.
Sample names are prefixed with object names to ensure uniqueness even when objects have identical sample names.
\end{Description}
%
\begin{Usage}
\begin{verbatim}
romicsCombineSpatialObjects(
  romics_objects,
  xMinShift = 100,
  yMinShift = 100,
  xElements = 2,
  yElements = 2,
  center_align = TRUE,
  add_object_factor = TRUE,
  preserve_factors = TRUE
)
\end{verbatim}
\end{Usage}
%
\begin{Arguments}
\begin{ldescription}
\item[\code{romics\_objects}] A named list of romics\_objects to be combined. All objects must have 'x' and 'y'
coordinates in their metadata and should contain the same features (rows in data layer).
Names will be used as prefixes for sample names.

\item[\code{xMinShift}] Numeric. Minimum spacing between objects along the x-axis. Default: 100

\item[\code{yMinShift}] Numeric. Minimum spacing between objects along the y-axis. Default: 100

\item[\code{xElements}] Integer. Number of objects to place horizontally before moving to next row. Default: 2

\item[\code{yElements}] Integer. Number of objects to place vertically. Default: 2

\item[\code{center\_align}] Logical. If TRUE (default), objects are centered within their grid positions.
If FALSE, objects are positioned at grid corner coordinates.

\item[\code{add\_object\_factor}] Logical. If TRUE (default), adds a factor indicating which
original object each sample came from.

\item[\code{preserve\_factors}] Logical. If TRUE (default), attempts to preserve all factors
from input objects. Conflicts are resolved by prefixing with object names.
\end{ldescription}
\end{Arguments}
%
\begin{Details}
This function combines spatial romics objects by:
1. Validating that all objects have compatible structure
2. Calculating grid positions for each object
3. Adjusting x,y coordinates to position objects in the grid
4. Prefixing sample names with object names to ensure uniqueness
5. Combining data, metadata, and other layers
6. Optionally preserving factors from original objects

Grid positioning:
- Objects are placed left-to-right, then bottom-to-top
- Position (0,0) is bottom-left of the grid
- Spacing ensures minimum distance between object boundaries
- Center alignment positions objects at grid cell centers

Sample name handling:
- All sample names are prefixed with their source object name
- Example: "voxel1" from object "heart" becomes "heart\_voxel1"
- This prevents conflicts when objects have identical sample names

Coordinate transformation:
- Each object's coordinates are shifted to fit in assigned grid position
- Original relative positions within each object are preserved
- Minimum bounding rectangles are calculated for proper spacing
\end{Details}
%
\begin{Value}
Returns a combined romics\_object containing all samples from input objects
with adjusted spatial coordinates and prefixed sample names. The object will contain
a new factor indicating source object (if add\_object\_factor = TRUE).
\end{Value}
%
\begin{Author}
Geremy Clair
\end{Author}
\HeaderA{romicsComplexHeatmap}{romicsComplexHeatmap()}{romicsComplexHeatmap}
%
\begin{Description}
Creates a heatmap with statistical annotations from a romics object
\end{Description}
%
\begin{Usage}
\begin{verbatim}
romicsComplexHeatmap(
  romics_object,
  comparison_prefix = NULL,
  factor_name = NULL,
  feature_list = NULL,
  p_threshold = 0.05,
  use_adjusted_p = TRUE,
  scale_data = TRUE,
  cluster_rows = TRUE,
  cluster_cols = TRUE,
  show_row_names = TRUE,
  show_column_names = TRUE,
  heatmap_colors = c("blue", "white", "red"),
  color_breaks = NULL,
  allow_missing = TRUE,
  missing_color = "grey",
  significance_color = "red"
)
\end{verbatim}
\end{Usage}
%
\begin{Arguments}
\begin{ldescription}
\item[\code{romics\_object}] A romics object containing \$data, \$statistics, and \$metadata layers

\item[\code{comparison\_prefix}] Character string of the comparison prefix (e.g., "Alveoli\_pleura\_vs\_Bronchi\_bronchiole"). If NULL, auto-detects from statistics columns

\item[\code{factor\_name}] Factor name from metadata for sample grouping. If NULL, uses main factor. Use romicsFactorNames() to see available factors

\item[\code{feature\_list}] Character vector of features to include in heatmap. If provided, creates subset using romicsFeatureSubset(). If NULL, uses all features

\item[\code{p\_threshold}] Numeric threshold for p-values (default 0.05)

\item[\code{use\_adjusted\_p}] TRUE or FALSE, if TRUE uses adjusted p-values, if FALSE uses raw p-values

\item[\code{scale\_data}] TRUE or FALSE, if TRUE scales the data by rows (z-score)

\item[\code{cluster\_rows}] TRUE or FALSE, if TRUE performs hierarchical clustering of rows

\item[\code{cluster\_cols}] TRUE or FALSE, if TRUE performs hierarchical clustering of columns

\item[\code{show\_row\_names}] TRUE or FALSE, if TRUE shows feature names on rows

\item[\code{show\_column\_names}] TRUE or FALSE, if TRUE shows sample names on columns

\item[\code{heatmap\_colors}] Character vector of colors for the heatmap. Default: c("blue", "white", "red")

\item[\code{color\_breaks}] Numeric vector of breaks for color mapping. If NULL, uses automatic breaks

\item[\code{allow\_missing}] TRUE or FALSE, if TRUE allows missing values in the heatmap

\item[\code{missing\_color}] Color for missing values in the heatmap (default: "grey")

\item[\code{significance\_color}] Color for significant results in all test annotations (default: "red")
\end{ldescription}
\end{Arguments}
%
\begin{Details}
Creates a heatmap with side annotations automatically detected from statistics columns (T-test, Wilcoxon, GLM binomial, ANOVA) and sample groups from metadata. If feature\_list is provided, automatically subsets the romics object first.
\end{Details}
%
\begin{Value}
ComplexHeatmap object
\end{Value}
%
\begin{Author}
Geremy Clair
\end{Author}
\HeaderA{romicsCountDetect}{romicsCountDetect()}{romicsCountDetect}
%
\begin{Description}
Calculates the in how many samples the features are detected and generates a table with these metrics

This function will detect in how many samples each feature is detected and create a summary data.frame from these counts
\end{Description}
%
\begin{Usage}
\begin{verbatim}
romicsCountDetect(romics_object)

romicsCountDetect(romics_object)
\end{verbatim}
\end{Usage}
%
\begin{Arguments}
\begin{ldescription}
\item[\code{romics\_object}] has to be a log transformed romics\_object created using romicsCreateObject() and transformed using the function log2transform() or log10transform()
\end{ldescription}
\end{Arguments}
%
\begin{Details}
TCalculates the in how many samples the features are detected and generates a table with these metrics in a data frame
\end{Details}
%
\begin{Value}
a data frame with the counts of proteins detected in different samples
\end{Value}
%
\begin{Author}
Geremy Clair
\end{Author}
\HeaderA{romicsCountDetectPlot}{romicsCountDetectPlot()}{romicsCountDetectPlot}
%
\begin{Description}
Calculates the in how many samples the features are detected and generates a plot with these metrics

This function will detect in how many samples each feature is detected and create a plot from these counts
\end{Description}
%
\begin{Usage}
\begin{verbatim}
romicsCountDetectPlot(romics_object, percent = TRUE)

romicsCountDetectPlot(romics_object, percent = TRUE)
\end{verbatim}
\end{Usage}
%
\begin{Arguments}
\begin{ldescription}
\item[\code{romics\_object}] has to be a log transformed romics\_object created using romicsCreateObject() and transformed using the function log2transform() or log10transform()
\end{ldescription}
\end{Arguments}
%
\begin{Details}
TCalculates the in how many samples the features are detected and generates a table with these metrics in a ggplot
\end{Details}
%
\begin{Value}
a ggplot plot
\end{Value}
%
\begin{Author}
Geremy Clair
\end{Author}
\HeaderA{romicsCreateDependencies}{romicsCreateDependencies()}{romicsCreateDependencies}
%
\begin{Description}
Creates the original dependencies of the romics\_object (only when it is created to add dependencies use the romicsAddDependency() function)
\end{Description}
%
\begin{Usage}
\begin{verbatim}
romicsCreateDependencies()
\end{verbatim}
\end{Usage}
%
\begin{Value}
This function will return a data.frame with two columns the required packages and the version at the time the code was run
\end{Value}
%
\begin{Author}
Geremy Clair
\end{Author}
\HeaderA{romicsCreateFactorFromROIs}{romicsCreateFactorFromROIs()}{romicsCreateFactorFromROIs}
%
\begin{Description}
Creates metadata factors from 2D Regions of Interest (ROI) definitions.
This function takes a list of ROI coordinates and creates a new factor in the romics\_object metadata
indicating which samples fall within each ROI. ROIs are defined as rectangular regions using
minimum coordinates and width parameters.
\end{Description}
%
\begin{Usage}
\begin{verbatim}
romicsCreateFactorFromROIs(
  romics_object,
  list_ROI,
  factor_name = "ROIs",
  separator = ";",
  outside_label = "Outside"
)
\end{verbatim}
\end{Usage}
%
\begin{Arguments}
\begin{ldescription}
\item[\code{romics\_object}] A romics\_object created using romicsCreateObject() that contains 'x' and 'y'
coordinates in its metadata layer

\item[\code{list\_ROI}] A named list where each element is a numeric vector of length 3: c(X, Y, W).
X = minimum x-coordinate, Y = minimum y-coordinate, W = width of the square ROI.
The maximum coordinates are calculated as X+W and Y+W.
Example: list(ROI1 = c(10, 20, 50), ROI2 = c(100, 150, 30))

\item[\code{factor\_name}] A character string specifying the name for the new factor to be created
in the metadata layer. Default: "ROIs"

\item[\code{separator}] A character string used to separate multiple ROI names when a sample
falls within multiple ROIs. Default: ";"

\item[\code{outside\_label}] A character string to label samples that don't fall within any ROI.
If NULL, these samples will be left as NA. Default: "Outside"
\end{ldescription}
\end{Arguments}
%
\begin{Details}
This function creates rectangular ROIs and assigns samples to them based on their
x and y coordinates stored in the metadata. Each ROI is defined by:
- X: minimum x-coordinate
- Y: minimum y-coordinate
- W: width (and height) of the square region
- Maximum coordinates: (X+W, Y+W)

Samples can belong to multiple ROIs if their coordinates fall within overlapping regions.
In such cases, ROI names are concatenated using the specified separator.

The function requires 'x' and 'y' factors to exist in the romics\_object metadata.
These coordinates are typically present in spatial omics datasets like imaging mass spectrometry,
spatial transcriptomics, or multiplexed immunofluorescence data.
\end{Details}
%
\begin{Value}
Returns the romics\_object with a new factor added to the metadata layer containing
ROI assignments for each sample. Samples outside all ROIs will be labeled according
to outside\_label parameter.
\end{Value}
%
\begin{Author}
Geremy Clair
\end{Author}
\HeaderA{romicsCreateROIsFromFactor}{romicsCreateROIsFromFactor()}{romicsCreateROIsFromFactor}
%
\begin{Description}
Creates a list of ROI definitions from factor levels in romics\_object metadata.
This function generates rectangular ROIs that encompass all samples belonging to each level
of a specified factor. Each ROI is defined as the minimum bounding rectangle containing
all samples of that factor level.
\end{Description}
%
\begin{Usage}
\begin{verbatim}
romicsCreateROIsFromFactor(
  romics_object,
  factor_name = NULL,
  padding = 0,
  min_roi_size = 1,
  check_overlap = TRUE,
  square_rois = TRUE
)
\end{verbatim}
\end{Usage}
%
\begin{Arguments}
\begin{ldescription}
\item[\code{romics\_object}] A romics\_object created using romicsCreateObject() that contains 'x', 'y'
coordinates and the specified factor in its metadata layer

\item[\code{factor\_name}] A character string specifying the name of the factor to use for ROI creation.
If NULL (default), uses the main\_factor of the romics\_object.

\item[\code{padding}] A numeric value specifying additional padding to add around each ROI boundary.
This ensures ROIs don't touch the exact sample boundaries. Default: 0

\item[\code{min\_roi\_size}] A numeric value specifying the minimum width/height for any ROI.
ROIs smaller than this will be expanded symmetrically. Default: 1

\item[\code{check\_overlap}] Logical. If TRUE (default), checks whether generated ROIs are fully distinct
and issues warnings for overlapping ROIs.

\item[\code{square\_rois}] Logical. If TRUE (default), creates square ROIs where width = height = max(width, height).
If FALSE, creates rectangular ROIs with independent width and height.
\end{ldescription}
\end{Arguments}
%
\begin{Details}
This function creates ROIs by:
1. Extracting x,y coordinates for each factor level
2. Computing the bounding rectangle for each level's samples
3. Adding optional padding and enforcing minimum size
4. Checking for overlaps between ROIs if requested

Each ROI is returned as a vector c(X, Y, W) where:
- X: minimum x-coordinate
- Y: minimum y-coordinate
- W: width (and height for square ROIs)

The function requires 'x' and 'y' factors in the metadata and will fail if
any factor level has no samples or samples with missing coordinates.

Overlap checking compares all ROI pairs and reports any overlapping regions,
which might indicate that the factor levels are not spatially distinct.
\end{Details}
%
\begin{Value}
Returns a named list where each element is a numeric vector c(X, Y, W) defining
an ROI for the corresponding factor level. Names correspond to factor levels.
\end{Value}
%
\begin{Author}
Geremy Clair
\end{Author}
%
\begin{SeeAlso}
\code{\LinkA{romicsCreateFactorFromROIs}{romicsCreateFactorFromROIs}} to create factors from ROI definitions
\end{SeeAlso}
%
\begin{Examples}
\begin{ExampleCode}
## Not run: 
# Create ROIs from main factor
roi_list <- romicsCreateROIsFromFactor(romics_obj)

# Create ROIs from specific factor with padding
roi_list <- romicsCreateROIsFromFactor(romics_obj,
                                       factor_name = "tissue_regions",
                                       padding = 5)

# Create rectangular ROIs without overlap checking
roi_list <- romicsCreateROIsFromFactor(romics_obj,
                                       factor_name = "cell_types",
                                       square_rois = FALSE,
                                       check_overlap = FALSE)

# Use the generated ROIs to create a new factor
romics_obj <- romicsCreateFactorFromROIs(romics_obj, list_ROI = roi_list)

## End(Not run)
\end{ExampleCode}
\end{Examples}
\HeaderA{romicsCVs}{romicsCVs()}{romicsCVs}
%
\begin{Description}
Calculates the CVs (works properly for non logged romics\_objects only) for each level of a metadata factor. This function can also be used to plot the CV boxplot and/or barplot on demand. The CVs are generated in a separate object and won’t be included in the statistics layer.
\end{Description}
%
\begin{Usage}
\begin{verbatim}
romicsCVs(romics_object, factor = "main", plot = "all")
\end{verbatim}
\end{Usage}
%
\begin{Arguments}
\begin{ldescription}
\item[\code{romics\_object}] A romics\_object created using romicsCreateObject()

\item[\code{factor}] A string indicating a factor from the romics\_object. The factors usable for a given romics\_object names can be obtain using the function romicsFactorNames(romics\_object)

\item[\code{plot}] can be either FALSE, 'boxplot', 'barplot', or 'all'
\end{ldescription}
\end{Arguments}
%
\begin{Details}
This function will calculate the CVs, the percentage CVs, and on demand will plot the barplot and/or the boxplot of the CVs per sampletype. The CVs are calculated on unlogged romics\_object only, if the data was logged transformed, it will be unlogged prior to perform the calculations.
\end{Details}
%
\begin{Author}
Geremy Clair
\end{Author}
\HeaderA{romicsDeleteRowMissing}{romicsDeleteRowMissing()}{romicsDeleteRowMissing}
%
\begin{Description}
Removes the rows that have any missing values from the romics\_object
\end{Description}
%
\begin{Usage}
\begin{verbatim}
romicsDeleteRowMissing(romics_object)
\end{verbatim}
\end{Usage}
%
\begin{Arguments}
\begin{ldescription}
\item[\code{romics\_object}] has to be an romics\_object created using romicsCreateObject()
\end{ldescription}
\end{Arguments}
%
\begin{Details}
This function will remove rows that have any missing values in the data and missingdata layers
\end{Details}
%
\begin{Value}
The function will return a filtered romics\_object with the rows of the data and missing data object removed when appropriate.
\end{Value}
%
\begin{Author}
Geremy Clair
\end{Author}
\HeaderA{romicsExportData}{romicsExportData()}{romicsExportData}
%
\begin{Description}
Creates an exportable data frame from the romics\_object. The generated data.frame contains the processed data, the statistics and the missing status of the data on demand.
\end{Description}
%
\begin{Usage}
\begin{verbatim}
romicsExportData(
  romics_object,
  data = TRUE,
  statistics = FALSE,
  missing_data = FALSE,
  IDs = FALSE
)
\end{verbatim}
\end{Usage}
%
\begin{Arguments}
\begin{ldescription}
\item[\code{romics\_object}] A romics\_object created using createRomicsObject()

\item[\code{statistics}] boolean, has to be TRUE or FALSE. Indicates if the statistics should be exported along with the data (FALSE by default).

\item[\code{missing\_data}] boolean, has to be TRUE or FALSE. Indicates if the missing status of the data should be exported along with the data (FALSE by default).

\item[\code{IDs}] boolean, has to be TRUE or FALSE. Indicates if the IDs of the data should be exported along with the data (FALSE by default).
\end{ldescription}
\end{Arguments}
%
\begin{Value}
This function will return an data frame containing the results of the processing, the statistics and the missingness status of the data as specified by the user.
\end{Value}
%
\begin{Author}
Geremy Clair
\end{Author}
\HeaderA{romicsExtractFactor}{romicsExtractFactor()}{romicsExtractFactor}
%
\begin{Description}
Extract a factor from the romics\_object
\end{Description}
%
\begin{Usage}
\begin{verbatim}
romicsExtractFactor(romics_object, factor = "main")
\end{verbatim}
\end{Usage}
%
\begin{Arguments}
\begin{ldescription}
\item[\code{romics\_object}] A romics\_object created using createRomicsObject()

\item[\code{factor}] has to be either 'main' to indicate that the main factor has to be used (default) or a factor contained in the romics\_object, the list of factors can be obtained using the function romicsFactorNames()
\end{ldescription}
\end{Arguments}
%
\begin{Details}
This function allows to quickly extract the content of a factor present in the romics\_object.
\end{Details}
%
\begin{Value}
a factor from an romics\_object
\end{Value}
%
\begin{Author}
Geremy Clair
\end{Author}
\HeaderA{romicsExtractFeature}{romicsExtractFeature()}{romicsExtractFeature}
%
\begin{Description}
Extract the data for a single feature from the romics\_object
\end{Description}
%
\begin{Usage}
\begin{verbatim}
romicsExtractFeature(romics_object, feature = "feature")
\end{verbatim}
\end{Usage}
%
\begin{Arguments}
\begin{ldescription}
\item[\code{romics\_object}] A romics\_object created using createRomicsObject()

\item[\code{feature}] has to be either a feature contained in the romics\_object, the list of features can be obtained using the function romicsFeatureNames()
\end{ldescription}
\end{Arguments}
%
\begin{Details}
This function allows to quickly extract the data for a romics\_object feature.
\end{Details}
%
\begin{Value}
a feature from an romics\_object
\end{Value}
%
\begin{Author}
Geremy Clair
\end{Author}
\HeaderA{romicsExtractSignificantFeatures}{romicsExtractSignificantFeatures()}{romicsExtractSignificantFeatures}
%
\begin{Description}
Extract the names of features that are significant based on any statistical test in the statistics layer
\end{Description}
%
\begin{Usage}
\begin{verbatim}
romicsExtractSignificantFeatures(
  romics_object,
  p_threshold = 0.05,
  padj_type = "p",
  specific_comparisons = NULL,
  test_types = NULL,
  directionality = "all",
  fc_threshold = 0
)
\end{verbatim}
\end{Usage}
%
\begin{Arguments}
\begin{ldescription}
\item[\code{romics\_object}] A romics\_object with calculated statistics

\item[\code{p\_threshold}] Numeric value for significance threshold. Default: 0.05

\item[\code{padj\_type}] Character string indicating which p-value type to use: "p" for raw p-values or "padj" for adjusted p-values. Default: "p"

\item[\code{specific\_comparisons}] Character vector of specific comparisons to consider (e.g., "Alport\_vs\_Ctrl"). If NULL, considers all comparisons. Default: NULL

\item[\code{test\_types}] Character vector of test types to consider (e.g., "Ttest", "Wilcox\_test", "glmBinomialTest"). If NULL, considers all tests. Default: NULL

\item[\code{directionality}] Character string indicating which direction of change to consider: "all" (both up and down), "up" (upregulated only), "down" (downregulated only). Default: "all"

\item[\code{fc\_threshold}] Numeric value for fold change threshold (used with directionality). For log-transformed data, this is the log fold change threshold. Default: 0 (no fold change filtering)
\end{ldescription}
\end{Arguments}
%
\begin{Details}
This function searches through statistical test columns in the statistics layer and identifies
features that meet the significance criteria and directionality requirements in any test.
For t-test and Wilcox test: up = log(FC) > 0 (or FC > 1), down = log(FC) < 0 (or FC < 1).
For GLM binomial test: up = directionality = 1, down = directionality = -1.
It automatically detects p-value columns based on the naming convention and the specified padj\_type.
\end{Details}
%
\begin{Value}
Character vector of significant feature names meeting the directionality criteria
\end{Value}
%
\begin{Author}
Geremy Clair
\end{Author}
\HeaderA{romicsFactorFromFeature}{romicsFactorFromFeature}{romicsFactorFromFeature}
%
\begin{Usage}
\begin{verbatim}
romicsFactorFromFeature(
  romics_object,
  feature,
  method = "presence",
  value = NULL,
  percentage = 50,
  n_groups = 2,
  factor_name = NULL,
  labels = NULL,
  na_label = "absent",
  na_as_low = FALSE,
  strict_inequality = FALSE,
  include_value_in_name = TRUE,
  verbose = TRUE
)
\end{verbatim}
\end{Usage}
%
\begin{Arguments}
\begin{ldescription}
\item[\code{romics\_object}] A romics\_object created using romicsCreateObject().

\item[\code{feature}] A character string specifying the name of the feature to use. Must be present in rownames of the data layer.

\item[\code{method}] Classification method: "presence" (present/absent), "value" (above/below a fixed value), "percentage" (above/below a percentile). Default: "presence".

\item[\code{value}] Numeric value for splitting when method = "value", Required when method = "value".

\item[\code{percentage}] Numeric between 0 and 100 for percentage method (e.g., 75 for top 25

\bsl{}itemn\_groupsInteger for number of groups when using percentage method with equal splits (2, 3, or 4). Default: 2. When > 2, overrides percentage parameter.

\bsl{}itemfactor\_nameCharacter string for the name of the new factor. If NULL, auto-generated based on method. Default: NULL.

\bsl{}itemlabelsCharacter vector of custom labels. Length depends on n\_groups. If NULL, auto-generated. For binary splits, order is c(high, low). Default: NULL.

\bsl{}itemna\_labelCharacter string for samples where the feature is NA/not detected. Default: "absent".

\bsl{}itemna\_as\_lowLogical. For binary methods, should NA values be treated as "low" rather than a separate category? Default: FALSE.

\bsl{}itemstrict\_inequalityLogical. Use strict inequality (>) or >= for high group. Default: FALSE (uses >).

\bsl{}iteminclude\_value\_in\_nameLogical. Include value/percentage in auto-generated factor name. Default: TRUE.

\bsl{}itemverboseLogical. Print detailed summary. Default: TRUE.

\end{ldescription}

Returns the romics\_object with the newly created factor added to its metadata.


Creates a new factor in the metadata based on a feature's intensity using various classification methods.


Geremy Clair

\end{Arguments}
\HeaderA{romicsFactorNames}{romicsFactorNames()}{romicsFactorNames}
%
\begin{Description}
Indicates the list of factor names from the romics\_object
\end{Description}
%
\begin{Usage}
\begin{verbatim}
romicsFactorNames(romics_object)
\end{verbatim}
\end{Usage}
%
\begin{Arguments}
\begin{ldescription}
\item[\code{romics\_object}] A romics\_object created using createRomicsObject()
\end{ldescription}
\end{Arguments}
%
\begin{Details}
This function allows to quickly get a vector containing all the factor names present in a romics\_object
\end{Details}
%
\begin{Value}
A character vector containing the list of factor contained in an romics\_object
\end{Value}
%
\begin{Author}
Geremy Clair
\end{Author}
\HeaderA{romicsFactorProportionTestPairwise}{romicsFactorProportionTestPairwise}{romicsFactorProportionTestPairwise}
%
\begin{Description}
Tests whether the proportions of each level of a given factor differ
significantly between pairs of conditions. Uses an omnibus chi-square test on the
full contingency table for each pairwise comparison, then derives per-level
directionality from standardized Pearson residuals.
\end{Description}
%
\begin{Usage}
\begin{verbatim}
romicsFactorProportionTestPairwise(
  romics_object,
  test_factor,
  condition_factor = "condition",
  method = "chisq",
  simulate_p = TRUE,
  B = 10000,
  padj_method = "BH",
  residual_threshold = 2,
  min_effect_size = 5,
  min_fold_change = 1.5
)
\end{verbatim}
\end{Usage}
%
\begin{Arguments}
\begin{ldescription}
\item[\code{romics\_object}] An romics\_object created using romicsCreateObject().

\item[\code{test\_factor}] A character string indicating the factor whose level proportions
will be tested (e.g., "Leiden\_clust\_names", "cell\_type", "AGER\_presence").

\item[\code{condition\_factor}] A character string indicating the condition factor (e.g., "condition").

\item[\code{method}] A character string specifying the test method: "chisq" (chi-square test) or
"fisher" (Fisher's exact test via monte-carlo simulation for tables > 2x2). Default: "chisq".

\item[\code{simulate\_p}] A logical indicating whether to simulate p-values (recommended when expected
cell frequencies are low, i.e. < 5). Default: TRUE.

\item[\code{B}] Integer. Number of replicates for p-value simulation. Default: 10000.

\item[\code{padj\_method}] Correction method for multiple testing across pairwise comparisons. Default: "BH".

\item[\code{residual\_threshold}] Minimum absolute standardized Pearson residual required to flag a
level as meaningfully over/under-represented. Conventionally 2.0. Default: 2.0.

\item[\code{min\_effect\_size}] Minimum absolute difference in proportions (
descriptive filter. Default: 5.

\item[\code{min\_fold\_change}] Minimum fold change in proportions required as a descriptive filter.
\end{ldescription}
\end{Arguments}
%
\begin{Author}
Geremy Clair
\end{Author}
\HeaderA{romicsFactorStackedBarplot}{romicsFactorStackedBarplot()}{romicsFactorStackedBarplot}
%
\begin{Usage}
\begin{verbatim}
romicsFactorStackedBarplot(
  romics_object,
  split_by_factor,
  frequency_for_factor,
  title = NULL,
  xlab = NULL,
  ylab = "Percentage (%)",
  colors = NULL,
  legend_title = NULL,
  show_percentages = TRUE,
  text_size = 3,
  text_color = "black",
  show_summary_table = TRUE,
  export_summary_table = FALSE,
  name_summary_table = "summary_table_FactorStackedBarplot"
)
\end{verbatim}
\end{Usage}
%
\begin{Arguments}
\begin{ldescription}
\item[\code{romics\_object}] A romics\_object created using romicsCreateObject().

\item[\code{split\_by\_factor}] A character vector (length=1) containing a factor of the romics\_object that will define the x-axis categories. The list of factors can be identified by using the function romicsFactorNames()

\item[\code{frequency\_for\_factor}] A character vector (length=1) containing a factor of the romics\_object that will be used to calculate frequency distributions within each level of split\_by\_factor. The list of factors can be identified by using the function romicsFactorNames()

\item[\code{title}] A character string specifying the plot title. If NULL (default), an automatic title will be generated.

\item[\code{xlab}] A character string specifying the x-axis label. If NULL (default), the split\_by\_factor name will be used.

\item[\code{ylab}] A character string specifying the y-axis label. Default is "Percentage (

\bsl{}itemcolorsA character vector of colors for the different levels of frequency\_for\_factor. If NULL (default), ROP\_colors will be used when available, otherwise rainbow colors will be used.

\bsl{}itemlegend\_titleA character string specifying the legend title. If NULL (default), the frequency\_for\_factor name will be used.

\bsl{}itemshow\_percentagesA logical value indicating whether to display percentage labels on the bars. Default is TRUE.

\bsl{}itemtext\_sizeA numeric value specifying the size of percentage labels on the bars. Default is 3.

\bsl{}itemtext\_colorA character string specifying the color of percentage labels. Default is "black".

\bsl{}itemshow\_summary\_tableA logical value indicating whether to print the summary table. Default is TRUE.

\bsl{}itemexport\_summary\_tableA logical value indicating whether to export the summary table to the main environment. Default is FALSE.

\bsl{}itemname\_summary\_tableA character string specifying the name of the exported summary table. Default is "summary\_table\_FactorStackedBarplot".

\end{ldescription}

A ggplot object representing the stacked barplot. If export\_summary\_table is TRUE, the summary table is attached as an attribute (accessible via attr(plot, "summary\_table")).


Creates and displays stacked barplots showing the percentage distribution of one factor's levels across different levels of another factor from a romics\_object


This function generates stacked barplots where each bar represents a level of the split\_by\_factor, and the stacked segments show the percentage distribution of the frequency\_for\_factor levels within each category. The function automatically converts extracted factors to factor type using as.factor(romicsExtractFactor()). Only complete cases (non-NA values) are included in the analysis. Colors are taken from ROP\_colors by default when sufficient colors are available.


Geremy Clair

\end{Arguments}
\HeaderA{romicsFactorSummaryTable}{romicsFactorSummaryTable()}{romicsFactorSummaryTable}
%
\begin{Description}
Generates a summary table showing the distribution of samples across two factor levels.
This is the data table version of romicsFactorStackedBarplot, useful for exporting or further analysis.
\end{Description}
%
\begin{Usage}
\begin{verbatim}
romicsFactorSummaryTable(
  romics_object,
  split_by_factor,
  frequency_for_factor,
  percentage = TRUE,
  wide_format = FALSE,
  include_totals = FALSE,
  round_percentages = 1
)
\end{verbatim}
\end{Usage}
%
\begin{Arguments}
\begin{ldescription}
\item[\code{romics\_object}] A romics\_object created using romicsCreateObject().

\item[\code{split\_by\_factor}] A character vector (length=1) containing a factor of the romics\_object that will define the rows/categories. The list of factors can be identified by using the function romicsFactorNames()

\item[\code{frequency\_for\_factor}] A character vector (length=1) containing a factor of the romics\_object that will be used to calculate frequency distributions within each level of split\_by\_factor. The list of factors can be identified by using the function romicsFactorNames()

\item[\code{percentage}] Logical. If TRUE, includes percentage columns in addition to counts. Default: TRUE.

\item[\code{wide\_format}] Logical. If TRUE, returns data in wide format (split\_by\_factor levels as rows, frequency\_for\_factor levels as columns).
If FALSE, returns long format (one row per split\_by\_factor-frequency\_for\_factor combination). Default: FALSE.

\item[\code{include\_totals}] Logical. If TRUE and wide\_format = TRUE, includes a "Total" column with row totals. Default: FALSE.

\item[\code{round\_percentages}] Numeric. Number of decimal places for percentage values. Default: 1.
\end{ldescription}
\end{Arguments}
%
\begin{Value}
A data.frame containing the summary table in either long or wide format.
\end{Value}
%
\begin{Author}
Geremy Clair
\end{Author}
\HeaderA{romicsFeatureCutoff}{romicsFeatureCutoff()}{romicsFeatureCutoff}
%
\begin{Description}
Applies cutoff thresholds to specified features in a romics\_object. Values outside the specified range are set to NA and the missingdata layer is updated accordingly.
\end{Description}
%
\begin{Usage}
\begin{verbatim}
romicsFeatureCutoff(
  romics_object,
  feature_name = NULL,
  low_cutoff = NULL,
  high_cutoff = NULL,
  apply_to_all = FALSE
)
\end{verbatim}
\end{Usage}
%
\begin{Arguments}
\begin{ldescription}
\item[\code{romics\_object}] A romics\_object created using romicsCreateObject().

\item[\code{feature\_name}] Character string or vector. Name(s) of the feature(s) to apply cutoffs to. Must match row names in the data exactly.

\item[\code{low\_cutoff}] Numeric. Lower threshold value. Values below this will be set to NA. If NULL, uses the minimum value of the feature. Default: NULL

\item[\code{high\_cutoff}] Numeric. Upper threshold value. Values above this will be set to NA. If NULL, uses the maximum value of the feature. Default: NULL

\item[\code{apply\_to\_all}] Logical. If TRUE, applies the same cutoffs to all features in the dataset. Overrides feature\_name. Default: FALSE
\end{ldescription}
\end{Arguments}
%
\begin{Details}
This function is useful for removing outliers or applying quality control thresholds to specific features.
When cutoffs are applied, the original data structure is preserved but values outside the range become missing values.
The missingdata layer is automatically updated to reflect the new missing values.
If low\_cutoff or high\_cutoff are NULL, they default to the min/max values of each feature respectively,
effectively meaning no cutoff is applied on that end.
\end{Details}
%
\begin{Value}
Returns a romics\_object with cutoffs applied and missingdata layer updated.
\end{Value}
%
\begin{Author}
Geremy Clair
\end{Author}
\HeaderA{romicsFeatureNames}{romicsFeatureNames()}{romicsFeatureNames}
%
\begin{Description}
Indicates the list of feature names from the romics\_object (e.g.,rownames )
\end{Description}
%
\begin{Usage}
\begin{verbatim}
romicsFeatureNames(romics_object)
\end{verbatim}
\end{Usage}
%
\begin{Arguments}
\begin{ldescription}
\item[\code{romics\_object}] A romics\_object created using createRomicsObject()
\end{ldescription}
\end{Arguments}
%
\begin{Details}
This function allows to quickly get a vector containing all the factor names present in a romics\_object
\end{Details}
%
\begin{Value}
A character vector containing the list of features contained in an romics\_object
\end{Value}
%
\begin{Author}
Geremy Clair
\end{Author}
\HeaderA{romicsFeatureSubset}{romicsFeatureSubset()}{romicsFeatureSubset}
%
\begin{Description}
This will create an romics\_object subseted for a list of features only
\end{Description}
%
\begin{Usage}
\begin{verbatim}
romicsFeatureSubset(
  romics_object,
  feature_list = c("feature1", "feature2", "etc")
)
\end{verbatim}
\end{Usage}
%
\begin{Arguments}
\begin{ldescription}
\item[\code{romics\_object}] has to be an romics\_object created using romicsCreateObject()

\item[\code{feature\_list}] has to be a character vector containing features to be extracted
\end{ldescription}
\end{Arguments}
%
\begin{Details}
This function will create a new romics object in which only the features contained in the feature\_list are contained.

The function supports exact matches only, if some features are missing, the function will still work but only will return an object with the existing features.
\end{Details}
%
\begin{Value}
This function returns a subseted romics\_object
\end{Value}
%
\begin{Author}
Geremy Clair
\end{Author}
\HeaderA{romicsFilterFeature}{romicsFilterFeature()}{romicsFilterFeature}
%
\begin{Description}
Filters Romics\_objects based on their feature names or statistics note that ALL the filters will be taken in consideration simultaneously
\end{Description}
%
\begin{Usage}
\begin{verbatim}
romicsFilterFeature(
  romics_object,
  ANOVA_filter = "none",
  p = 0.05,
  feature_names = "none",
  statCol = "none",
  statCol_filter = "<=0.05",
  statCol2 = "none",
  statCol2_filter = "<=0.05",
  mode = "keep"
)
\end{verbatim}
\end{Usage}
%
\begin{Arguments}
\begin{ldescription}
\item[\code{romics\_object}] A romics\_object created using romicsCreateObject()

\item[\code{ANOVA\_filter}] Either 'none', 'p' or 'padj'. Indicates if an the ANOVA filter has to be used to filter features (only the features below the filter will be retained)

\item[\code{p}] Numerical of length 1 indicating the value of the ANOVA\_filter cutoff (anything below this value will be conserved).

\item[\code{feature\_names}] A character vector that enable to filter based on the names of the features, if feature\_names set to 'none' the filter won't be applied.

\item[\code{statCol}] A column contained in the statistical layer of the romics\_object, the list of columns can be obtained by using the function romicsCalculatedStats().

\item[\code{statCol\_filter}] Character to indicate how this column should be filtered (e.g. '<=0.05','>0.05','==1', '==TRUE', '>2')

\item[\code{statCol2}] A column contained in the statistical layer of the romics\_object, the list of columns can be obtained by using the function romicsCalculatedStats().

\item[\code{statCol2\_filter}] Character to indicate how this column should be filtered (e.g. '<=0.05','>0.05','==1', '==TRUE', '>2')

\item[\code{mode}] Either 'keep' or 'drop' to indicate if the features should be kept or dropped based on the filters.
\end{ldescription}
\end{Arguments}
%
\begin{Details}
This function applies multiple filters simultaneously to retain or remove features based on:
- ANOVA p-values or adjusted p-values
- Feature name patterns
- Statistical column criteria
All filters are combined using AND logic (all conditions must be met).
\end{Details}
%
\begin{Value}
Returns a romics\_object with filtered features
\end{Value}
%
\begin{Author}
Geremy Clair
\end{Author}
\HeaderA{romicsFilterMissing}{romicsFilterMissing()}{romicsFilterMissing}
%
\begin{Description}
Filters out the variables of the romics\_object below the choosen percentage of completeness. The percentage of completeness can either be global or by factor of a given factor (in this later case if the percentage of completeness is achieved for at least one level of the factor it the variable will be kept).
\end{Description}
%
\begin{Usage}
\begin{verbatim}
romicsFilterMissing(
  romics_object,
  percentage_completeness = 50,
  factor = "main",
  all_groups = FALSE
)
\end{verbatim}
\end{Usage}
%
\begin{Arguments}
\begin{ldescription}
\item[\code{romics\_object}] A romics\_object created using romicsCreateObject()

\item[\code{percentage\_completeness}] Numerical vector indicating the minimum percentage of data to be considered

\item[\code{factor}] has to be either "main", "none" or a factor of an romics\_object created using romicsCreateObject() the list of factors can be obtained by running the function romicsFactorNames() on the romics\_object

\item[\code{all\_groups}] if this parameter is TRUE the completeness requirement is for each and every group, if not, the completeness requirement is for at least one group.
\end{ldescription}
\end{Arguments}
%
\begin{Details}
This function will use the completeness of the features in the overall samples (when none is used as factor), or of a given level of a specific defined factor (in this case the factor has to be set to either "main" or to the given factor of filtering). By default factor is the main factor of the object, the percentage\_completeness is set at 50
\end{Details}
%
\begin{Value}
The function will return a filtered romics\_object with the rows of the data and missing data object removed when appropriate.
\end{Value}
%
\begin{Author}
Geremy Clair, Nicholas Day
\end{Author}
\HeaderA{romicsFilterMissingSamples}{romicsFilterMissingSamples()}{romicsFilterMissingSamples}
%
\begin{Description}
Filters out the samples of the romics\_object\$data layer based on their level of missingness below a chosen percentage of completeness for the different features
\end{Description}
%
\begin{Usage}
\begin{verbatim}
romicsFilterMissingSamples(romics_object, percentage_completeness = 50)
\end{verbatim}
\end{Usage}
%
\begin{Arguments}
\begin{ldescription}
\item[\code{romics\_object}] A romics\_object created using romicsCreateObject().

\item[\code{percentage\_completeness}] Numerical value indicating the minimum percentage of completeness required for a sample to be kept.
\end{ldescription}
\end{Arguments}
%
\begin{Details}
This function removes columns (samples) from the romics\_object where the percentage of missing data exceeds the chosen threshold. By default, percentage\_completeness is set at 50
\end{Details}
%
\begin{Value}
The function will return a filtered romics\_object with samples (columns) removed based on their level of missingness.
\end{Value}
%
\begin{Author}
Geremy Clair, Nicholas Day, Modified by AI Assistant
\end{Author}
\HeaderA{romicsFoldChangeBarplot}{romicsFoldChangeBarplot()}{romicsFoldChangeBarplot}
%
\begin{Description}
Creates fold change barplots with significance annotations from a romics object
\end{Description}
%
\begin{Usage}
\begin{verbatim}
romicsFoldChangeBarplot(
  romics_object,
  feature_list = NULL,
  comparison = "all",
  test_type = "auto",
  use_adjusted_p = FALSE,
  p_threshold = 0.05,
  make_horizontal = FALSE,
  order_by = "provided",
  significant_only = FALSE,
  bar_colors = NULL,
  show_values = FALSE,
  title = NULL,
  fc_limits = NULL,
  use_theme_ROP = TRUE
)
\end{verbatim}
\end{Usage}
%
\begin{Arguments}
\begin{ldescription}
\item[\code{romics\_object}] A romics object containing \$data, \$statistics, and \$metadata layers

\item[\code{feature\_list}] Character vector of features to include in barplot. If NULL, uses all features (limited to 100 max)

\item[\code{comparison}] Character vector of comparisons to include. Default "all" includes all available comparisons. Format: c("groupA\_vs\_groupB", "groupB\_vs\_groupC")

\item[\code{test\_type}] Which test to use for fold changes. Options: "ttest", "wilcoxon", "auto" (default: "auto" - uses first available)

\item[\code{use\_adjusted\_p}] TRUE or FALSE, if TRUE uses adjusted p-values, if FALSE uses raw p-values (default: TRUE)

\item[\code{p\_threshold}] Numeric threshold for p-values (default 0.05)

\item[\code{make\_horizontal}] TRUE or FALSE, if TRUE creates horizontal barplot (features on y-axis), if FALSE creates vertical (features on x-axis) (default: FALSE)

\item[\code{order\_by}] How to order features. Options: "provided" (default), "alphabetical", "pvalue", "fold\_change"

\item[\code{significant\_only}] TRUE or FALSE, if TRUE only shows features significant in at least one comparison (default: FALSE)

\item[\code{bar\_colors}] Character vector of colors for different comparisons. If NULL, uses default palette

\item[\code{show\_values}] TRUE or FALSE, if TRUE shows fold change values on bars (default: FALSE)

\item[\code{title}] Character string for plot title. If NULL, auto-generates title

\item[\code{fc\_limits}] Numeric vector of length 2 specifying y-axis limits for fold changes c(min, max). If NULL, uses automatic limits

\item[\code{use\_theme\_ROP}] TRUE or FALSE, if TRUE applies theme\_ROP() to the plot, if FALSE uses theme\_classic() (default: TRUE)
\end{ldescription}
\end{Arguments}
%
\begin{Details}
Creates barplots showing fold changes with significance stars. Uses T-test or Wilcoxon test results from statistics layer.
\end{Details}
%
\begin{Value}
ggplot2 object
\end{Value}
%
\begin{Author}
Geremy Clair
\end{Author}
\HeaderA{romicsFreqMissingFeatures}{romicsFreqMissingFeatures()}{romicsFreqMissingFeatures}
%
\begin{Description}
Plots the frequency of the feature missingness for all features contained in an romics\_object.
\end{Description}
%
\begin{Usage}
\begin{verbatim}
romicsFreqMissingFeatures(romics_object = romics_object, binwidth = 1)
\end{verbatim}
\end{Usage}
%
\begin{Arguments}
\begin{ldescription}
\item[\code{romics\_object}] has to be an romics\_object created using romicsCreateObject()

\item[\code{binwidth}] has to be a numeric value indicating the width of the desired frequency bins, 1 will be used by default.
\end{ldescription}
\end{Arguments}
%
\begin{Details}
This function does not alter the romics\_object, it generates a table containing the percentage of missingness for each feature of the romics\_object.
\end{Details}
%
\begin{Value}
This function will return a ggplot2 geom\_bar plot. it can then be further visually adjusted using the ggplot2 commands.
\end{Value}
%
\begin{Author}
Geremy Clair
\end{Author}
\HeaderA{romicsGLMheatmap}{romicsGLMheatmap()}{romicsGLMheatmap}
%
\begin{Description}
Generate a heatmap displaying the presence/absence pattern of features across samples with GLM binomial test significance annotations
\end{Description}
%
\begin{Usage}
\begin{verbatim}
romicsGLMheatmap(
  romics_object,
  filter_p = TRUE,
  p_type = c("p", "padj"),
  p = 0.05,
  mode = c("vs", "enrichment"),
  hclust_row = TRUE,
  row_annotation = TRUE,
  column_annotation = TRUE,
  show_rownames = FALSE,
  factor = NULL
)
\end{verbatim}
\end{Usage}
%
\begin{Arguments}
\begin{ldescription}
\item[\code{romics\_object}] A romics\_object created using romicsCreateObject() containing calculated GLM binomial test statistics.

\item[\code{filter\_p}] logical (TRUE or FALSE) to indicate if features should be filtered by GLM test p-values.

\item[\code{p\_type}] 'p' or 'padj' to indicate the type of p-value to use for filtering and significance determination.

\item[\code{p}] numeric value indicating the p-value threshold for significance (between 0 and 1).

\item[\code{mode}] 'vs' to show up/down regulation patterns or 'enrichment' to show only enrichment (positive directionality).

\item[\code{hclust\_row}] logical (TRUE or FALSE) to indicate whether to perform hierarchical clustering on rows.

\item[\code{row\_annotation}] logical (TRUE or FALSE) to indicate whether to display row annotations showing significance status for each comparison (red = significant, grey = not significant).

\item[\code{column\_annotation}] logical (TRUE or FALSE) to indicate whether to display column annotation bar showing the factor level for each sample.

\item[\code{show\_rownames}] logical (TRUE or FALSE) to indicate whether to display feature names on the left side of the heatmap.

\item[\code{factor}] character or NULL to specify which factor to use for column annotations. If NULL, uses the main factor from the romics\_object.
\end{ldescription}
\end{Arguments}
%
\begin{Details}
Generate a ComplexHeatmap showing presence/absence of features with GLM binomial test results. Row annotations show significance status using a unified color scheme (red/grey), and column annotations show factor levels using ROP\_colors palette.
\end{Details}
%
\begin{Value}
A ComplexHeatmap object
\end{Value}
%
\begin{Author}
Geremy Clair
\end{Author}
\HeaderA{romicsHclust}{romicsHclust()}{romicsHclust}
%
\begin{Description}
Plots the hierarchical clustering of the samples contaned in the romics\_object calculated using the functions dist() and hclust() (see documentation of those R functions for more details). The colors used for the plotting will correspond to the main\_factor of the romics\_object.
\end{Description}
%
\begin{Usage}
\begin{verbatim}
romicsHclust(
  romics_object,
  method_dist = "euclidean",
  method_hclust = "ward.D"
)
\end{verbatim}
\end{Usage}
%
\begin{Arguments}
\begin{ldescription}
\item[\code{romics\_object}] has to be a log transformed romics\_object created using romicsCreateObject() and transformed using the function log2transform() or log10transform()

\item[\code{method\_dist}] the distance measure to be used. This must be one of "euclidean", "maximum", "manhattan", "canberra", "binary" or "minkowski". Any unambiguous substring can be given.

\item[\code{method\_hclust}] the agglomeration method to be used. This should be (an unambiguous abbreviation of) one of "ward.D", "ward.D2", "single", "complete", "average" (= UPGMA), "mcquitty" (= WPGMA), "median" (= WPGMC) or "centroid" (= UPGMC).
\end{ldescription}
\end{Arguments}
%
\begin{Details}
This function uses the dist() and hclust() functions to calculate the hierachical clustering and then plots the hclust with colors based on the current main\_factor of the romics\_object.
\end{Details}
%
\begin{Value}
a hierarchical clustering tree plot with its branches colored by factor.
\end{Value}
%
\begin{Author}
Geremy Clair
\end{Author}
\HeaderA{romicsHeatmap}{romicsHeatmap()}{romicsHeatmap}
%
\begin{Description}
Plots a scaled heatmap of the data layer from a romics\_object using heatmap.2 from the gplots package. Data can be filtered based on the statistics layer. Optionally records cluster assignments to the statistics layer. Can also collapse samples by factor level before plotting.
\end{Description}
%
\begin{Usage}
\begin{verbatim}
romicsHeatmap(
  romics_object,
  color_palette = NULL,
  color_boundaries = c(-2, 2),
  sample_hclust = TRUE,
  sample_hclust_method_dist = "euclidean",
  sample_hclust_method_hclust = "ward.D",
  feature_hclust = TRUE,
  feature_hclust_method_dist = "euclidean",
  feature_hclust_method_hclust = "ward.D",
  feature_hclust_number = 1,
  show_feature_clusters = FALSE,
  record_feature_clusters = FALSE,
  feature_cluster_colors = NULL,
  show_cluster_legend = TRUE,
  cluster_legend_position = "topright",
  cluster_legend_cex = 0.8,
  collapse_by_factor = "none",
  collapse_method = "mean",
  statCol = "none",
  statCol_filter = "<=0.05",
  statCol2 = "none",
  statCol2_filter = "<=0.05",
  feature_list = "all",
  notecol = "black",
  density.info = "none",
  trace = "none",
  labRow = FALSE,
  cexCol = 1,
  margins = c(15, 5),
  key.title = "Scaled Heatmap",
  key.xlab = "Z-scores",
  ...
)
\end{verbatim}
\end{Usage}
%
\begin{Arguments}
\begin{ldescription}
\item[\code{romics\_object}] A romics\_object created using romicsCreateObject().

\item[\code{color\_palette}] Character vector of colors. Default: viridis(20)

\item[\code{color\_boundaries}] Numerical vector of length 2 for min and max of color scale. Default: c(-2,2)

\item[\code{sample\_hclust}] Boolean. Whether to cluster samples. Default: TRUE

\item[\code{sample\_hclust\_method\_dist}] Distance method for samples. Default: "euclidean"

\item[\code{sample\_hclust\_method\_hclust}] Agglomeration method for samples. Default: "ward.D"

\item[\code{feature\_hclust}] Boolean. Whether to cluster features. Default: TRUE

\item[\code{feature\_hclust\_method\_dist}] Distance method for features. Default: "euclidean"

\item[\code{feature\_hclust\_method\_hclust}] Agglomeration method for features. Default: "ward.D"

\item[\code{feature\_hclust\_number}] Number of clusters for feature coloring. Default: 1

\item[\code{show\_feature\_clusters}] Boolean. Whether to show cluster assignments as a sidebar. Default: FALSE. Requires feature\_hclust\_number > 1.

\item[\code{record\_feature\_clusters}] Boolean. Whether to save cluster assignments to the statistics layer as 'hclust\_clusters' column. Default: FALSE. Requires feature\_hclust\_number > 1.

\item[\code{feature\_cluster\_colors}] Character vector of colors for cluster sidebar. If NULL, uses a default color palette. Length should match number of unique clusters.

\item[\code{show\_cluster\_legend}] Boolean. Whether to show a legend for cluster colors when show\_feature\_clusters = TRUE. Default: TRUE.

\item[\code{cluster\_legend\_position}] Position for cluster legend: "topright", "topleft", "bottomright", "bottomleft", "top", "bottom", "left", "right". Default: "topright"

\item[\code{cluster\_legend\_cex}] Size of legend text. Default: 0.8

\item[\code{collapse\_by\_factor}] Character. Name of the factor to collapse samples by (average within each factor level). Default: "none" (no collapsing). If specified, samples are grouped by factor level and averaged before plotting.

\item[\code{collapse\_method}] Character. Method for collapsing: "mean" or "median". Default: "mean"

\item[\code{statCol}] Statistical column for filtering. Default: "none"

\item[\code{statCol\_filter}] Filter expression for statCol. Default: "<=0.05"

\item[\code{statCol2}] Second statistical column for filtering. Default: "none"

\item[\code{statCol2\_filter}] Filter expression for statCol2. Default: "<=0.05"

\item[\code{feature\_list}] Character vector of feature names to include in the heatmap. Default: "all"

\item[\code{notecol}] Color for cell note text. Default: "black"

\item[\code{density.info}] Density plot on color key: 'histogram', 'density', or 'none'. Default: "none"

\item[\code{trace}] Trace lines: "column", "row", "both", or "none". Default: "none"

\item[\code{labRow}] Show row labels. Can be TRUE, FALSE, or character vector. Default: FALSE

\item[\code{cexCol}] Size of column labels. Default: 1

\item[\code{margins}] Margins for labels c(bottom, right). Default: c(15, 5)

\item[\code{key.title}] Title for color key. Default: "Scaled Heatmap"

\item[\code{key.xlab}] X-axis label for color key. Default: "Z-scores"

\item[\code{...}] Additional parameters for heatmap.2
\end{ldescription}
\end{Arguments}
%
\begin{Details}
When record\_feature\_clusters = TRUE, the function will save cluster assignments to romics\_object\$statistics\$hclust\_clusters
and return the updated romics\_object. The heatmap is still displayed.

When record\_feature\_clusters = FALSE (default), the function only displays the heatmap and returns NULL invisibly.

Note: Cluster assignments are only recorded for features that pass the filtering criteria and appear in the heatmap.
Features not in the heatmap will have NA in the hclust\_clusters column.

When collapse\_by\_factor is specified, samples are grouped by the specified factor level and averaged
(or median-collapsed) before plotting. This reduces the number of columns in the heatmap and allows
for easier visualization of group-level patterns. The original romics\_object is not modified.
\end{Details}
%
\begin{Value}
If record\_feature\_clusters = TRUE, returns the updated romics\_object. Otherwise, returns NULL invisibly. The heatmap is displayed in both cases.
\end{Value}
%
\begin{Author}
Geremy Clair and Harsh Bhotika
\end{Author}
\HeaderA{romicsIDsNames}{romicsIDsNames()}{romicsIDsNames}
%
\begin{Description}
Indicates the list of ID names from the romics\_object
\end{Description}
%
\begin{Usage}
\begin{verbatim}
romicsIDsNames(romics_object)
\end{verbatim}
\end{Usage}
%
\begin{Arguments}
\begin{ldescription}
\item[\code{romics\_object}] A romics\_object created using createRomicsObject()
\end{ldescription}
\end{Arguments}
%
\begin{Details}
This function allows to quickly get a vector containing all the ID names present in a romics\_object
\end{Details}
%
\begin{Value}
A character vector containing the list of IDs contained in an romics\_object
\end{Value}
%
\begin{Author}
Geremy Clair
\end{Author}
\HeaderA{romicsKmeansSamples}{romicsKmeansSamples()}{romicsKmeansSamples}
%
\begin{Description}
Generate K-means clustering for the samples in a romics\_object. The clustering can be based on
the full data matrix or dimensionally reduced data from PCA, UMAP, or t-SNE embeddings.
\end{Description}
%
\begin{Usage}
\begin{verbatim}
romicsKmeansSamples(
  romics_object,
  centers = 3,
  cluster_using = c("data", "pca", "umap", "tsne"),
  factor_name = "K_clust",
  iter.max = 100,
  algorithm = c("Lloyd", "Hartigan-Wong", "Forgy", "MacQueen"),
  lock.seed = TRUE,
  seed = 42,
  ...
)
\end{verbatim}
\end{Usage}
%
\begin{Arguments}
\begin{ldescription}
\item[\code{romics\_object}] A romics\_object created using romicsCreateObject()

\item[\code{centers}] Either the number of clusters, or a set of initial (distinct) cluster centers.
If a number, a random set of (distinct) rows in the data is chosen as the initial centers.

\item[\code{cluster\_using}] Character string specifying which data to use for clustering:
'data' (full data matrix), 'pca' (PCA embeddings), 'umap' (UMAP embeddings), or 'tsne' (t-SNE embeddings). Default: 'data'

\item[\code{factor\_name}] Character string specifying the name for the new factor to be added to metadata.
Default: 'K\_clust'

\item[\code{iter.max}] The maximum number of iterations allowed. Default: 100

\item[\code{algorithm}] Character string specifying the K-means algorithm to use:
'Hartigan-Wong', 'Lloyd', 'Forgy', or 'MacQueen'. Default: 'Lloyd'

\item[\code{lock.seed}] Logical. If TRUE, uses a fixed seed for reproducible clustering. Default: TRUE

\item[\code{seed}] Numeric value for the random seed when lock.seed=TRUE. Default: 42

\item[\code{...}] Additional parameters passed to the kmeans() function.
\end{ldescription}
\end{Arguments}
%
\begin{Details}
This function performs K-means clustering on samples and adds the cluster assignments
as a new factor in the romics\_object's metadata. Using dimensionally-reduced data
(PCA, UMAP, or t-SNE) can improve computation time but may reduce the accuracy of the clusters.
The function ensures that "colors\_romics" factor remains the last row in the metadata.
Any NA values in cluster assignments will be replaced with "undefined cluster".
\end{Details}
%
\begin{Value}
Returns the romics\_object with a new factor added to metadata containing the cluster assignments.
Clusters are named as "factor\_name\_001", "factor\_name\_002", etc.
The "colors\_romics" factor is guaranteed to be the last row in metadata.
\end{Value}
%
\begin{Author}
Geremy Clair
\end{Author}
\HeaderA{romicsKmeansSamplesEval}{romicsKmeansSamplesEval()}{romicsKmeansSamplesEval}
%
\begin{Description}
Evaluates the optimal number of clusters for K-means using the elbow method.
Plots the within-cluster sum of squares (WCSS) for different numbers of clusters.
\end{Description}
%
\begin{Usage}
\begin{verbatim}
romicsKmeansSamplesEval(
  romics_object,
  iter.max = 100,
  algorithm = c("Lloyd", "Hartigan-Wong", "Forgy", "MacQueen"),
  cluster_using = c("data", "pca", "umap", "tsne"),
  max_clust = 50,
  ...
)
\end{verbatim}
\end{Usage}
%
\begin{Arguments}
\begin{ldescription}
\item[\code{romics\_object}] A romics\_object created using romicsCreateObject()

\item[\code{iter.max}] The maximum number of iterations allowed. Default: 100

\item[\code{algorithm}] Character string specifying the K-means algorithm to use:
'Hartigan-Wong', 'Lloyd', 'Forgy', or 'MacQueen'. Default: 'Lloyd'

\item[\code{cluster\_using}] Character string specifying which data to use for clustering:
'data' (full data matrix), 'pca' (PCA embeddings), 'umap' (UMAP embeddings), or 'tsne' (t-SNE embeddings). Default: 'data'

\item[\code{max\_clust}] Maximum number of clusters to evaluate. Default: 50

\item[\code{...}] Additional parameters passed to kmeans().
\end{ldescription}
\end{Arguments}
%
\begin{Value}
Returns a plot showing the WCSS for clusters from 1 to max\_clust.
\end{Value}
%
\begin{Author}
Geremy Clair
\end{Author}
\HeaderA{romicsLeidenSamples}{romicsLeidenSamples()}{romicsLeidenSamples}
%
\begin{Description}
OPTIMIZED VERSION: Generate Leiden community clustering for samples using fast k-nearest neighbor graph construction.
Uses vectorized operations and avoids redundant distance calculations for significantly improved performance on large datasets.
The clustering can be based on the full data matrix or dimensionally reduced data from PCA, UMAP, or t-SNE embeddings.
\end{Description}
%
\begin{Usage}
\begin{verbatim}
romicsLeidenSamples(
  romics_object,
  cluster_using = c("data", "pca", "umap", "tsne"),
  k = 20,
  resolution = 1,
  target_clusters = NULL,
  max_iterations = 20,
  tolerance = 1,
  factor_name = "Leiden_clust",
  weights = NULL,
  lock.seed = TRUE,
  seed = 42,
  use_approximate = TRUE,
  objective_function = c("modularity", "CPM", "significance"),
  n_iterations = -1,
  ...
)
\end{verbatim}
\end{Usage}
%
\begin{Arguments}
\begin{ldescription}
\item[\code{romics\_object}] A romics\_object created using romicsCreateObject().

\item[\code{cluster\_using}] Character string specifying which data to use for clustering. Options: 'data', 'pca', 'umap', 'tsne'.
Default: 'data'

\item[\code{k}] Number of nearest neighbors to use in graph construction. Default: 20
For datasets >50k samples, consider reducing this value (e.g., k=10) to improve memory efficiency.

\item[\code{resolution}] Resolution parameter that determines the granularity of the clustering. Default: 1.0

\item[\code{target\_clusters}] Optional integer specifying the desired number of clusters. Default: NULL

\item[\code{max\_iterations}] Maximum number of iterations when trying to reach target\_clusters. Default: 20

\item[\code{tolerance}] Tolerance for target\_clusters - clustering will stop if within this range. Default: 1

\item[\code{factor\_name}] Character string specifying the name for the new factor to be added to metadata.
Default: 'Leiden\_clust'

\item[\code{weights}] Optional edge weights for the graph. If NULL, distance-based weighting is used automatically.

\item[\code{lock.seed}] Logical. If TRUE, uses a fixed seed for reproducible clustering. Default: TRUE

\item[\code{seed}] Numeric value for the random seed when lock.seed=TRUE. Default: 42

\item[\code{use\_approximate}] Logical. If TRUE and FNN available, uses fast approximate KNN. Default: TRUE

\item[\code{objective\_function}] Character string: 'modularity', 'CPM', or 'significance'. Default: 'modularity'

\item[\code{n\_iterations}] Maximum number of Leiden iterations. Default: -1 (run until convergence)

\item[\code{...}] Additional parameters passed to igraph::cluster\_leiden().
\end{ldescription}
\end{Arguments}
%
\begin{Details}
OPTIMIZED IMPLEMENTATION:
- Pre-allocates edge matrices instead of vector concatenation (10-50x faster)
- Reuses KNN distances directly instead of recalculating (5-10x faster)
- Vectorizes weight calculations (3-5x faster)
- Overall expected speedup: 20-100x on large datasets

Uses the same interface as romicsLeidenSamples\_old but with significantly improved performance.
\end{Details}
%
\begin{Value}
Returns the romics\_object with cluster assignments added to metadata.
\end{Value}
%
\begin{Author}
Geremy Clair
\end{Author}
\HeaderA{romicsLeidenSamples\_old}{romicsLeidenSamples()}{romicsLeidenSamples.Rul.old}
%
\begin{Description}
Generate Leiden community clustering for samples in a romics\_object using memory-efficient algorithms.
The clustering can be based on the full data matrix or dimensionally reduced data from PCA, UMAP, or t-SNE embeddings.
For large datasets (>10k samples), this function uses optimized nearest neighbor search to reduce memory usage.
The Leiden algorithm is an improvement over Louvain that can escape local optima.

DEPRECATED: Original (slow) version of Leiden clustering. Use romicsLeidenSamples() instead for significantly better performance.
This function is kept for backward compatibility only.
\end{Description}
%
\begin{Usage}
\begin{verbatim}
romicsLeidenSamples_old(
  romics_object,
  cluster_using = c("data", "pca", "umap", "tsne"),
  k = 20,
  resolution = 1,
  target_clusters = NULL,
  max_iterations = 20,
  tolerance = 1,
  factor_name = "Leiden_clust",
  weights = NULL,
  lock.seed = TRUE,
  seed = 42,
  use_approximate = TRUE,
  objective_function = c("modularity", "CPM", "significance"),
  n_iterations = -1,
  ...
)
\end{verbatim}
\end{Usage}
%
\begin{Arguments}
\begin{ldescription}
\item[\code{romics\_object}] A romics\_object created using romicsCreateObject().

\item[\code{cluster\_using}] Character string specifying which data to use for clustering.

\item[\code{k}] Number of nearest neighbors to use in graph construction.

\item[\code{resolution}] Resolution parameter for the clustering.

\item[\code{target\_clusters}] Optional integer specifying the desired number of clusters.

\item[\code{max\_iterations}] Maximum number of iterations when trying to reach target\_clusters.

\item[\code{tolerance}] Tolerance for target\_clusters.

\item[\code{factor\_name}] Character string specifying the name for the new factor.

\item[\code{weights}] Optional edge weights for the graph.

\item[\code{lock.seed}] Logical. If TRUE, uses a fixed seed for reproducible clustering.

\item[\code{seed}] Numeric value for the random seed.

\item[\code{use\_approximate}] Logical. If TRUE and FNN available, uses fast approximate KNN.

\item[\code{objective\_function}] Character string: 'modularity', 'CPM', or 'significance'.

\item[\code{n\_iterations}] Maximum number of Leiden iterations.

\item[\code{...}] Additional parameters passed to igraph::cluster\_leiden().
\end{ldescription}
\end{Arguments}
%
\begin{Details}
This function performs Leiden community detection on samples and adds the cluster assignments
as a new factor in the romics\_object's metadata. The function constructs a k-nearest neighbor
graph from the data using memory-efficient algorithms, then applies the Leiden algorithm to identify communities.

The Leiden algorithm is an improvement over the Louvain algorithm that can escape local optima
and provides better quality partitions. It guarantees well-connected communities.

When target\_clusters is specified, the function will iteratively adjust the resolution parameter
to try to achieve approximately that number of clusters. Note that exact control is not guaranteed
due to the nature of the Leiden algorithm.

For large datasets, the function automatically uses optimized approaches:
- FNN package for fast approximate nearest neighbor search (if available)
- Edge list representation instead of full adjacency matrices
- Chunked processing to control memory usage
- Periodic garbage collection to free memory

Memory recommendations for large datasets:
- Use dimensionally reduced data (cluster\_using = "pca" or "umap")
- Reduce k parameter (e.g., k=10 for >100k samples)
- Install FNN package: install.packages("FNN")
- Ensure sufficient system memory (recommend >16GB for >100k samples)
\end{Details}
%
\begin{Value}
Returns the romics\_object with a new factor added to metadata containing the cluster assignments.
Clusters are named as "factor\_name\_001", "factor\_name\_002", etc.
A message will indicate the number of clusters found and confirm the factor was added.

Returns the romics\_object with cluster assignments added to metadata.
\end{Value}
%
\begin{Author}
Geremy Clair
\end{Author}
%
\begin{SeeAlso}
\code{\LinkA{romicsLouvainSamples}{romicsLouvainSamples}} for Louvain clustering,
\code{\LinkA{romicsPCA}{romicsPCA}}, \code{\LinkA{romicsUMAP}{romicsUMAP}} for dimensionality reduction
romicsLeidenSamples\_old()
\end{SeeAlso}
%
\begin{Examples}
\begin{ExampleCode}
## Not run: 
# Basic usage
romics_obj <- romicsLeidenSamples(romics_obj)

# Target specific number of clusters
romics_obj <- romicsLeidenSamples(romics_obj, target_clusters = 5)

# For large datasets, use PCA embeddings with reduced k
romics_obj <- romicsLeidenSamples(romics_obj,
                                  cluster_using = "pca",
                                  k = 10,
                                  target_clusters = 8)

# Custom factor name and resolution with CPM objective
romics_obj <- romicsLeidenSamples(romics_obj,
                                  cluster_using = "umap",
                                  resolution = 1.5,
                                  factor_name = "UMAP_Leiden",
                                  objective_function = "CPM")

## End(Not run)
\end{ExampleCode}
\end{Examples}
\HeaderA{romicsLMM}{romicsLMM()}{romicsLMM}
%
\begin{Description}
Performs linear mixed model analysis for each variable using specified fixed and random effects. The results are added as new columns in the statistics layer.
\end{Description}
%
\begin{Usage}
\begin{verbatim}
romicsLMM(
  romics_object,
  fixed_effects,
  random_effects,
  cluster_factor = "none",
  method = "satterthwaite",
  padj = TRUE,
  padj_method = "BH",
  percentage_completeness = 0,
  REML = NULL,
  handle_singular = "simplify",
  suppress_warnings = TRUE,
  pairwise_comparisons = TRUE
)
\end{verbatim}
\end{Usage}
%
\begin{Arguments}
\begin{ldescription}
\item[\code{romics\_object}] An romics\_object created with the function romicsCreateObject().

\item[\code{fixed\_effects}] A character string specifying the fixed effects formula (e.g., "condition" or "condition + treatment").

\item[\code{random\_effects}] A character string specifying the random effects formula (e.g., "(1|donor)" or "(1|donor) + (1|batch)").

\item[\code{cluster\_factor}] A character string indicating a clustering factor to subset data by (e.g., "Leiden\_clust\_names"). If "none", analysis is performed on the full dataset. Default: "none".

\item[\code{method}] A character string specifying the method for p-value calculation: "satterthwaite" (using lmerTest) or "likelihood\_ratio" (comparing nested models). Default: "satterthwaite".

\item[\code{padj}] A logical variable indicating whether to perform p-value adjustment. Default: TRUE.

\item[\code{padj\_method}] Correction method. Must be in c("holm", "hochberg", "hommel", "bonferroni", "BH", "BY", "fdr"). Default: "BH".

\item[\code{percentage\_completeness}] A numerical value between 0 and 100 indicating the minimum completeness required to calculate the LMM (if not met, p-value will be NA). Default: 0.

\item[\code{REML}] A logical indicating whether to use REML estimation. Default: TRUE for satterthwaite method, FALSE for likelihood\_ratio method.

\item[\code{handle\_singular}] A character string specifying how to handle singular fits: "keep" (keep results with warning), "simplify" (try simpler random effects), "drop" (return NA). Default: "simplify".

\item[\code{suppress\_warnings}] A logical indicating whether to suppress convergence warnings. Default: TRUE.

\item[\code{pairwise\_comparisons}] A logical indicating whether to perform all pairwise comparisons for factors with 3+ levels. Default: TRUE.
\end{ldescription}
\end{Arguments}
%
\begin{Details}
This function fits linear mixed models to test fixed effects while accounting for random effects structure. When pairwise\_comparisons is TRUE and a factor has 3 or more levels, all pairwise comparisons are computed. When models fail to converge or produce singular fits, the function can automatically simplify the random effects structure or fall back to fixed effects only models. When cluster\_factor is specified, separate models are fitted for each cluster level.
\end{Details}
%
\begin{Value}
An romics\_object with the statistics layer containing the newly generated LMM results including effect sizes (log fold changes or differences), p-values, and adjusted p-values.
\end{Value}
%
\begin{Author}
Geremy Clair
\end{Author}
\HeaderA{romicsLogCheck}{romicsLogCheck()}{romicsLogCheck}
%
\begin{Description}
Identifies if the romics\_object is log\_transformed or not
\end{Description}
%
\begin{Usage}
\begin{verbatim}
romicsLogCheck(romics_object)
\end{verbatim}
\end{Usage}
%
\begin{Arguments}
\begin{ldescription}
\item[\code{romics\_object}] A romics\_object created using createRomicsObject()
\end{ldescription}
\end{Arguments}
%
\begin{Value}
This function will return TRUE or FALSE indicating if the object was or not log transformed using the function log2transform
\end{Value}
%
\begin{Author}
Geremy Clair
\end{Author}
\HeaderA{romicsLouvainSamples}{romicsLouvainSamples()}{romicsLouvainSamples}
%
\begin{Description}
OPTIMIZED VERSION: Generate Louvain community clustering for samples using fast k-nearest neighbor graph construction.
Uses vectorized operations and avoids redundant distance calculations for significantly improved performance on large datasets.
The clustering can be based on the full data matrix or dimensionally reduced data from PCA, UMAP, or t-SNE embeddings.
\end{Description}
%
\begin{Usage}
\begin{verbatim}
romicsLouvainSamples(
  romics_object,
  cluster_using = c("data", "pca", "umap", "tsne"),
  k = 20,
  resolution = 1,
  target_clusters = NULL,
  max_iterations = 20,
  tolerance = 1,
  factor_name = "Louvain_clust",
  weights = NULL,
  lock.seed = TRUE,
  seed = 42,
  use_approximate = TRUE,
  ...
)
\end{verbatim}
\end{Usage}
%
\begin{Arguments}
\begin{ldescription}
\item[\code{romics\_object}] A romics\_object created using romicsCreateObject().

\item[\code{cluster\_using}] Character string specifying which data to use for clustering. Options: 'data', 'pca', 'umap', 'tsne'.
Default: 'data'

\item[\code{k}] Number of nearest neighbors to use in graph construction. Default: 20
For datasets >50k samples, consider reducing this value (e.g., k=10) to improve memory efficiency.

\item[\code{resolution}] Resolution parameter that determines the granularity of the clustering.
Higher values lead to more clusters. Default: 1.0

\item[\code{target\_clusters}] Optional integer specifying the desired number of clusters. If provided, the function will
iteratively adjust the resolution parameter to approximate this number. Default: NULL

\item[\code{max\_iterations}] Maximum number of iterations when trying to reach target\_clusters. Default: 20

\item[\code{tolerance}] Tolerance for target\_clusters - clustering will stop if within this range. Default: 1

\item[\code{factor\_name}] Character string specifying the name for the new factor to be added to metadata.
Default: 'Louvain\_clust'

\item[\code{weights}] Optional edge weights for the graph. If NULL, distance-based weighting is used automatically.

\item[\code{lock.seed}] Logical. If TRUE, uses a fixed seed for reproducible clustering. Default: TRUE

\item[\code{seed}] Numeric value for the random seed when lock.seed=TRUE. Default: 42

\item[\code{use\_approximate}] Logical. If TRUE and FNN available, uses fast approximate KNN. Default: TRUE

\item[\code{...}] Additional parameters passed to igraph::cluster\_louvain().
\end{ldescription}
\end{Arguments}
%
\begin{Details}
OPTIMIZED IMPLEMENTATION:
- Pre-allocates edge matrices instead of vector concatenation (10-50x faster)
- Reuses KNN distances directly instead of recalculating (5-10x faster)
- Vectorizes weight calculations (3-5x faster)
- Overall expected speedup: 20-100x on large datasets

Uses the same interface as romicsLouvainSamples\_old but with significantly improved performance.
\end{Details}
%
\begin{Value}
Returns the romics\_object with cluster assignments added to metadata.
\end{Value}
%
\begin{Author}
Geremy Clair
\end{Author}
\HeaderA{romicsLouvainSamples\_old}{romicsLouvainSamples()}{romicsLouvainSamples.Rul.old}
%
\begin{Description}
Generate Louvain community clustering for samples in a romics\_object using memory-efficient algorithms.
The clustering can be based on the full data matrix or dimensionally reduced data from PCA, UMAP, or t-SNE embeddings.
For large datasets (>10k samples), this function uses optimized nearest neighbor search to reduce memory usage.

DEPRECATED: Original (slow) version of Louvain clustering. Use romicsLouvainSamples() instead for significantly better performance.
This function is kept for backward compatibility only.
\end{Description}
%
\begin{Usage}
\begin{verbatim}
romicsLouvainSamples_old(
  romics_object,
  cluster_using = c("data", "pca", "umap", "tsne"),
  k = 20,
  resolution = 1,
  target_clusters = NULL,
  max_iterations = 20,
  tolerance = 1,
  factor_name = "Louvain_clust",
  weights = NULL,
  lock.seed = TRUE,
  seed = 42,
  use_approximate = TRUE,
  ...
)
\end{verbatim}
\end{Usage}
%
\begin{Arguments}
\begin{ldescription}
\item[\code{romics\_object}] A romics\_object created using romicsCreateObject().

\item[\code{cluster\_using}] Character string specifying which data to use for clustering.

\item[\code{k}] Number of nearest neighbors to use in graph construction.

\item[\code{resolution}] Resolution parameter for the clustering.

\item[\code{target\_clusters}] Optional integer specifying the desired number of clusters.

\item[\code{max\_iterations}] Maximum number of iterations when trying to reach target\_clusters.

\item[\code{tolerance}] Tolerance for target\_clusters.

\item[\code{factor\_name}] Character string specifying the name for the new factor.

\item[\code{weights}] Optional edge weights for the graph.

\item[\code{lock.seed}] Logical. If TRUE, uses a fixed seed for reproducible clustering.

\item[\code{seed}] Numeric value for the random seed.

\item[\code{use\_approximate}] Logical. If TRUE and FNN available, uses fast approximate KNN.

\item[\code{...}] Additional parameters passed to igraph::cluster\_louvain().
\end{ldescription}
\end{Arguments}
%
\begin{Details}
This function performs Louvain community detection on samples and adds the cluster assignments
as a new factor in the romics\_object's metadata. The function constructs a k-nearest neighbor
graph from the data using memory-efficient algorithms, then applies the Louvain algorithm to identify communities.

When target\_clusters is specified, the function will iteratively adjust the resolution parameter
to try to achieve approximately that number of clusters. Note that exact control is not guaranteed
due to the nature of the Louvain algorithm.

For large datasets, the function automatically uses optimized approaches:
- FNN package for fast approximate nearest neighbor search (if available)
- Edge list representation instead of full adjacency matrices
- Chunked processing to control memory usage
- Periodic garbage collection to free memory

The Louvain algorithm is a popular method for community detection in large networks,
optimizing modularity to find densely connected groups of nodes.

Memory recommendations for large datasets:
- Use dimensionally reduced data (cluster\_using = "pca" or "umap")
- Reduce k parameter (e.g., k=10 for >100k samples)
- Install FNN package: install.packages("FNN")
- Ensure sufficient system memory (recommend >16GB for >100k samples)
\end{Details}
%
\begin{Value}
Returns the romics\_object with a new factor added to metadata containing the cluster assignments.
Clusters are named as "factor\_name\_001", "factor\_name\_002", etc.
A message will indicate the number of clusters found and confirm the factor was added.

Returns the romics\_object with cluster assignments added to metadata.
\end{Value}
%
\begin{Author}
Geremy Clair
romicsLouvainSamples\_old()

Geremy Clair
\end{Author}
\HeaderA{romicsMean}{romicsMean()}{romicsMean}
%
\begin{Description}
Calculates the means of each variable within each level of the selected factor and add the generated columns in the statistics layer of the romics\_object.
\end{Description}
%
\begin{Usage}
\begin{verbatim}
romicsMean(romics_object, factor = "main")
\end{verbatim}
\end{Usage}
%
\begin{Arguments}
\begin{ldescription}
\item[\code{romics\_object}] A object created using the function romicsCreateObject().

\item[\code{main\_factor}] Either 'main' OR any factor from the romics\_object, the list of factors from a romics object can be obtained using the function romicsFactorNames().
\end{ldescription}
\end{Arguments}
%
\begin{Details}
Adds the means columns to the statistics layer. enable to choose a different factor than the main one to do those calculation
\end{Details}
%
\begin{Value}
This function returns a modified romics object, containing mean columns in the statistics layer.
\end{Value}
%
\begin{Author}
Geremy Clair
\end{Author}
\HeaderA{romicsMeanByLevel}{romicsMeanByLevel()}{romicsMeanByLevel}
%
\begin{Description}
The function will condensate the values based on a romics factor, it will reduce the number of sample to one by level using the mean for each feature.
\end{Description}
%
\begin{Usage}
\begin{verbatim}
romicsMeanByLevel(romics_object, averaging_factor, main_factor)
\end{verbatim}
\end{Usage}
%
\begin{Arguments}
\begin{ldescription}
\item[\code{romics\_object}] has to be an romics\_object created using romicsCreateObject()

\item[\code{averaging\_factor}] as to be a factor of an romics\_object factors can be identified using the function romicsFactorNames()

\item[\code{main\_factor}] as to be either "main" or a factor of the romics\_object that can be identified using the function romicsFactorNames()
\end{ldescription}
\end{Arguments}
%
\begin{Details}
The function will average the data inside the factor, NA will be not considered. a new metadata file will be created, factors with discrepancies while the averaging is done will be removed.
\end{Details}
%
\begin{Value}
This function returns the transformed romics\_object with updated data layer
\end{Value}
%
\begin{Author}
Geremy Clair, Igor Estevao
\end{Author}
\HeaderA{romicsMetadataRelations}{romicsMetadataRelations()}{romicsMetadataRelations}
%
\begin{Description}
display relationships between different factors contained in the metadata layers of an Romics\_object
\end{Description}
%
\begin{Usage}
\begin{verbatim}
romicsMetadataRelations(
  romics_object = romics_object,
  factors = c("factor1", "factor2", "factor3"),
  textsize = 3,
  colors = NULL
)
\end{verbatim}
\end{Usage}
%
\begin{Arguments}
\begin{ldescription}
\item[\code{romics\_object}] A romics\_object created using romicsCreateObject().

\item[\code{factors}] A character vector (length>2) containing the list of factors of the romics\_object. the list of factors can be identified by using the function romicsFactorNames()

\item[\code{textsize}] Size of the text labels (default: 3)

\item[\code{colors}] Optional vector of colors for the strata. If NULL, colors are generated automatically.
\end{ldescription}
\end{Arguments}
%
\begin{Details}
generates a plot displaying relationships between different factors
\end{Details}
%
\begin{Value}
a ggplot
\end{Value}
%
\begin{Author}
Geremy Clair
\end{Author}
\HeaderA{romicsMissingFeatures}{romicsMissingFeatures()}{romicsMissingFeatures}
%
\begin{Description}
returns a table containing the missingness of each feature contained in the romics\_object.
\end{Description}
%
\begin{Usage}
\begin{verbatim}
romicsMissingFeatures(
  romics_object = romics_object,
  reorder = c(NULL, "descending", "ascending")
)
\end{verbatim}
\end{Usage}
%
\begin{Arguments}
\begin{ldescription}
\item[\code{romics\_object}] has to be an romics\_object created using romicsCreateObject()

\item[\code{reorder}] has to be either empty (original order will be used), 'ascending', or 'descending' to indicate in what order the features should be displayed on the plot.
\end{ldescription}
\end{Arguments}
%
\begin{Details}
This function does not alter the romics\_object, it generates a table containing the percentage of missingness for each feature of the romics\_object
\end{Details}
%
\begin{Value}
This function will return a ggplot2 geom\_bar plot. it can then be further visually adjusted using the ggplot2 commands.
\end{Value}
%
\begin{Author}
Geremy Clair
\end{Author}
\HeaderA{romicsNameClusters}{romicsNameClusters()}{romicsNameClusters}
%
\begin{Description}
Rename clusters in any factor of a romics\_object's metadata.
The original factor values are preserved in a new row with "\_original" appended to the factor name.
Optionally create a new factor instead of modifying the existing one.
\end{Description}
%
\begin{Usage}
\begin{verbatim}
romicsNameClusters(
  romics_object,
  cluster_factor = NULL,
  cluster_new_names = NULL,
  new_factor_name = NULL,
  alphabetical_mapping = FALSE
)
\end{verbatim}
\end{Usage}
%
\begin{Arguments}
\begin{ldescription}
\item[\code{romics\_object}] A romics\_object created using romicsCreateObject()

\item[\code{cluster\_factor}] Character string specifying the name of an existing factor in metadata containing the clusters to rename.
Must be a valid factor name that can be verified with romicsFactorNames().
If missing, defaults to the main factor of the romics\_object.

\item[\code{cluster\_new\_names}] Optional character vector containing the new names for each cluster.
If provided, must have the same length as the number of unique clusters.
If NULL (default), the function will interactively prompt for names for each cluster.

\item[\code{new\_factor\_name}] Optional character string. If provided, creates a new factor with this name instead of modifying the existing factor.
The original factor remains unchanged. If NULL (default), modifies the existing factor.

\item[\code{alphabetical\_mapping}] Logical. If TRUE and cluster\_new\_names is provided, maps new names alphabetically to original clusters.
If FALSE (default), maps based on the order of appearance or natural sorting of original clusters.
\end{ldescription}
\end{Arguments}
%
\begin{Value}
Returns the romics\_object with renamed clusters and the original clusters preserved (if modifying existing factor)
\end{Value}
%
\begin{Author}
Geremy Clair
\end{Author}
\HeaderA{romicsNumericHeaderAdjust}{romicsNumericHeaderAdjust()}{romicsNumericHeaderAdjust}
%
\begin{Description}
Adjusts sample names in a romics\_object by adding a prefix to prevent issues with numeric-only sample names.
This function is particularly useful when sample names are purely numeric, which can cause issues in R data processing,
plotting, and statistical analysis. The function adds a specified prefix to all sample names across all layers of the romics\_object.
\end{Description}
%
\begin{Usage}
\begin{verbatim}
romicsNumericHeaderAdjust(
  romics_object,
  prefix = "sample",
  separator = "_",
  force_adjustment = FALSE
)
\end{verbatim}
\end{Usage}
%
\begin{Arguments}
\begin{ldescription}
\item[\code{romics\_object}] A romics\_object created using createRomicsObject()

\item[\code{prefix}] A character string to be added as a prefix to all sample names. Default is "sample"

\item[\code{separator}] A character string used to separate the prefix from the original sample name. Default is "\_"

\item[\code{force\_adjustment}] Logical. If TRUE, adds prefix to all sample names regardless of whether they are numeric.
If FALSE (default), only adds prefix to purely numeric sample names.
\end{ldescription}
\end{Arguments}
%
\begin{Details}
This function modifies sample names (column names) across all relevant layers of the romics\_object
\end{Details}
%
\begin{Value}
Returns the romics\_object with adjusted sample names. A message will indicate
how many sample names were modified.
\end{Value}
%
\begin{Author}
Geremy Clair
\end{Author}
\HeaderA{romicsOutlierEval}{romics\_outlier\_eval()}{romicsOutlierEval}
%
\begin{Description}
Plots the the pmartR method for the evaluation of outliers (require the installation of the package pmartR). To remove the outlier below a certain treshold use the function Romics\_outlier\_remove()
\end{Description}
%
\begin{Usage}
\begin{verbatim}
romicsOutlierEval(
  romics_object,
  seed = 42,
  metrics = c("Correlation", "Proportion_Missing", "MAD", "Skewness"),
  pvalue_threshold = 0.01,
  label = FALSE
)
\end{verbatim}
\end{Usage}
%
\begin{Arguments}
\begin{ldescription}
\item[\code{romics\_object}] A log transformed romics\_object created using romicsCreateObject() and transformed using the function log2transform() or log10transform()

\item[\code{seed}] An integer of length 1, by default 42 will be used

\item[\code{metrics}] A character vector containing the following terms to indicate which parameters to use for the filtering of the data : 'Correlation', 'Proportion\_Missing','MAD', 'Skewness'. By defaults all parameters will be used.

\item[\code{pvalue\_threshold}] A numeric vector of lenght 1 indicating the pvalue threshold to be used.

\item[\code{label}] Either TRUE or FALSE to indicate if the sample labels have to be plotted. Default: FALSE
\end{ldescription}
\end{Arguments}
%
\begin{Details}
This function requires the package 'pmartR' to be installed and loaded to be excecuted. It will calculate and plot the samples to be filtered out using the function Romics\_outlier\_eval().
\end{Details}
%
\begin{Value}
This function will print the pmartR filtering details and will return 2 plots the first one is a scatter plot of the pvalue by log2(Robust Mahalanobis Distance) the second one is a scatter plot of the log2(Robust Mahalanobis Distance) per sample.
\end{Value}
%
\begin{Author}
Geremy Clair
\end{Author}
%
\begin{References}
Matzke, M., Waters, K., Metz, T., Jacobs, J., Sims, A., Baric, R., Pounds, J., and Webb-Robertson, B.J. (2011), Improved quality control processing of peptide-centric LC-MS proteomics data. Bioinformatics. 27(20): 2866-2872.
\end{References}
\HeaderA{romicsOutlierRemove}{Romics\_outlier\_remove()}{romicsOutlierRemove}
%
\begin{Description}
Removes outliers below a certain pvalue threshold. The pmartR method for the removal of outliers is used (require the installation of the package pmartR). The evaluation of the outlier removal can be done prior to apply this function using the function Romics\_outlier\_eval().
\end{Description}
%
\begin{Usage}
\begin{verbatim}
romicsOutlierRemove(
  romics_object,
  seed = 42,
  metrics = c("Correlation", "Proportion_Missing", "MAD", "Skewness"),
  pvalue_threshold = 0.01
)
\end{verbatim}
\end{Usage}
%
\begin{Arguments}
\begin{ldescription}
\item[\code{romics\_object}] A log transformed romics\_object created using romicsCreateObject() and transformed using the function log2transform() or log10transform()

\item[\code{seed}] An integer of length 1, by default 42 will be used

\item[\code{metrics}] A character vector containing the following terms to indicate which parameters to use for the filtering of the data : 'Correlation', 'Proportion\_Missing','MAD', 'Skewness'. By defaults all parameters will be used.

\item[\code{pvalue\_threshold}] A numeric vector of lenght 1 indicating the pvalue threshold to be used.
\end{ldescription}
\end{Arguments}
%
\begin{Details}
This function requires the package 'pmartR' to be installed and loaded to be excecuted. It will calculate and plot the samples to be filtered out using the function Romics\_outlier\_eval().
\end{Details}
%
\begin{Value}
This function will print the pmartR filtering details and will return 2 plots the first one is a scatter plot of the pvalue by log2(Robust Mahalanobis Distance) the second one is a scatter plot of the log2(Robust Mahalanobis Distance) per sample.
\end{Value}
%
\begin{Author}
Geremy Clair
\end{Author}
%
\begin{References}
Matzke, M., Waters, K., Metz, T., Jacobs, J., Sims, A., Baric, R., Pounds, J., and Webb-Robertson, B.J. (2011), Improved quality control processing of peptide-centric LC-MS proteomics data. Bioinformatics. 27(20): 2866-2872.
\end{References}
\HeaderA{romicsPCA}{romicsPCA()}{romicsPCA}
%
\begin{Description}
Calculate the PCA of the data layer of the romics\_object using the package FactoMineR.
If the data layer contains some missing values those will be imputed using the missMDA::imputePCA() method.
This function updates the romics\_object by storing the PCA embeddings in the embeddings layer.
\end{Description}
%
\begin{Usage}
\begin{verbatim}
romicsPCA(
  romics_object,
  verbose = TRUE,
  ncp,
  n_components = NULL,
  scale = TRUE,
  method = c("Regularized", "EM"),
  ncp.min = 0,
  ncp.max = 5,
  method.cv,
  ...
)
\end{verbatim}
\end{Usage}
%
\begin{Arguments}
\begin{ldescription}
\item[\code{romics\_object}] has to be a romics\_object created using romicsCreateObject()

\item[\code{verbose}] boolean. If TRUE, also assigns the complete PCA results to PCA\_results in the global environment

\item[\code{ncp}] number of principal components to compute. Default is 5 if method.cv is not provided.
If method.cv is provided but ncp is not, the optimal number is estimated automatically.

\item[\code{n\_components}] number of principal components to store in the embeddings layer. If NULL (default),
stores all computed components. Must be <= ncp.

\item[\code{scale}] boolean. TRUE implies a same weight for each variable for the missMDA data imputation)

\item[\code{method}] "Regularized" (by default) or "EM". Used for missing value imputation using missMDA.

\item[\code{ncp.min}] used only if ncp is automatically estimated. Integer corresponding to the minimum number of components to test

\item[\code{ncp.max}] used only if ncp is automatically estimated. Integer corresponding to the maximum number of components to test

\item[\code{method.cv}] method for cross-validation when automatically determining the optimal number of components.
One of "gcv" (generalized cross-validation), "loo" (leave-one-out), or "Kfold" (K-fold cross-validation).
If provided and ncp is missing, triggers automatic estimation of the optimal number of components.

\item[\code{...}] further arguments passed to missMDA functions.

\item[\code{row.w}] row weights (by default, a vector of 1 for uniform row weights)
\end{ldescription}
\end{Arguments}
%
\begin{Details}
This function calculates PCA and stores the sample embeddings in the romics\_object\$embeddings layer.
If missing values are present in the data, they are imputed using missMDA::imputePCA().
When method.cv is provided but ncp is not, the optimal number of components is estimated using missMDA::estim\_ncpPCA().
If verbose=TRUE, the complete PCA results are assigned to PCA\_results in the global environment.
\end{Details}
%
\begin{Value}
Returns the romics\_object with PCA results added to the embeddings layer.
\end{Value}
%
\begin{Author}
Geremy Clair
\end{Author}
\HeaderA{romicsPCAplot}{romicsPCAplot()}{romicsPCAplot}
%
\begin{Description}
Plots the samples on the PCA sample plot in 2D or 3D. When in 2D mode, can also show percentage of explained variance.
The plot can be colored by any factor in the romics\_object's metadata or by feature values.
This function uses the PCA embeddings stored in the romics\_object and checks for PCA\_results in the global environment for variance plots.
\end{Description}
%
\begin{Usage}
\begin{verbatim}
romicsPCAplot(
  romics_object,
  Xcomp = 1,
  Ycomp = 2,
  Zcomp = NULL,
  label = FALSE,
  plot_type = "dual",
  feature = NULL,
  factor_name = NULL,
  color_palette = NULL,
  standardize = TRUE,
  size = 3,
  alpha = 0.8,
  plotly = FALSE,
  show_outline = TRUE
)
\end{verbatim}
\end{Usage}
%
\begin{Arguments}
\begin{ldescription}
\item[\code{romics\_object}] A romics\_object that has PCA embeddings in its embeddings layer (generated using romicsPCA)

\item[\code{Xcomp}] numerical/double. Indicate the component to plot on the X axis (default: 1)

\item[\code{Ycomp}] numerical/double. Indicate the component to plot on the Y axis (default: 2)

\item[\code{Zcomp}] numerical/double. If provided, creates a 3D plot with this component on the Z axis (default: NULL)

\item[\code{label}] boolean. Indicate if the sample name labels should be plotted (default: FALSE, applies only to 2D plots)

\item[\code{plot\_type}] should be one of the following options to indicate the type of plot to be returned:
'dual' (for both sample plot and variance plot), 'individual', or 'percentage' (default: "dual").
Note: The 'dual' and 'percentage' options only work for 2D plots (when Zcomp is NULL).

\item[\code{feature}] character string. If provided, the plot will be colored according to the values of this feature.
If NULL (default), coloring will be based on the factor specified by factor\_name.

\item[\code{factor\_name}] character string. The name of the factor in romics\_object\$metadata to use for coloring the points
when feature is NULL. If NULL (default), the romics\_object's main\_factor will be used.

\item[\code{color\_palette}] a color palette to use. For categorical factors, should be a vector with one color per level.
For continuous features (when feature is specified), should be a gradient color palette.
Default: NULL (uses colors\_romics for factors, viridis for features).

\item[\code{standardize}] boolean. Indicates if the feature values should be scaled when coloring by feature (default: TRUE)

\item[\code{size}] numeric. Size of the points in the plot (default: 3)

\item[\code{alpha}] numeric. Transparency of the points in the plot (default: 0.8)

\item[\code{plotly}] boolean. If TRUE, generates an interactive plotly plot instead of a static ggplot (default: FALSE)

\item[\code{show\_outline}] boolean. If FALSE, removes point outlines and shows only fill colors (default: TRUE)
\end{ldescription}
\end{Arguments}
%
\begin{Details}
This function plots PCA results using the embeddings stored in the romics\_object.
If Zcomp is provided, a 3D interactive plot is generated using plotly. In 3D mode, only the sample plot
is available (plot\_type settings are ignored).
If Zcomp is NULL, a 2D plot is generated using ggplot2, and can include the variance plot depending on plot\_type.
For the variance percentage plot (2D mode only), the function looks for PCA\_results in the global environment.
When using a non-main factor for coloring, the function automatically updates the romics\_object
using romicsChangeFactor() to ensure consistent colors across plots.
Coloring is determined automatically:
- If feature is provided, the plot is colored by feature values
- If feature is NULL, the plot is colored by factor values
When show\_outline is FALSE, points are displayed without borders for cleaner appearance.
\end{Details}
%
\begin{Value}
Returns either a ggplot2 plot, a grid.arrange combined plot (for 2D dual plots), or a plotly 3D plot.
\end{Value}
%
\begin{Author}
Geremy Clair
\end{Author}
\HeaderA{romicsPercentComplete}{romicsPercentComplete()}{romicsPercentComplete}
%
\begin{Description}
Calculates the percentage of completeness for each level of the selected factor and add the generated columns in the statistics layer of the romics\_object.
\end{Description}
%
\begin{Usage}
\begin{verbatim}
romicsPercentComplete(romics_object, factor = "main")
\end{verbatim}
\end{Usage}
%
\begin{Arguments}
\begin{ldescription}
\item[\code{romics\_object}] A object created using the function romicsCreateObject().

\item[\code{main\_factor}] Either 'main' OR any factor from the romics\_object, the list of factors from a romics object can be obtained using the function romicsFactorNames().
\end{ldescription}
\end{Arguments}
%
\begin{Details}
Adds the percentage\_completeness columns to the statistics layer. Enable to choose a different factor than the main one to do those calculation.
\end{Details}
%
\begin{Value}
This function returns a modified romics object, containing percentage\_completeness columns in the statistics layer.
\end{Value}
%
\begin{Author}
Geremy Clair
\end{Author}
\HeaderA{romicsPercentMissing}{romicsPercentMissing()}{romicsPercentMissing}
%
\begin{Description}
Calculates the percentage of Missingness for each level of the selected factor and add the generated columns in the statistics layer of the romics\_object.
\end{Description}
%
\begin{Usage}
\begin{verbatim}
romicsPercentMissing(romics_object, factor = "main")
\end{verbatim}
\end{Usage}
%
\begin{Arguments}
\begin{ldescription}
\item[\code{romics\_object}] A object created using the function romicsCreateObject().

\item[\code{main\_factor}] Either 'main' OR any factor from the romics\_object, the list of factors from a romics object can be obtained using the function romicsFactorNames().
\end{ldescription}
\end{Arguments}
%
\begin{Details}
Adds the percentage\_missingness columns to the statistics layer. Enable to choose a different factor than the main one to do those calculation.
\end{Details}
%
\begin{Value}
This function returns a modified romics object, containing percentage\_missingness columns in the statistics layer.
\end{Value}
%
\begin{Author}
Geremy Clair
\end{Author}
\HeaderA{romicsPlotMissing}{romicsPlotMissing()}{romicsPlotMissing}
%
\begin{Description}
Plots the missingness of each sample contained in the romics\_object. The colors used for the plotting will correspond to the main\_factor of the romics\_object.
\end{Description}
%
\begin{Usage}
\begin{verbatim}
romicsPlotMissing(romics_object, custom_colors = "colorlist")
\end{verbatim}
\end{Usage}
%
\begin{Arguments}
\begin{ldescription}
\item[\code{romics\_object}] has to be an romics\_object created using romicsCreateObject()

\item[\code{main\_factor}] has to be either "main", "none" or a factor of an romics\_object created using romicsCreateObject() the list of factors can be obtained by running the function romicsFactorNames()
\end{ldescription}
\end{Arguments}
%
\begin{Details}
This function does not alter the romics\_object, it plots the the missingness of each sample in a barplot.
\end{Details}
%
\begin{Value}
This function will return a ggplot2 geom\_bar plot. it can then be further visually adjusted using the ggplot2 commands
\end{Value}
%
\begin{Author}
Geremy Clair
\end{Author}
\HeaderA{romicsPlotMissingFeatures}{romicsPlotMissingFeatures()}{romicsPlotMissingFeatures}
%
\begin{Description}
Plots the missingness of each feature contained in the romics\_object.
\end{Description}
%
\begin{Usage}
\begin{verbatim}
romicsPlotMissingFeatures(
  romics_object = romics_object,
  features = c(feature_a, feature_b),
  reorder = c(NULL, "descending", "ascending")
)
\end{verbatim}
\end{Usage}
%
\begin{Arguments}
\begin{ldescription}
\item[\code{romics\_object}] has to be an romics\_object created using romicsCreateObject()

\item[\code{features}] has to be either empty or a character vector containing a list of features. if a character list of vector is used only the requested features will be shown.

\item[\code{reorder}] has to be either empty (original order will be used), 'ascending', or 'descending' to indicate in what order the features should be displayed on the plot.
\end{ldescription}
\end{Arguments}
%
\begin{Details}
This function does not alter the romics\_object, it plots the the missingness of each feature for all the samples in a barplot.
\end{Details}
%
\begin{Value}
This function will return a ggplot2 geom\_bar plot. it can then be further visually adjusted using the ggplot2 commands.
\end{Value}
%
\begin{Author}
Geremy Clair
\end{Author}
\HeaderA{romicsPlotPresent}{romicsPlotPresent()}{romicsPlotPresent}
%
\begin{Description}
Plots the percentage of proteins with valid (non-missing) values for each sample in the romics\_object. The colors used for the plotting will correspond to the main\_factor of the romics\_object.
\end{Description}
%
\begin{Usage}
\begin{verbatim}
romicsPlotPresent(romics_object, custom_colors = "colorlist")
\end{verbatim}
\end{Usage}
%
\begin{Arguments}
\begin{ldescription}
\item[\code{romics\_object}] has to be an romics\_object created using romicsCreateObject()

\item[\code{custom\_colors}] has to be either "colorlist" or a character vector containing the colors to be used for the plotting
\end{ldescription}
\end{Arguments}
%
\begin{Details}
This function does not alter the romics\_object, it plots the percentage of valid data for each sample in a barplot. This is complementary to romicsPlotMissing() and shows data availability instead of data absence.
\end{Details}
%
\begin{Value}
This function will return a ggplot2 geom\_bar plot. it can then be further visually adjusted using the ggplot2 commands
\end{Value}
%
\begin{Author}
Geremy Clair
\end{Author}
\HeaderA{romicsPlotSignificantFeatures}{romicsPlotSignificantFeatures()}{romicsPlotSignificantFeatures}
%
\begin{Description}
Create a barplot showing the number of significant features for each statistical test
\end{Description}
%
\begin{Usage}
\begin{verbatim}
romicsPlotSignificantFeatures(
  romics_object,
  p_threshold = 0.05,
  padj_type = "p",
  specific_comparisons = NULL,
  test_types = NULL,
  split_directions = FALSE,
  fc_threshold = 0,
  colors = NULL,
  title = NULL,
  rotate_x_labels = TRUE
)
\end{verbatim}
\end{Usage}
%
\begin{Arguments}
\begin{ldescription}
\item[\code{romics\_object}] A romics\_object with calculated statistics

\item[\code{p\_threshold}] Numeric value for significance threshold. Default: 0.05

\item[\code{padj\_type}] Character string indicating which p-value type to use: "p" for raw p-values or "padj" for adjusted p-values. Default: "p"

\item[\code{specific\_comparisons}] Character vector of specific comparisons to consider (e.g., "Alport\_vs\_Ctrl"). If NULL, considers all comparisons. Default: NULL

\item[\code{test\_types}] Character vector of test types to consider (e.g., "Ttest", "Wilcox\_test", "glmBinomialTest"). If NULL, considers all tests. Default: NULL

\item[\code{split\_directions}] Logical indicating whether to split up and downregulated features into separate bars. Default: FALSE

\item[\code{fc\_threshold}] Numeric value for fold change threshold (used with directionality). Default: 0

\item[\code{colors}] Character vector of colors for bars. If split\_directions=TRUE, provide colors for c("Up", "Down"). If FALSE, provide single color. Default: c("\#E31A1C", "\#1F78B4") for up/down or "\#2cbcb2" for total

\item[\code{title}] Character string for plot title. If NULL, generates automatic title. Default: NULL

\item[\code{rotate\_x\_labels}] Logical indicating whether to rotate x-axis labels 45 degrees. Default: TRUE
\end{ldescription}
\end{Arguments}
%
\begin{Details}
This function uses romicsExtractSignificantFeatures() to count significant features for each test type
and creates a barplot using ggplot2 with theme\_ROP(). Can show total counts or split by up/down regulation.
\end{Details}
%
\begin{Value}
ggplot object
\end{Value}
%
\begin{Author}
[Your Name]
\end{Author}
\HeaderA{romicsPmartR}{romicsPmartR()}{romicsPmartR}
%
\begin{Description}
Converts romics\_object to pmartR object (requires the "pmartR' package to be installed).
\end{Description}
%
\begin{Usage}
\begin{verbatim}
romicsPmartR(romics_object, type = "proData")
\end{verbatim}
\end{Usage}
%
\begin{Arguments}
\begin{ldescription}
\item[\code{romics\_object}] A romics\_object created using the function romicsCreateObject()

\item[\code{type}] Has to be "lipidData","proData","pepData","metabData" to indicate what data type to use for the pmartR object
\end{ldescription}
\end{Arguments}
%
\begin{Details}
This function converts an romics\_object to a pmartR object
\end{Details}
%
\begin{Value}
return the pmartR object
\end{Value}
%
\begin{Author}
Geremy Clair
\end{Author}
\HeaderA{romicsProportionHeatmap}{romicsProportionHeatmap}{romicsProportionHeatmap}
%
\begin{Description}
Plots a heatmap of signed proportion differences (
output of romicsProportionTestPairwise(). Rows are factor levels
ordered alphabetically, columns are pairwise comparisons. Fill colour
encodes Signed\_Diff (proportion in cond1 minus proportion in cond2).
Each tile displays the rounded signed difference followed by a significance level.
\end{Description}
%
\begin{Usage}
\begin{verbatim}
romicsProportionHeatmap(
  proportion_test_result,
  title = "Proportion Difference Heatmap",
  midpoint = 0,
  low_color = "blue",
  mid_color = "white",
  high_color = "red",
  label_size = 3,
  text_size = 11,
  interactive = FALSE
)
\end{verbatim}
\end{Usage}
%
\begin{Arguments}
\begin{ldescription}
\item[\code{proportion\_test\_result}] The list returned by romicsProportionTestPairwise.

\item[\code{title}] A character string for the plot title. Default:
"Proportion Difference Heatmap".

\item[\code{midpoint}] Numeric. Centre of the diverging colour scale. Default: 0.

\item[\code{low\_color}] Colour for negative differences (depleted in cond1).
Default: "blue".

\item[\code{mid\_color}] Colour for zero difference. Default: "white".

\item[\code{high\_color}] Colour for positive differences (enriched in cond1).
Default: "red".

\item[\code{label\_size}] Numeric. Font size of the tile labels. Default: 3.

\item[\code{text\_size}] Numeric. Base font size for the plot. Default: 11.

\item[\code{interactive}] Logical. If TRUE returns a plotly object; if FALSE
returns a ggplot2 object. Default: FALSE.
\end{ldescription}
\end{Arguments}
%
\begin{Value}
A ggplot2 object, or a plotly object if interactive = TRUE.
\end{Value}
%
\begin{Author}
Geremy Clair
\end{Author}
\HeaderA{romicsPseudoBulk}{romicsPseudoBulk}{romicsPseudoBulk}
%
\begin{Description}
Creates pseudobulk samples by aggregating observations within groups
\end{Description}
%
\begin{Usage}
\begin{verbatim}
romicsPseudoBulk(
  romics_object,
  bulk_by,
  method = "mean",
  min_observations = 3,
  keep_factors = NULL,
  suffix = "_bulk"
)
\end{verbatim}
\end{Usage}
%
\begin{Arguments}
\begin{ldescription}
\item[\code{romics\_object}] A romics\_object created using romicsCreateObject()

\item[\code{bulk\_by}] Factor name to aggregate by

\item[\code{method}] Aggregation method: "mean", "median", or "sum". Default: "mean"

\item[\code{min\_observations}] Minimum observations required per group. Default: 3

\item[\code{keep\_factors}] Character vector of factor names to preserve in metadata. Default: NULL

\item[\code{suffix}] Suffix to add to sample names. Default: "\_bulk"
\end{ldescription}
\end{Arguments}
%
\begin{Value}
A new romics\_object with pseudobulk data
\end{Value}
%
\begin{Author}
Geremy Clair
\end{Author}
\HeaderA{romicsRemoveEmbeddings}{romicsRemoveEmbeddings()}{romicsRemoveEmbeddings}
%
\begin{Description}
Remove embeddings from a romics\_object to reduce memory usage.
\end{Description}
%
\begin{Usage}
\begin{verbatim}
romicsRemoveEmbeddings(romics_object, embeddings = NULL)
\end{verbatim}
\end{Usage}
%
\begin{Arguments}
\begin{ldescription}
\item[\code{romics\_object}] A romics\_object containing embeddings to be removed

\item[\code{embeddings}] Character vector of embedding names to remove. If NULL, removes all embeddings.
\end{ldescription}
\end{Arguments}
%
\begin{Value}
Returns the romics\_object with specified embeddings removed
\end{Value}
%
\begin{Author}
Geremy Clair
\end{Author}
\HeaderA{romicsRemoveFactor}{romicsRemoveFactor}{romicsRemoveFactor}
%
\begin{Description}
This function removes one or more factors from the metadata of a romics object.
\end{Description}
%
\begin{Usage}
\begin{verbatim}
romicsRemoveFactor(romics_object, factors_to_remove, force = FALSE)
\end{verbatim}
\end{Usage}
%
\begin{Arguments}
\begin{ldescription}
\item[\code{romics\_object}] A romics\_object created using romicsCreateObject().

\item[\code{factors\_to\_remove}] A character vector containing the name(s) of factor(s) to remove. The list of factor names can be identified by using the function romicsFactorNames().

\item[\code{force}] A logical indicating whether to allow removal of the main factor. Default: FALSE (prevents accidental removal of main factor).
\end{ldescription}
\end{Arguments}
%
\begin{Details}
This function removes specified factors from the metadata layer of a romics object. By default, it prevents removal of the main factor to avoid breaking the object structure. Set force=TRUE to override this protection.
\end{Details}
%
\begin{Value}
A romics\_object with the specified factor(s) removed from the metadata.
\end{Value}
%
\begin{Author}
Geremy Clair
\end{Author}
%
\begin{Examples}
\begin{ExampleCode}
## Not run: 
# Remove a single factor
romics_object <- romicsRemoveFactor(romics_object, factors_to_remove = "old_factor")

# Remove multiple factors
romics_object <- romicsRemoveFactor(romics_object, factors_to_remove = c("factor1", "factor2"))

# Remove main factor (use with caution!)
romics_object <- romicsRemoveFactor(romics_object, factors_to_remove = "condition", force = TRUE)

## End(Not run)
\end{ExampleCode}
\end{Examples}
\HeaderA{romicsRemoveUnknown}{romicsRemoveUnknown()}{romicsRemoveUnknown}
%
\begin{Description}
Remove list of unkown features from the dataset
\end{Description}
%
\begin{Usage}
\begin{verbatim}
romicsRemoveUnknown(romics_object, unknown_identifier_txt = "Unknown")
\end{verbatim}
\end{Usage}
%
\begin{Arguments}
\begin{ldescription}
\item[\code{romics\_object}] has to be an romics\_object created using romicsCreateObject()

\item[\code{unknown\_identifier\_txt}] string that indicates which features are unknown and are to be removed by default the string 'Unknown'
\end{ldescription}
\end{Arguments}
%
\begin{Details}
This function will create a new romics object in which only the unkown features are removed from the data, missingdata and statistics layers.

The function removes unkown features
\end{Details}
%
\begin{Value}
This function returns an romics\_object
\end{Value}
%
\begin{Author}
Geremy Clair
\end{Author}
\HeaderA{romicsRestoreMissing}{romicsRestoreMissing()}{romicsRestoreMissing}
%
\begin{Description}
Restores the missingness of the data layer of the romics\_object based on the content of the missingdata layer.
\end{Description}
%
\begin{Usage}
\begin{verbatim}
romicsRestoreMissing(romics_object)
\end{verbatim}
\end{Usage}
%
\begin{Arguments}
\begin{ldescription}
\item[\code{romics\_object}] has to be an romics\_object created using romicsCreateObject()
\end{ldescription}
\end{Arguments}
%
\begin{Details}
This function will renmove any imputed value based on the content of the missingdata layer.
\end{Details}
%
\begin{Value}
The function will return a modified romics\_object that will have NA instead of the imputed data.
\end{Value}
%
\begin{Author}
Geremy Clair
\end{Author}
\HeaderA{romicsSampleNames}{romicsSampleNames()}{romicsSampleNames}
%
\begin{Description}
Indicates the list of sample names from the romics\_object (e.g.,colnames )
\end{Description}
%
\begin{Usage}
\begin{verbatim}
romicsSampleNames(romics_object)
\end{verbatim}
\end{Usage}
%
\begin{Arguments}
\begin{ldescription}
\item[\code{romics\_object}] A romics\_object created using createRomicsObject()
\end{ldescription}
\end{Arguments}
%
\begin{Details}
This function allows to quickly get a vector containing all the factor names present in a romics\_object
\end{Details}
%
\begin{Value}
A character vector containing the list of factor contained in an romics\_object
\end{Value}
%
\begin{Author}
Geremy Clair
\end{Author}
\HeaderA{romicsSd}{romicsSd()}{romicsSd}
%
\begin{Description}
Calculates the standard deviation of each variable within each level of the selected factor and add the generated columns in the statistics layer of the romics\_object.
\end{Description}
%
\begin{Usage}
\begin{verbatim}
romicsSd(romics_object, factor = "main")
\end{verbatim}
\end{Usage}
%
\begin{Arguments}
\begin{ldescription}
\item[\code{romics\_object}] A object created using the function romicsCreateObject().

\item[\code{main\_factor}] Either 'main' OR any factor from the romics\_object, the list of factors from an romics object can be obtained using the function romicsFactorNames().
\end{ldescription}
\end{Arguments}
%
\begin{Details}
Adds the sd columns to the statistics layer. enable to choose a different factor than the main one to do those calculation
\end{Details}
%
\begin{Value}
This function returns a modified romics object, containing sd columns in the statistics layer.
\end{Value}
%
\begin{Author}
Geremy Clair
\end{Author}
\HeaderA{romicsSplitFactor}{romicsSplitFactor}{romicsSplitFactor}
%
\begin{Description}
This function will split different elements contained in a factor (separated by a common separating character string)
\end{Description}
%
\begin{Usage}
\begin{verbatim}
romicsSplitFactor(
  romics_object = romics_object,
  factor = f,
  splitting_string = ";"
)
\end{verbatim}
\end{Usage}
%
\begin{Arguments}
\begin{ldescription}
\item[\code{romics\_object}] A romics\_object created using romicsCreateObject().

\item[\code{factor}] A character vector (lenght=1) containing a factors of the romics\_object. the list of factors can be identified by using the function romicsFactorNames()

\item[\code{splitting\_string}] A character vector (lenght=1) that indicates the separator of the the different factors to be used (when multiple factors to be created are contained in a given factor)
\end{ldescription}
\end{Arguments}
%
\begin{Details}
generate new factor from an existing factor, it will indicates if the new factors are <present> or <absent> from the original factor, if a splitting string is used, then it will create additional factors based on the original factor content.
\end{Details}
%
\begin{Value}
a Romics\_object with new factors
\end{Value}
%
\begin{Author}
Geremy Clair, Naina Beishembieva
\end{Author}
\HeaderA{romicsStatLayer}{romicsStatLayer()}{romicsStatLayer}
%
\begin{Description}
This function will add a statistic layer if it does not exist, it will also verify if the existing stat layer has the right number of rows (same as the data), if not it will create a blank stat layer.
\end{Description}
%
\begin{Usage}
\begin{verbatim}
romicsStatLayer(romics_object = romics_object)
\end{verbatim}
\end{Usage}
%
\begin{Arguments}
\begin{ldescription}
\item[\code{romics\_object}] An romics\_object containing a statistic layer
\end{ldescription}
\end{Arguments}
%
\begin{Details}
Adds the statistics layer.
\end{Details}
%
\begin{Value}
This function returns a modified romics object, containing statistics layer
\end{Value}
%
\begin{Author}
Geremy Clair
\end{Author}
\HeaderA{romicsSteps}{romicsSteps()}{romicsSteps}
%
\begin{Description}
Displays the content steps layer of the romics\_object
\end{Description}
%
\begin{Usage}
\begin{verbatim}
romicsSteps(romics_object, show_dates = TRUE, show_details = TRUE)
\end{verbatim}
\end{Usage}
%
\begin{Arguments}
\begin{ldescription}
\item[\code{romics\_object}] A romics\_object created using createRomicsObject()

\item[\code{show\_dates}] Boolean indicating if the dates have to be displayed

\item[\code{show\_details}] Boolean indicating if the details have to be displayed
\end{ldescription}
\end{Arguments}
%
\begin{Value}
This function will return the steps of an romics\_object
\end{Value}
\HeaderA{romicsSubset}{romicsSubset()}{romicsSubset}
%
\begin{Description}
Keeps or drops a subset of specific elements/columns from the romics\_object

Keeps or drops a subset of specific elements/columns from the romics\_object
\end{Description}
%
\begin{Usage}
\begin{verbatim}
romicsSubset(
  romics_object,
  subset_vector,
  filter_mode = "keep",
  by = "colnames",
  factor = "main",
  handle_embeddings = "remove"
)

romicsSubset(
  romics_object,
  subset_vector,
  filter_mode = "keep",
  by = "colnames",
  factor = "main",
  handle_embeddings = "remove"
)
\end{verbatim}
\end{Usage}
%
\begin{Arguments}
\begin{ldescription}
\item[\code{romics\_object}] A romics\_object created using createRomicsObject()

\item[\code{subset\_vector}] A character vector of factor levels or colnames to keep in the object.

\item[\code{filter\_mode}] Either 'keep' or 'drop' to indicate if you want to conserve or to drop the elements from a given factor

\item[\code{by}] Either 'colnames' or 'level' to indicate what you want to keep or drop.

\item[\code{factor}] A factor contained in the metadata of the romics\_object, to obtain the list of factors please use the function romicsFactorNames()

\item[\code{handle\_embeddings}] Either 'remove' (default) or 'subset'. How to handle embeddings when subsetting.
\end{ldescription}
\end{Arguments}
%
\begin{Details}
This function creates a new object based on a previous romics\_object to include or drop a list of specified columns from the original object. The created object will have a new step object created that will indicate the name of the original object to be subsetted and the log/non-log status of the object.

Note that this function will remove the stat layer from your object

This function creates a new object based on a previous romics\_object to include or drop a list of specified columns from the original object. The created object will have a new step object created that will indicate the name of the original object to be subsetted and the log/non-log status of the object.

Note that this function will remove the stat layer from your object
\end{Details}
%
\begin{Value}
This function generates a subsetted romics\_object

This function generates a subsetted romics\_object
\end{Value}
%
\begin{Author}
Geremy Clair
\end{Author}
\HeaderA{romicsTransferEmbeddings}{romicsTransferEmbeddings()}{romicsTransferEmbeddings}
%
\begin{Description}
This function enables to transfer PCA, UMAP, and t-SNE coordinates/embeddings from one romics\_object
to another one providing added flexibility in embeddings generation. The origin and receiving object
have to contain the same <samples> (data columns).
\end{Description}
%
\begin{Usage}
\begin{verbatim}
romicsTransferEmbeddings(
  origin_romics_object,
  receiving_romics_object,
  methods = c("pca", "umap", "tsne")
)
\end{verbatim}
\end{Usage}
%
\begin{Arguments}
\begin{ldescription}
\item[\code{origin\_romics\_object}] has to be a romics\_object created using romicsCreateObject()

\item[\code{receiving\_romics\_object}] has to be a romics\_object created using romicsCreateObject() and containing the same samples as the origin\_romics\_object.

\item[\code{methods}] a vector indicating which embeddings to transfer. Options are "pca", "umap", and "tsne". Default: c("pca","umap","tsne")
\end{ldescription}
\end{Arguments}
%
\begin{Value}
Returns the receiving\_romics\_object with the transferred embeddings.
\end{Value}
%
\begin{Author}
Geremy Clair
\end{Author}
\HeaderA{romicsTransferFactor}{romicsTransferFactor()}{romicsTransferFactor}
%
\begin{Description}
Transfer one or more metadata factors from one romics\_object to another.
The function ensures that factors are properly aligned based on sample names.
If the source object is a subset, samples present only in the target will be
labeled as "not\_present\_in\_subset".
\end{Description}
%
\begin{Usage}
\begin{verbatim}
romicsTransferFactor(
  source_romics_object,
  target_romics_object,
  factors,
  overwrite = FALSE,
  allow_partial = FALSE,
  missing_label = "not_present_in_subset"
)
\end{verbatim}
\end{Usage}
%
\begin{Arguments}
\begin{ldescription}
\item[\code{source\_romics\_object}] A romics\_object containing the factors to be transferred

\item[\code{target\_romics\_object}] A romics\_object that will receive the transferred factors

\item[\code{factors}] Character vector specifying the names of factors to transfer from source to target

\item[\code{overwrite}] Boolean. If TRUE, existing factors in target will be overwritten. Default: FALSE

\item[\code{allow\_partial}] Boolean. If TRUE, allows transfer even when source is a subset of target. Default: FALSE

\item[\code{missing\_label}] Character string to use for samples not present in source. Default: "not\_present\_in\_subset"
\end{ldescription}
\end{Arguments}
%
\begin{Details}
This function transfers metadata factors between romics\_objects. By default, both objects
must contain the same samples. When allow\_partial = TRUE, the source can be a subset of the target,
and missing samples will be labeled with missing\_label.
\end{Details}
%
\begin{Value}
Returns the target\_romics\_object with the transferred factors added to its metadata
\end{Value}
%
\begin{Author}
Geremy Clair
\end{Author}
\HeaderA{romicsTrend}{romicsTrend()}{romicsTrend}
%
\begin{Description}
This function extract the linear and quadratic trend pvalues and direction based on the pvalue it attributes the best fitted trend
\end{Description}
%
\begin{Usage}
\begin{verbatim}
romicsTrend(
  romics_object,
  factor = "main",
  log_factor = FALSE,
  analysis_type = "both",
  p = 0.05
)
\end{verbatim}
\end{Usage}
%
\begin{Arguments}
\begin{ldescription}
\item[\code{romics\_object}] has to be a romics\_object created using the function romicsCreateObject() (see dedicated documentation) and has to be numeric.

\item[\code{factor}] has to be a factor contained in the romics\_object created, the list of factor can be extracted using the function romicsFactorNames() (see dedicated documentation)

\item[\code{log\_factor}] Boolean, indicate if the factor needs to be log transformed (TRUE), or not (FALSE). By default the factor is not transformed

\item[\code{analysis\_type}] indicates if linear and/or quadratic trends analysis have to be performed, has to be 'both', 'linear',or 'quadratic'.

\item[\code{p}] indicates the minimal fitting p-value to be considered.
\end{ldescription}
\end{Arguments}
%
\begin{Details}
This function allows to calculate the linear and quadratic p value and estimate the best fitting model
\end{Details}
%
\begin{Value}
This function returns a romics\_object with new columns in the stat layer corresponding to the best fitting model, and the linear and/or quadratic p-values
\end{Value}
\HeaderA{romicsTrendHeatmap}{romicsTrendHeatmap()}{romicsTrendHeatmap}
%
\begin{Description}
This function plots a heatmap of all the featuress by trend type
\end{Description}
%
\begin{Usage}
\begin{verbatim}
romicsTrendHeatmap(romics_object, factor = "main", log_factor = FALSE, ...)
\end{verbatim}
\end{Usage}
%
\begin{Arguments}
\begin{ldescription}
\item[\code{romics\_object}] has to be a romics\_object created using the function romicsCreateObject() (see dedicated documentation) and has to be numeric.

\item[\code{factor}] has to be a factor contained in the romics\_object created, the list of factor can be extracted using the function romicsFactorNames() (see dedicated documentation)

\item[\code{log\_factor}] Boolean, indicate if the factor needs to be log transformed (TRUE), or not (FALSE). By default the factor is not transformed

\item[\code{...}] arguments to be passed to heatmap.2() see this function documentation for more details
\end{ldescription}
\end{Arguments}
%
\begin{Details}
this function enables to calculate and extract the linear or quardratic trend p values and the best fitted trend (e.g. linear\_increasing, linear\_decreasing, quadratic\_concave,quadratic\_convex)
\end{Details}
%
\begin{Value}
This function returns a list
\end{Value}
%
\begin{Author}
Geremy Clair and Feng Song
\end{Author}
\HeaderA{romicsTsne}{romicsTsne()}{romicsTsne}
%
\begin{Description}
This function utilizes Rtsne::Rtsne() function to calculate the t-SNE embedding based on the data contained in a Romics\_object and export the embedding in the embeddings layer of the romics\_object.
You can use a subset of features to generate the embedding, all the parameters of the Rtsne::Rtsne() function are available, see the documentation of this function for more details.
The results are recorded in the embeddings layer of the romics\_object, with rows named 'tsne\_component\_1', 'tsne\_component\_2', etc.
If a t-SNE was run previously, the embedding will replace the current embedding.
\end{Description}
%
\begin{Usage}
\begin{verbatim}
romicsTsne(
  romics_object,
  verbose = TRUE,
  perplexity = 30,
  dims = 2,
  initial_dims = 50,
  standardize = TRUE,
  theta = 0.5,
  max_iter = 1000,
  pca = TRUE,
  pca_center = TRUE,
  pca_scale = FALSE,
  check_duplicates = TRUE,
  lock.seed = TRUE,
  seed = 42,
  ...
)
\end{verbatim}
\end{Usage}
%
\begin{Arguments}
\begin{ldescription}
\item[\code{romics\_object}] has to be an romics\_object created using romicsCreateObject() note that the data cannot contain missing values, missing values (e.g., NA) can be imputed using a RomicsProcessor::impute function. After performing the t-SNE the imputation can be cancelled using the romicsRestoreMissing() function.

\item[\code{verbose}] Logical; If TRUE, also assigns the complete t-SNE results to TSNE\_results in the global environment and shows progress

\item[\code{perplexity}] Numeric; Perplexity parameter (should not be bigger than 3 * perplexity < nrow(X) - 1, see details for interpretation). Default: 30

\item[\code{dims}] Integer; Output dimensionality (default: 2)

\item[\code{initial\_dims}] Integer; the number of dimensions that should be retained in the initial PCA step (default: 50)

\item[\code{standardize}] Logical; Whether to scale and center the data before applying t-SNE. Default: TRUE

\item[\code{theta}] Numeric; Speed/accuracy trade-off (increase for less accuracy), set to 0.0 for exact t-SNE (default: 0.5)

\item[\code{max\_iter}] Integer; Number of iterations (default: 1000)

\item[\code{pca}] Logical; Whether an initial PCA step should be performed (default: TRUE)

\item[\code{pca\_center}] Logical; Should data be centered before pca is applied? (default: TRUE)

\item[\code{pca\_scale}] Logical; Should data be scaled before pca is applied? (default: FALSE)

\item[\code{check\_duplicates}] Logical; Checks whether duplicates are present. It is best to make sure there are no duplicates present and set this option to FALSE, especially for large datasets (default: TRUE)

\item[\code{lock.seed}] Logical; TRUE or FALSE to indicate if the seed is locked to enable reproducing results

\item[\code{seed}] Numeric; value indicating what seed to use when random seed are used

\item[\code{...}] Arguments passed to Rtsne function.
\end{ldescription}
\end{Arguments}
%
\begin{Details}
This function calculates t-SNE embeddings and stores them in the romics\_object. It can also store the complete
t-SNE results in the global environment if verbose=TRUE. t-SNE is particularly good at preserving local structure and revealing clusters.
\end{Details}
%
\begin{Value}
Returns a romics\_object with t-SNE embeddings added to the embeddings layer.
\end{Value}
%
\begin{Author}
Geremy Clair
\end{Author}
\HeaderA{romicsTsnePlot}{romicsTsnePlot()}{romicsTsnePlot}
%
\begin{Description}
Plots the samples on a t-SNE plot in 2D or 3D.
The plot can be colored by any factor in the romics\_object's metadata or by feature values.
This function uses the t-SNE embeddings stored in the romics\_object.
\end{Description}
%
\begin{Usage}
\begin{verbatim}
romicsTsnePlot(
  romics_object,
  Xcomp = 1,
  Ycomp = 2,
  Zcomp = NULL,
  label = FALSE,
  feature = NULL,
  factor_name = NULL,
  color_palette = NULL,
  standardize = TRUE,
  size = 3,
  alpha = 0.8,
  plotly = FALSE,
  show_outline = TRUE
)
\end{verbatim}
\end{Usage}
%
\begin{Arguments}
\begin{ldescription}
\item[\code{romics\_object}] A romics\_object that has t-SNE embeddings in its embeddings layer (generated using romicsTsne)

\item[\code{Xcomp}] numerical/double. Indicate the component to plot on the X axis (default: 1)

\item[\code{Ycomp}] numerical/double. Indicate the component to plot on the Y axis (default: 2)

\item[\code{Zcomp}] numerical/double. If provided, creates a 3D plot with this component on the Z axis (default: NULL)

\item[\code{label}] boolean. Indicate if the sample name labels should be plotted (default: FALSE, applies only to 2D plots)

\item[\code{feature}] character string. If provided, the plot will be colored according to the values of this feature.
If NULL (default), coloring will be based on the factor specified by factor\_name.

\item[\code{factor\_name}] character string. The name of the factor in romics\_object\$metadata to use for coloring the points
when feature is NULL. If NULL (default), the romics\_object's main\_factor will be used.

\item[\code{color\_palette}] a color palette to use. For categorical factors, should be a vector with one color per level.
For continuous features (when feature is specified), should be a gradient color palette.
Default: NULL (uses colors\_romics for factors, viridis for features).

\item[\code{standardize}] boolean. Indicates if the feature values should be scaled when coloring by feature (default: TRUE)

\item[\code{size}] numeric. Size of the points in the plot (default: 3)

\item[\code{alpha}] numeric. Transparency of the points in the plot (default: 0.8)

\item[\code{plotly}] boolean. If TRUE, generates an interactive plotly plot instead of a static ggplot (default: FALSE)

\item[\code{show\_outline}] boolean. If FALSE, removes point outlines and shows only fill colors (default: TRUE)
\end{ldescription}
\end{Arguments}
%
\begin{Details}
This function plots t-SNE results using the embeddings stored in the romics\_object.
If Zcomp is provided, a 3D interactive plot is generated using plotly.
If Zcomp is NULL, a 2D plot is generated using ggplot2.
When using a non-main factor for coloring, the function automatically updates the romics\_object
using romicsChangeFactor() to ensure consistent colors across plots.
Coloring is determined automatically:
- If feature is provided, the plot is colored by feature values
- If feature is NULL, the plot is colored by factor values
When show\_outline is FALSE, points are displayed without borders for cleaner appearance.
\end{Details}
%
\begin{Value}
Returns either a ggplot2 plot or a plotly interactive plot.
\end{Value}
%
\begin{Author}
Geremy Clair
\end{Author}
\HeaderA{romicsUmap}{romicsUmap()}{romicsUmap}
%
\begin{Description}
This function utilizes uwot::umap() function to calculate the umap embedding based on the data contained in a Romics\_object and export the embedding in the metadata layer of the romics\_object.
You can use a subset of features to generate the embedding, all the parameters of the uwot::umap() function are available, see the documentation of this function for more details.
The results are recorded in the embeddings layer of the romics\_object, with rows named 'umap\_component\_1', 'umap\_component\_2', etc.
If a umap was run previously, the embedding will replace the current embedding.
\end{Description}
%
\begin{Usage}
\begin{verbatim}
romicsUmap(
  romics_object,
  verbose = TRUE,
  n_neighbors = 30,
  n_components = 2,
  metric = "euclidean",
  standardize = TRUE,
  init,
  use_pca_embeddings = FALSE,
  pca_ncp,
  lock.seed = TRUE,
  seed = 42,
  ...
)
\end{verbatim}
\end{Usage}
%
\begin{Arguments}
\begin{ldescription}
\item[\code{romics\_object}] has to be a romics\_object created using romicsCreateObject() note that the data cannot contain missing values, missing values (e.g., NA) can be imputed using a RomicsProcessor::impute function. After performing the umap the imputation can be cancelled using the romicsRestoreMissing() function.

\item[\code{verbose}] boolean. If TRUE, also assigns the complete UMAP model to UMAP\_results in the global environment

\item[\code{n\_neighbors}] The size of local neighborhood (in terms of number of neighboring sample points) used for manifold approximation. Larger values result in more global views of the manifold, while smaller values result in more local data being preserved. In general values should be in the range 2 to 100.

\item[\code{n\_components}] The dimension of the space to embed into. This defaults to 2 to provide easy visualization, but can reasonably be set to any integer value in the range 2 to 100.

\item[\code{metric}] Type of distance metric to use to find nearest neighbors. For nn\_method = "annoy" this can be one of: "euclidean" (the default), "cosine", "manhattan", "hamming","correlation" (a distance based on the Pearson correlation), "categorical" (see below).

\item[\code{standardize}] Scaling to apply to the data layer of the romics\_object: "none" or FALSE or NULL No scaling; "Z" or "scale" or TRUE Scale each feature to zero mean and variance 1; "maxabs" Center each feature to mean 0, then divide each element by the maximum absolute value over the entire matrix; "range" Range scale the entire matrix, so the smallest element is 0 and the largest is 1; "colrange" Scale each feature in the range (0,1); For UMAP, the default is "none".

\item[\code{init}] Type of initialization for the coordinates. Options are:
\begin{itemize}

\item{} "spectral" Spectral embedding using the normalized Laplacian of the fuzzy 1-skeleton, with Gaussian noise added
\item{} "normlaplacian" Spectral embedding using the normalized Laplacian of the fuzzy 1-skeleton, without noise
\item{} "random" Coordinates assigned using a uniform random distribution between -10 and 10
\item{} "lvrandom" Coordinates assigned using a Gaussian distribution with standard deviation 1e-4, as used in LargeVis (Tang et al., 2016) and t-SNE
\item{} "laplacian" Spectral embedding using the Laplacian Eigenmap (Belkin and Niyogi, 2002)
\item{} "pca" The first two principal components from PCA of X if X is a data frame, and from a 2-dimensional classical MDS if X is of class "dist"
\item{} "spca" Like "pca", but each dimension is then scaled so the standard deviation is 1e-4, to give a distribution similar to that used in t-SNE. This is an alias for init = "pca", init\_sdev = 1e-4
\item{} "agspectral" An "approximate global" modification of "spectral" which all edges in the graph to a value of 1, and then sets a random number of edges (negative\_sample\_rate edges per vertex) to 0.1, to approximate the effect of non-local affinities
\item{} "pca\_embeddings" Use existing PCA embeddings from the romics\_object embeddings layer (see use\_pca\_embeddings parameter)
\item{} A matrix of initial coordinates

\end{itemize}

For spectral initializations, ("spectral", "normlaplacian", "laplacian", "agspectral"), if more than one connected component is identified, no spectral initialization is attempted. Instead a PCA-based initialization is attempted. If verbose = TRUE the number of connected components are logged to the console. The existence of multiple connected components implies that a global view of the data cannot be attained with this initialization. Increasing the value of n\_neighbors may help.

\item[\code{use\_pca\_embeddings}] Logical. If TRUE and init is not specified, will attempt to use PCA embeddings from the romics\_object embeddings layer for initialization.
If TRUE and no PCA embeddings are found, will run romicsPCA() first. Default is FALSE.
When TRUE, the number of PCA components used will match n\_components.

\item[\code{pca\_ncp}] Integer. Number of principal components to calculate if PCA needs to be run for initialization.
Only used when use\_pca\_embeddings = TRUE and PCA embeddings don't exist. Default is max(10, n\_components).

\item[\code{lock.seed}] has to be TRUE or FALSE to indicate if the seed is locked to enable reproducing plotting

\item[\code{seed}] numeric value indicating what seed to use when random seed are used

\item[\code{...}] Arguments passed to umap function.
\end{ldescription}
\end{Arguments}
%
\begin{Details}
This function calculates UMAP embeddings and stores them in the romics\_object. It can also store the complete
UMAP model in the global environment if verbose=TRUE.
\end{Details}
%
\begin{Value}
Returns a romics\_object with UMAP embeddings added to the embeddings layer.
\end{Value}
%
\begin{Author}
Geremy Clair
\end{Author}
\HeaderA{romicsUmapDensityPlot}{romicsUmapDensityPlot()}{romicsUmapDensityPlot}
%
\begin{Description}
Creates a UMAP plot with density overlay showing regions of high point concentration.
This is useful for identifying clusters or areas with many samples in UMAP space.
\end{Description}
%
\begin{Usage}
\begin{verbatim}
romicsUmapDensityPlot(
  romics_object,
  Xcomp = 1,
  Ycomp = 2,
  density_type = "filled",
  bins = 30,
  n_contours = 10,
  color_palette = NULL,
  bandwidth = NULL,
  internal_margin = 0.5
)
\end{verbatim}
\end{Usage}
%
\begin{Arguments}
\begin{ldescription}
\item[\code{romics\_object}] A romics\_object that has UMAP embeddings in its embeddings layer

\item[\code{Xcomp}] numerical/double. Component to plot on the X axis (default: 1)

\item[\code{Ycomp}] numerical/double. Component to plot on the Y axis (default: 2)

\item[\code{density\_type}] character. Type of density visualization: "contour", "filled", "hex", or "2d" (default: "filled")

\item[\code{bins}] numeric. Number of bins for hexagonal binning (used when density\_type = "hex", default: 30)

\item[\code{n\_contours}] numeric. Number of contour levels (used for contour plots, default: 10)

\item[\code{color\_palette}] character vector. Color palette for density gradient (default: viridis)

\item[\code{bandwidth}] numeric. Bandwidth for density estimation (default: NULL for automatic)

\item[\code{internal\_margin}] numeric. Absolute value to extend axis limits beyond data range (default: 0.5)
\end{ldescription}
\end{Arguments}
%
\begin{Value}
Returns a ggplot2 object showing density distribution in UMAP space
\end{Value}
%
\begin{Author}
Geremy Clair
\end{Author}
\HeaderA{romicsUmapPlot}{romicsUmapPlot()}{romicsUmapPlot}
%
\begin{Description}
Plots the samples on a UMAP plot in 2D or 3D.
The plot can be colored by any factor in the romics\_object's metadata or by feature values.
This function uses the UMAP embeddings stored in the romics\_object.
\end{Description}
%
\begin{Usage}
\begin{verbatim}
romicsUmapPlot(
  romics_object,
  Xcomp = 1,
  Ycomp = 2,
  Zcomp = NULL,
  label = FALSE,
  feature = NULL,
  factor_name = NULL,
  color_palette = NULL,
  standardize = TRUE,
  size = 3,
  alpha = 0.8,
  plotly = FALSE,
  show_outline = TRUE
)
\end{verbatim}
\end{Usage}
%
\begin{Arguments}
\begin{ldescription}
\item[\code{romics\_object}] A romics\_object that has UMAP embeddings in its embeddings layer (generated using romicsUmap)

\item[\code{Xcomp}] numerical/double. Indicate the component to plot on the X axis (default: 1)

\item[\code{Ycomp}] numerical/double. Indicate the component to plot on the Y axis (default: 2)

\item[\code{Zcomp}] numerical/double. If provided, creates a 3D plot with this component on the Z axis (default: NULL)

\item[\code{label}] boolean. Indicate if the sample name labels should be plotted (default: FALSE, applies only to 2D plots)

\item[\code{feature}] character string. If provided, the plot will be colored according to the values of this feature.
If NULL (default), coloring will be based on the factor specified by factor\_name.

\item[\code{factor\_name}] character string. The name of the factor in romics\_object\$metadata to use for coloring the points
when feature is NULL. If NULL (default), the romics\_object's main\_factor will be used.

\item[\code{color\_palette}] a color palette to use. For categorical factors, should be a vector with one color per level.
For continuous features (when feature is specified), should be a gradient color palette.
Default: NULL (uses colors\_romics for factors, viridis for features).

\item[\code{standardize}] boolean. Indicates if the feature values should be scaled when coloring by feature (default: TRUE)

\item[\code{size}] numeric. Size of the points in the plot (default: 3)

\item[\code{alpha}] numeric. Transparency of the points in the plot (default: 0.8)

\item[\code{plotly}] boolean. If TRUE, generates an interactive plotly plot instead of a static ggplot (default: FALSE)

\item[\code{show\_outline}] boolean. If FALSE, removes point outlines and shows only fill colors (default: TRUE)
\end{ldescription}
\end{Arguments}
%
\begin{Details}
This function plots UMAP results using the embeddings stored in the romics\_object.
If Zcomp is provided, a 3D interactive plot is generated using plotly.
If Zcomp is NULL, a 2D plot is generated using ggplot2.

Color palette behavior:
- For categorical factors: If color\_palette is provided and has enough colors, it will be used.
Otherwise, colors\_romics or ROP\_colors will be used (with warning if insufficient).
- For continuous features: If color\_palette is NULL, viridis is used. Otherwise, the provided palette is used.

When using a non-main factor for coloring without a custom color\_palette, the function automatically
updates the romics\_object using romicsChangeFactor() to ensure consistent colors across plots.
\end{Details}
%
\begin{Value}
Returns either a ggplot2 plot or a plotly interactive plot.
\end{Value}
%
\begin{Author}
Geremy Clair
\end{Author}
\HeaderA{romicsUpdateColor}{romicsUpdateColor()}{romicsUpdateColor}
%
\begin{Description}
Updates the colors layer contained in the romics\_object
\end{Description}
%
\begin{Usage}
\begin{verbatim}
romicsUpdateColor(romics_object)
\end{verbatim}
\end{Usage}
%
\begin{Arguments}
\begin{ldescription}
\item[\code{romics\_object}] A romics\_object created using createRomicsObject()
\end{ldescription}
\end{Arguments}
%
\begin{Value}
This function returns a romics\_object with updated colors.
\end{Value}
%
\begin{Author}
Geremy Clair
\end{Author}
\HeaderA{romicsUpdateSteps}{romicsUpdateSteps()}{romicsUpdateSteps}
%
\begin{Description}
Updates the steps of the romics\_object, require to have recorded the argument in earlier steps of the function
\end{Description}
%
\begin{Usage}
\begin{verbatim}
romicsUpdateSteps(romics_object, arguments)
\end{verbatim}
\end{Usage}
%
\begin{Arguments}
\begin{ldescription}
\item[\code{romics\_object}] A romics\_object created using createRomicsObject()

\item[\code{arguments}] the arguments of a function are required to read the user input of a function, this user input will be used to generate the steps, the arguments are obtained by running the following code <arguments<-as.list(match.call())> in the first line of a function
\end{ldescription}
\end{Arguments}
%
\begin{Details}
The goal of Romics processor is to provide a trackable and reproducible pipeline for processing omics data. Subsequently it is necessary when a function is created to implement a way to record the user input that will be recorded in the steps layer of the Romics\_object.

This function will enable to simplify the work of developers who want to contribute to Romics by simplifying this process. Only two lines of codes are then necessary to update the steps.

The first line of code has to be placed in the first line after the function declaration : <arguments<-as.list(match.call())>

The second line of code has to be <romics\_object<-stepUpdater(romics\_object,arguments)> placed at the end of the function code (ideally right before returning the processed romics\_object or graphic generated by the function)
\end{Details}
%
\begin{Value}
This function add the description of the processing to the step layer of an Romics object
\end{Value}
%
\begin{Author}
Geremy Clair
\end{Author}
\HeaderA{romicsValidFeatures}{romicsValidFeatures()}{romicsValidFeatures}
%
\begin{Description}
Calculates and returns the number of features with valid (non-missing) values for each sample in the romics\_object.
\end{Description}
%
\begin{Usage}
\begin{verbatim}
romicsValidFeatures(
  romics_object,
  return_type = "table",
  custom_colors = "colorlist"
)
\end{verbatim}
\end{Usage}
%
\begin{Arguments}
\begin{ldescription}
\item[\code{romics\_object}] has to be a romics\_object created using romicsCreateObject()

\item[\code{return\_type}] Character string specifying the output format. Options are: "table" (default) returns a data frame with sample names and counts; "vector" returns a named numeric vector; "plot" returns a ggplot2 bar plot.

\item[\code{custom\_colors}] Character string or vector. If "colorlist" (default), uses the colors from romics\_object. If a vector of colors is provided, it will be used for the plot (only applicable when return\_type = "plot").
\end{ldescription}
\end{Arguments}
%
\begin{Details}
This function counts the number of non-missing values (valid features) for each sample (column) in the romics\_object\$data layer. The results can be returned as a table, vector, or plot depending on the return\_type parameter.
\end{Details}
%
\begin{Value}
Depending on return\_type:
\begin{itemize}

\item{} "table": A data frame with columns: Sample (sample names), Valid\_features (count), Percent\_valid (percentage)
\item{} "vector": A named numeric vector with sample names as names and counts as values
\item{} "plot": A ggplot2 bar plot showing valid features per sample, colored by the main factor

\end{itemize}

\end{Value}
%
\begin{Author}
Geremy Clair
\end{Author}
%
\begin{Examples}
\begin{ExampleCode}
## Not run: 
# Get a table of valid features per sample
valid_table <- romicsValidFeatures(romics_object)

# Get a vector of valid features per sample
valid_vector <- romicsValidFeatures(romics_object, return_type = "vector")

# Plot the valid features per sample
romicsValidFeatures(romics_object, return_type = "plot")

## End(Not run)
\end{ExampleCode}
\end{Examples}
\HeaderA{romicsValidFeaturesPercent}{romicsValidFeaturesPercent()}{romicsValidFeaturesPercent}
%
\begin{Usage}
\begin{verbatim}
romicsValidFeaturesPercent(
  romics_object,
  return_type = "table",
  reorder = NULL,
  features = NULL
)
\end{verbatim}
\end{Usage}
%
\begin{Arguments}
\begin{ldescription}
\item[\code{romics\_object}] has to be a romics\_object created using romicsCreateObject()

\item[\code{return\_type}] Character string specifying the output format. Options are: "table" (default) returns a data frame with feature names and percentages; "vector" returns a named numeric vector; "plot" returns a ggplot2 bar plot.

\item[\code{reorder}] Character string indicating how to order features in the output. Options are: NULL (original order), "ascending" (lowest to highest valid 

\bsl{}itemfeaturesCharacter vector containing a list of specific features to analyze. If NULL (default), all features are included.

\end{ldescription}

Calculates and returns the percentage of samples that have valid (non-missing) values for each feature in the romics\_object.


This function calculates the percentage of samples (columns) that have valid (non-missing) values for each feature (row) in the romics\_object\$data layer. This is complementary to romicsValidFeatures() which counts valid features per sample.


Geremy Clair

\end{Arguments}
\HeaderA{romicsVolcano}{romicsVolcano()}{romicsVolcano}
%
\begin{Description}
generate one or multiple volcano plots from t.test, wilcox.test, or LMM run using the functions romicsTtest(), romicsWilcoxTest(), or romicsLMM() respectively
\end{Description}
%
\begin{Usage}
\begin{verbatim}
romicsVolcano(
  romics_object,
  p_type = "p",
  p = 0.05,
  min_fold_change = 0.6,
  colors = c("#2cbcb2", "#242021", "#d44e28"),
  stat_type = "auto",
  plot = "all",
  plotly = FALSE,
  label_features = FALSE,
  top_features = 10,
  top_features_by = "abundance",
  xlim = NULL,
  ylim = NULL,
  size = 2,
  alpha = 0.5
)
\end{verbatim}
\end{Usage}
%
\begin{Arguments}
\begin{ldescription}
\item[\code{p\_type}] 'p' or 'padj' to indicate the type of pvalue to consider for the volcano plots.

\item[\code{p}] numeric value to indicate the maximum pvalue to use for the coloring of the volcano plot

\item[\code{colors}] numeric vector of lenght 3 used to color the features lower, not significantly changing and higher in the considered comparisons.

\item[\code{stat\_type}] 't.test', 'wilcox.test', 'LMM' or 'auto' to indicate what statistics to use for the volcano plots generation. If 'auto', will detect available tests.

\item[\code{plot}] either 'all' if all paired comparisons have to be displayed OR a vector of numeric values comprised between 1 and the maximum number of possible plots generated.

\item[\code{plotly}] logical (TRUE or FALSE) to indicate whether to return interactive plotly plots (TRUE) or static ggplot plots (FALSE).

\item[\code{label\_features}] logical (TRUE or FALSE) to indicate whether to label top features on the plot.

\item[\code{top\_features}] numeric value indicating the number of top upregulated and downregulated features to label (only those passing the filters).

\item[\code{top\_features\_by}] 'abundance' or 'p' to indicate the criteria for selecting top features. 'abundance' selects by fold change magnitude, 'p' selects by p-value significance.

\item[\code{xlim}] numeric vector of length 2 specifying the x-axis limits (e.g., c(-2, 2)). If NULL, limits are determined automatically.

\item[\code{ylim}] numeric vector of length 2 specifying the y-axis limits (e.g., c(0, 10)). If NULL, limits are determined automatically.

\item[\code{size}] numeric value for point size in the plots. Default is 2.

\item[\code{alpha}] numeric value for point opacity in the plots (0-1). Default is 0.5.
\end{ldescription}
\end{Arguments}
%
\begin{Details}
generate one or multiple volcano plots from t.test, wilcox.test, or LMM run using the functions romicsTtest(), romicsWilcoxTest(), or romicsLMM() respectively. Automatically detects whether data is clustered or not. If xlim or ylim are specified, a warning is issued if any data points fall outside these limits.
\end{Details}
%
\begin{Value}
This function will print different plots requested
\end{Value}
%
\begin{Author}
Geremy Clair
\end{Author}
\HeaderA{romicsVolcanoByCluster}{romicsVolcanoByCluster()}{romicsVolcanoByCluster}
%
\begin{Description}
generate volcano plots for all comparisons within a specific cluster
\end{Description}
%
\begin{Usage}
\begin{verbatim}
romicsVolcanoByCluster(
  romics_object,
  within,
  multipanel = TRUE,
  stat_type = "auto",
  p_type = "p",
  p = 0.05,
  min_fold_change = 0.6,
  colors = c("#2cbcb2", "#242021", "#d44e28"),
  plotly = FALSE,
  label_features = FALSE,
  top_features = 10,
  top_features_by = "abundance",
  n_cols = NULL,
  return_plot = FALSE,
  xlim = NULL,
  ylim = NULL,
  size = 2,
  alpha = 0.5
)
\end{verbatim}
\end{Usage}
%
\begin{Arguments}
\begin{ldescription}
\item[\code{romics\_object}] a romics object containing statistics

\item[\code{within}] character string specifying the cluster to filter (e.g., "Leiden\_clust\_01")

\item[\code{multipanel}] logical (TRUE or FALSE) to indicate whether to display all plots in a panel (TRUE) or one by one (FALSE)

\item[\code{stat\_type}] 't.test', 'wilcox.test', 'LMM' or 'auto' to indicate what statistics to use. If 'auto', will detect available tests.

\item[\code{p\_type}] 'p' or 'padj' to indicate the type of pvalue to consider for the volcano plots

\item[\code{p}] numeric value to indicate the maximum pvalue to use for the coloring of the volcano plot

\item[\code{min\_fold\_change}] numeric value indicating the minimum fold change threshold

\item[\code{colors}] character vector of length 3 used to color the features (down, non-significant, up)

\item[\code{plotly}] logical (TRUE or FALSE) to indicate whether to return interactive plotly plots (TRUE) or static ggplot plots (FALSE)

\item[\code{label\_features}] logical (TRUE or FALSE) to indicate whether to label top features on the plot

\item[\code{top\_features}] numeric value indicating the number of top upregulated and downregulated features to label

\item[\code{top\_features\_by}] 'abundance' or 'p' to indicate the criteria for selecting top features

\item[\code{n\_cols}] numeric value indicating the number of columns in the multipanel layout (NULL for auto square-ish grid)

\item[\code{return\_plot}] logical (TRUE or FALSE) to indicate whether to return the plot object instead of printing it (default: FALSE)

\item[\code{xlim}] numeric vector of length 2 specifying the x-axis limits (e.g., c(-2, 2)). If NULL, limits are determined automatically.

\item[\code{ylim}] numeric vector of length 2 specifying the y-axis limits (e.g., c(0, 10)). If NULL, limits are determined automatically.

\item[\code{size}] numeric value for point size in the plots. Default is 2.

\item[\code{alpha}] numeric value for point opacity in the plots (0-1). Default is 0.5.
\end{ldescription}
\end{Arguments}
%
\begin{Details}
generates volcano plots for all comparisons within a specific cluster from any statistical test (t.test, wilcox.test, or LMM). If xlim or ylim are specified, a warning is issued if any data points fall outside these limits.
\end{Details}
%
\begin{Value}
Returns plot list invisibly if return\_plot=FALSE, or returns the combined plot object if return\_plot=TRUE
\end{Value}
%
\begin{Author}
Geremy Clair
\end{Author}
\HeaderA{romicsZeroToMissing}{romicsZeroToMissing()}{romicsZeroToMissing}
%
\begin{Description}
Replaces zeros in the romics\_object by NA values
\end{Description}
%
\begin{Usage}
\begin{verbatim}
romicsZeroToMissing(romics_object)
\end{verbatim}
\end{Usage}
%
\begin{Arguments}
\begin{ldescription}
\item[\code{romics\_object}] has to be an romics\_object created using romicsCreateObject()
\end{ldescription}
\end{Arguments}
%
\begin{Details}
This function will convert 0 values to NA in the data and missingdata layers
\end{Details}
%
\begin{Value}
This function returns the transformed romics\_object with updated data and missingdata layers
\end{Value}
%
\begin{Author}
Geremy Clair
\end{Author}
\HeaderA{romicsZscores}{romicsZscores()}{romicsZscores}
%
\begin{Description}
Calculates the Zscores for each cell of the data layer of the romics\_object. the generated results are added to the statistics Layer as new columns.
\end{Description}
%
\begin{Usage}
\begin{verbatim}
romicsZscores(romics_object)
\end{verbatim}
\end{Usage}
%
\begin{Arguments}
\begin{ldescription}
\item[\code{romics\_object}] A romics\_object created with the function romicsCreateObject(),
\end{ldescription}
\end{Arguments}
%
\begin{Details}
adds the Zscores columns to the statistics Layer of a romics\_object
\end{Details}
%
\begin{Value}
an romics\_object with the statistical layer containing the newly generated Zscores columns
\end{Value}
%
\begin{Author}
Geremy Clair
\end{Author}
\HeaderA{SampleCorrelation}{SampleCorrelation()}{SampleCorrelation}
%
\begin{Description}
calculate the correlation matrix for the samples of a romics\_object
\end{Description}
%
\begin{Usage}
\begin{verbatim}
SampleCorrelation(romics_object, corr_type = "pearson", use = "everything")
\end{verbatim}
\end{Usage}
%
\begin{Arguments}
\begin{ldescription}
\item[\code{romics\_object}] an object of type romics\_object

\item[\code{corr\_type}] has to be either 'pearson', 'kendall',or 'spearman', indicates the type of correlation calculated.

\item[\code{use}] has to be either 'everything', 'all.obs', 'complete.obs', 'na.or.complete', or 'pairwise.complete.obs'
\end{ldescription}
\end{Arguments}
%
\begin{Details}
calculate a correlation matrix for the samples of an Romics\_object
\end{Details}
%
\begin{Value}
This function will return a correlation matrix for the samples of an romics\_object
\end{Value}
%
\begin{Author}
Geremy Clair
\end{Author}
\HeaderA{SampleCorrelationHclust}{SampleCorrelationHclust()}{SampleCorrelationHclust}
%
\begin{Description}
calculate the sample correlation matrix for an romics\_object and generate a hclust object.
\end{Description}
%
\begin{Usage}
\begin{verbatim}
SampleCorrelationHclust(
  romics_object,
  corr_type = "pearson",
  use = "everything",
  hclust_method = "ward.D",
  dist_method = "euclidean"
)
\end{verbatim}
\end{Usage}
%
\begin{Arguments}
\begin{ldescription}
\item[\code{romics\_object}] an object of type romics\_object

\item[\code{corr\_type}] has to be either 'pearson', 'kendall',or 'spearman', indicates the type of correlation calculated.

\item[\code{use}] has to be either 'everything', 'all.obs', 'complete.obs', 'na.or.complete', or 'pairwise.complete.obs'

\item[\code{hclust\_method}] has to be in 'correlation.dissimilarity','ward.D', 'ward.D2', 'single', 'complete', 'average', 'mcquitty', 'median', or 'centroid'

\item[\code{dist\_method}] has to be in 'correlation.dissimilarity','euclidean', 'maximum', 'manhattan', 'canberra', 'binary' or 'minkowski'.
\end{ldescription}
\end{Arguments}
%
\begin{Details}
calculate the sample correlation matrix for the samples of an romics object and generate a hclust object.
\end{Details}
%
\begin{Value}
This function will return a hclust result that can be used to generate hclust plots.
\end{Value}
%
\begin{Author}
Geremy Clair
\end{Author}
\HeaderA{SampleCorrelationHclustPlot}{SampleCorrelationHclustPlot()}{SampleCorrelationHclustPlot}
%
\begin{Description}
calculate the Sample correlation matrix for the samples of an romics\_object and generate a hclust plot.
\end{Description}
%
\begin{Usage}
\begin{verbatim}
SampleCorrelationHclustPlot(
  romics_object,
  corr_type = "pearson",
  use = "everything",
  hclust_method = "ward.D",
  dist_method = "euclidean",
  plot_type = "dendrogram",
  number_of_clusters = 0,
  cex = 1,
  cluster_colors = c("c1", "c2", "c3", "c4")
)
\end{verbatim}
\end{Usage}
%
\begin{Arguments}
\begin{ldescription}
\item[\code{romics\_object}] an object of type romics\_object

\item[\code{corr\_type}] has to be either 'pearson', 'kendall',or 'spearman', indicates the type of correlation calculated.

\item[\code{use}] has to be either 'everything', 'all.obs', 'complete.obs', 'na.or.complete', or 'pairwise.complete.obs'

\item[\code{hclust\_method}] has to be in 'ward.D', 'ward.D2', 'single', 'complete', 'average', 'mcquitty', 'median', or 'centroid'

\item[\code{dist\_method}] has to be in 'correlation.dissimilarity','euclidean', 'maximum', 'manhattan', 'canberra', 'binary' or 'minkowski'.

\item[\code{plot\_type}] has to be either 'dendrogram', 'unrooted', or 'fan'

\item[\code{number\_of\_clusters}] has to be a numeric value to indicate the number of clusters

\item[\code{cex}] has to be a numeric value to indicate the label font size

\item[\code{cluster\_colors}] has to be a character vector containing the list of colors used for the clusters.
\end{ldescription}
\end{Arguments}
%
\begin{Details}
calculate the Sample correlation matrix for the samples of an romics object and generate a hclust plot
\end{Details}
%
\begin{Value}
This function will return a hclust plot generated with the base R plot() function.
\end{Value}
%
\begin{Author}
Geremy Clair
\end{Author}
\HeaderA{SampleCorrelationHeatmap}{SampleCorrelationHeatmap()}{SampleCorrelationHeatmap}
%
\begin{Description}
calculate the sample correlation matrix of an romics\_object and hclust the results are displayed as an heatmap.
\end{Description}
%
\begin{Usage}
\begin{verbatim}
SampleCorrelationHeatmap(
  romics_object,
  corr_type = "pearson",
  use = "everything",
  number_of_clusters = 0,
  cex = 1,
  cluster_colors = c("c1", "c2", "c3", "c4"),
  heatmap_col = viridis(20),
  dist_method = "euclidean",
  hclust_method = "ward.D",
  cellnote = F
)
\end{verbatim}
\end{Usage}
%
\begin{Arguments}
\begin{ldescription}
\item[\code{romics\_object}] an object of type romics\_object

\item[\code{corr\_type}] has to be either 'pearson', 'kendall',or 'spearman', indicates the type of correlation calculated.

\item[\code{use}] has to be either 'everything', 'all.obs', 'complete.obs', 'na.or.complete', or 'pairwise.complete.obs'

\item[\code{number\_of\_clusters}] has to be a numeric value to indicate the number of clusters

\item[\code{cex}] has to be a numeric value to indicate the label font size

\item[\code{cluster\_colors}] has to be a character vector containing the list of colors used for the clusters.

\item[\code{heatmap\_col}] has to be a color gradient (e.g.viridis(20), greenred(50))

\item[\code{dist\_method}] has to be in 'correlation.dissimilarity','euclidean', 'maximum', 'manhattan', 'canberra', 'binary' or 'minkowski'.

\item[\code{hclust\_method}] has to be in 'ward.D', 'ward.D2', 'single', 'complete', 'average', 'mcquitty', 'median', or 'centroid'

\item[\code{cellnote}] has to be TRUE or FALSE to indicate if the correlation numbers should be displayed onto the heatmap.
\end{ldescription}
\end{Arguments}
%
\begin{Value}
This function will return a an heatmap generated with heatmap.2().
\end{Value}
%
\begin{Author}
Geremy Clair
\end{Author}
\HeaderA{singleFeaturePlot}{singleFeaturePlot()}{singleFeaturePlot}
%
\begin{Description}
Creates a single feature plot from a romics\_object with optional statistical comparison brackets. Supports multiple plot types including jitter, boxplot, violin, and combinations.
\end{Description}
%
\begin{Usage}
\begin{verbatim}
singleFeaturePlot(
  romics_object,
  feature = "feature",
  plot_type = "jb",
  factor = "main",
  ylim = NULL,
  title = "auto",
  pval = "none",
  y_bracket_pos = 0.5,
  test_priority = c("Ttest", "ttest", "wilcoxon", "glm_binomial", "ANOVA"),
  significance_threshold = 0.05
)
\end{verbatim}
\end{Usage}
%
\begin{Arguments}
\begin{ldescription}
\item[\code{romics\_object}] A romics\_object created using romicsCreateObject().

\item[\code{feature}] Character. Name of the feature to plot. Must match exactly (case-sensitive). Default: "feature"

\item[\code{plot\_type}] Character. Plot type: "jitter", "boxplot", "violin", "jb" (jitter+boxplot), or "jv" (jitter+violin). Default: "jb"

\item[\code{factor}] Character. Factor name from romics\_object to use for grouping. Use "main" for the main factor. Default: "main"

\item[\code{ylim}] Numeric vector of length 2. Y-axis limits c(min, max). Default: NULL (auto-scaling)

\item[\code{title}] Character. Plot title. Use "auto" for feature name or provide custom title. Default: "auto"

\item[\code{pval}] Character. P-value type for statistical brackets: "none", "p", or "padj". Default: "none"

\item[\code{y\_bracket\_pos}] Numeric. Vertical spacing between statistical brackets. Default: 0.5

\item[\code{test\_priority}] Character vector. Priority order for statistical tests when multiple are available. Default: c("Ttest", "ttest", "wilcoxon", "glm\_binomial", "ANOVA")

\item[\code{significance\_threshold}] Numeric. P-value threshold for highlighting significant results in red. Default: 0.05
\end{ldescription}
\end{Arguments}
%
\begin{Details}
The function automatically detects available statistical comparisons in the romics\_object's statistics layer
and adds brackets showing p-values. When multiple statistical tests are available for the same comparison,
the function uses the test\_priority order to select which test to display.

Significant p-values (below significance\_threshold) are displayed in red, while non-significant p-values
are displayed in black.
\end{Details}
%
\begin{Value}
A ggplot2 object
\end{Value}
%
\begin{Author}
Geremy Clair
\end{Author}
\HeaderA{singleVariableTrend}{singleVariableTrend()}{singleVariableTrend}
%
\begin{Description}
This function plots the best fitted trend and a dotplot of the analytes intensities.
\end{Description}
%
\begin{Usage}
\begin{verbatim}
singleVariableTrend(
  romics_object,
  variable = "variable",
  factor = "main",
  log_factor = FALSE,
  title = "auto"
)
\end{verbatim}
\end{Usage}
%
\begin{Arguments}
\begin{ldescription}
\item[\code{romics\_object}] has to be a romics\_object created using the function romicsCreateObject() (see dedicated documentation) and has to be numeric.

\item[\code{variable}] has to be a character vector. Romics will try to find an exact match or a partial match of the variable in the analytes names (colnames(romics\_object\$data))

\item[\code{factor}] has to be a factor contained in the romics\_object created, the list of factor can be extracted using the function romicsFactorNames() (see dedicated documentation)

\item[\code{log\_factor}] Boolean, indicate if the factor needs to be log transformed (TRUE), or not (FALSE). By default the factor is not transformed

\item[\code{title}] this has to be a character vector, either "auto" (in this case the rowname containing the variable will be displayed as title) or a title of your choice.
\end{ldescription}
\end{Arguments}
%
\begin{Details}
this function enables to calculate and extract the linear or quardratic trend p values and the best fitted trend (e.g. linear\_increasing, linear\_decreasing, quadratic\_concave,quadratic\_convex)
\end{Details}
%
\begin{Value}
This function returns a list
\end{Value}
%
\begin{Author}
Geremy Clair and Feng Song
\end{Author}
\HeaderA{theme\_map}{theme\_map()}{theme.Rul.map}
%
\begin{Description}
This function is a ggplot2 theme function for the spatial maps that can be created using romics
\end{Description}
%
\begin{Usage}
\begin{verbatim}
theme_map()
\end{verbatim}
\end{Usage}
%
\begin{Details}
This function is the ROP theme for ggplot2, it utilization is similar to any other ggplot2 theme function
\end{Details}
%
\begin{Author}
Geremy Clair
\end{Author}
\HeaderA{theme\_ROP}{theme\_ROP()}{theme.Rul.ROP}
%
\begin{Description}
This function is a ggplot2 theme function
\end{Description}
%
\begin{Usage}
\begin{verbatim}
theme_ROP()
\end{verbatim}
\end{Usage}
%
\begin{Details}
This function is the ROP theme for ggplot2, it utilization is similar to any other ggplot2 theme function
\end{Details}
%
\begin{Author}
Geremy Clair
\end{Author}
\HeaderA{unlog10data}{unlog10data()}{unlog10data}
%
\begin{Description}
Reverses the log10 tranformation of the romics\_object data layer.
\end{Description}
%
\begin{Usage}
\begin{verbatim}
unlog10data(romics_object)
\end{verbatim}
\end{Usage}
%
\begin{Arguments}
\begin{ldescription}
\item[\code{romics\_object}] has to be an romics\_object created using romicsCreateObject()
\end{ldescription}
\end{Arguments}
%
\begin{Details}
will reverse the log10 transformation of the romics\_object
\end{Details}
%
\begin{Value}
This function returns the transformed romics\_object with updated data layer
\end{Value}
%
\begin{Author}
Geremy Clair
\end{Author}
\HeaderA{unlog2data}{unlog2data()}{unlog2data}
%
\begin{Description}
Reverses the log2 tranformation of the romics\_object data layer.
\end{Description}
%
\begin{Usage}
\begin{verbatim}
unlog2data(romics_object)
\end{verbatim}
\end{Usage}
%
\begin{Arguments}
\begin{ldescription}
\item[\code{romics\_object}] has to be an romics\_object created using romicsCreateObject()
\end{ldescription}
\end{Arguments}
%
\begin{Details}
will reverse the log2 transformation of the romics\_object
\end{Details}
%
\begin{Value}
This function returns the transformed romics\_object with updated data layer
\end{Value}
%
\begin{Author}
Geremy Clair
\end{Author}
\HeaderA{upperQuartileNormSample}{upperQuartileNormSample()}{upperQuartileNormSample}
%
\begin{Description}
This function Normalizes the data using the UpperQuartile value (i.e., third quartile).
\end{Description}
%
\begin{Usage}
\begin{verbatim}
upperQuartileNormSample(romics_object)
\end{verbatim}
\end{Usage}
%
\begin{Arguments}
\begin{ldescription}
\item[\code{romics\_object}] has to be an romics\_object created using romicsCreateObject() that has not been previously log-transformed using log10transform()
\end{ldescription}
\end{Arguments}
%
\begin{Details}
UpperQuartile-normalizes the dataset.
\end{Details}
%
\begin{Value}
This function returns the transformed romics\_object with updated data layer.
\end{Value}
%
\begin{Author}
Geremy Clair
\end{Author}
\printindex{}
\end{document}
